% La logique classique — Philosophie 1re Tech — Cours signé OnBuch+
% Programme officiel MINESEC. Compiler avec : tectonic N17-logique-classique.tex
\newcommand{\DOCMATIERE}{Philosophie}
\newcommand{\DOCNIVEAU}{1re Tech}
\newcommand{\DOCMODULE}{Philosophie – classe de Première technique}
\newcommand{\DOCLECON}{N17}
\newcommand{\DOCTITRE}{La logique classique}
% OnBuch+ — préambule « cours signés » ENRICHI (compilé avec tectonic / XeLaTeX)
% Système visuel « Pop » d'OnBuch+ : fond crème (jamais blanc), bordures encre
% épaisses, ombres « dures » décalées, titres Archivo Black en capitales.
% Version MATHÉMATIQUES : graphes (pgfplots/tikz), boîtes pédagogiques
% (définition, propriété, méthode, attention, le savais-tu, exercice résolu,
% exercice, TP, bilan), page de garde et signature OnBuch+ ancrée en bas.
%
% Macros attendues dans main.tex AVANT \input{preamble} :
%   \DOCMATIERE  \DOCNIVEAU  \DOCMODULE  \DOCLECON (numéro)  \DOCTITRE
\documentclass[11pt]{article}

\usepackage{fontspec}
\usepackage[french]{babel}
\usepackage[a4paper,margin=2cm,top=3.4cm,bottom=2.8cm,headheight=1.4cm,footskip=1.3cm]{geometry}
\usepackage{amsmath,amssymb,amsfonts}
\usepackage{mathtools} % \xmapsto, \coloneqq... souvent utilisés par les modèles
\usepackage{cancel} % simplifications barrées (bilans de forces)
% \abs{z} (module, valeur absolue) et \norm{u} (norme), utilisés par les modèles
\providecommand{\abs}{}\let\abs\relax\DeclarePairedDelimiter{\abs}{\lvert}{\rvert}
\providecommand{\norm}{}\let\norm\relax\DeclarePairedDelimiter{\norm}{\lVert}{\rVert}
\usepackage[table]{xcolor} % [table] : fournit aussi \rowcolors (alternance de lignes)
\usepackage{pagecolor}
\usepackage{tikz}
\usetikzlibrary{arrows.meta,positioning,calc,shapes.geometric,shapes.misc,shapes.arrows,decorations.pathmorphing,decorations.pathreplacing,decorations.markings,patterns,shadows,mindmap,trees,fit,backgrounds,matrix,intersections,angles,quotes}
\usepackage{pgfplots}
\usepackage[version=4]{mhchem} % équations nucléaires \ce{^{235}_{92}U}
\usepackage[european,siunitx]{circuitikz} % schémas électriques
\pgfplotsset{compat=1.18}
\usepgfplotslibrary{fillbetween,groupplots} % aires entre courbes (calcul intégral)
\usepackage{enumitem}
\usepackage{fancyhdr}
% [most] charge toutes les bibliothèques tcolorbox (listings, minted, raster,
% documentation...) et gonfle inutilement la mémoire TeX sur un document long
% (~300 boîtes) : on ne charge que ce qui est réellement utilisé.
\usepackage[skins,breakable]{tcolorbox}
\usepackage{graphicx}
\usepackage{colortbl}
\usepackage{booktabs}
\usepackage{tabularx}
\usepackage{multirow}
\usepackage{diagbox} % case d'en-tête diagonale des tableaux à double entrée
% Colonnes extensibles centrée / gauche / droite (C, L, R), souvent utilisées par les modèles
\newcolumntype{C}{>{\centering\arraybackslash}X}
\newcolumntype{L}{>{\raggedright\arraybackslash}X}
\newcolumntype{R}{>{\raggedleft\arraybackslash}X}
\usepackage{array}
\usepackage{multicol}
\usepackage{siunitx}
% parse-numbers=false : accepte aussi \SI{2,5 \times 10^{-3}}{...}
\sisetup{locale=FR,output-decimal-marker={,},group-minimum-digits=4,parse-numbers=false}
\DeclareSIUnit\ohm{\ensuremath{\symup{\Omega}}} % U+2126 absent de la police
\DeclareSIUnit\division{div} % réglages d'oscilloscope (ms/div, V/div)
\DeclareSIUnit\tr{tr} % tours (vitesse de rotation, ex. \SI{1500}{\tr\per\minute})
\DeclareSIUnit\molar{M} % concentration molaire (ex. \SI{3}{\molar}, \SI{1}{\micro\molar})
\DeclareSIUnit\an{an} % année (le modèle confond parfois avec \year, un compteur TeX) — ex. \SI{5}{\kilo\meter\per\an}
\DeclareSIUnit\FCFA{FCFA} % franc CFA (ex. \SI{500}{\FCFA\per\kilo\gram})
\DeclareSIUnit\degreeBrix{\textdegree Brix} % teneur en sucre d'un sirop/jus (réfractométrie)
\DeclareSIUnit\month{mois} % le modèle utilise parfois \month, un compteur TeX, comme unité de durée
\DeclareSIUnit\microliter{\micro\liter} % le modèle écrit parfois \microliter en un seul token
\DeclareSIUnit\million{millions} % dénombrement (ex. \SI{15}{\million\per\mL} spermatozoïdes)
\usepackage{qrcode}
% \degree : commande fréquemment utilisée par les modèles pour un angle ou une
% température en dehors de \SI{}{} (non fournie par les paquets chargés).
\providecommand{\degree}{^{\circ}}
% \qed (fin de démonstration), utilisé par les modèles sans amsthm
\providecommand{\qed}{\hfill$\square$}
% \barre : double barre des valeurs interdites dans un tableau de variations (tkz-tab)
\providecommand{\barre}{\ensuremath{\|}}
% environnement proof (amsthm non chargé) utilisé par les modèles
\ifdefined\proof\else\newenvironment{proof}[1][Démonstration]{\par\noindent\textit{#1.}\ }{\qed\par}\fi
\usepackage{lastpage}
\usepackage{hyperref}
\hypersetup{hidelinks}

% ============================================================
% Palette « Pop » OnBuch+ — mêmes hex que la classe P (plus_ui.dart)
% ============================================================
\definecolor{poporange}{HTML}{F59321}
\definecolor{popdark}{HTML}{E07A0C}
\definecolor{popcream}{HTML}{FFF6E8}
\definecolor{popink}{HTML}{1C1714}
\definecolor{poppurple}{HTML}{7A5AE0}
\definecolor{popgreen}{HTML}{2E9E6A}
\definecolor{poppink}{HTML}{E0457B}
\definecolor{popblue}{HTML}{2F7DE1}
\definecolor{popgold}{HTML}{C98A12}
\definecolor{popmuted}{HTML}{8A7F76}
\definecolor{popline}{HTML}{EADFD2}
\definecolor{popgray}{HTML}{8A7F76}
\colorlet{popgrey}{popgray}
% Alias tolérants (le modèle nomme parfois les couleurs différemment)
\colorlet{popwhite}{white}
\colorlet{popblack}{popink}
\colorlet{popred}{poppink}
\colorlet{popyellow}{popgold}
% Teintes claires (fonds de tableaux/encadrés) — le modèle nomme parfois
% les couleurs avec un suffixe L (ex. popblueL) sans les avoir définies.
\colorlet{poporangeL}{poporange!20}
\colorlet{popdarkL}{popdark!20}
\colorlet{poppurpleL}{poppurple!20}
\colorlet{popgreenL}{popgreen!20}
\colorlet{poppinkL}{poppink!20}
\colorlet{popblueL}{popblue!20}
\colorlet{popgoldL}{popgold!20}
\colorlet{popredL}{popred!20}
\colorlet{popgrayL}{popgray!20}
% Teintes claires pour fonds de tableaux / zones
\colorlet{popblueL}{popblue!10!white}
\colorlet{poporangeL}{poporange!14!white}
\colorlet{popgreenL}{popgreen!12!white}
\colorlet{poppurpleL}{poppurple!12!white}
\colorlet{poppinkL}{poppink!12!white}
\colorlet{popgoldL}{popgold!14!white}

\pagecolor{popcream}

% ============================================================
% Typographie
% ============================================================
\defaultfontfeatures{Ligatures=TeX}
\setmainfont{PlusJakartaSans-Medium.ttf}[Path=../../../fonts/,BoldFont=PlusJakartaSans-Bold.ttf]
% Maths en sans-serif (Fira Math) avec chiffres et lettres droites de Plus
% Jakarta, pour que les formules se fondent dans le texte.
\usepackage[math-style=ISO,bold-style=ISO]{unicode-math}
\setmathfont{FiraMath-Regular.otf}[Path=../../../fonts/]
\setmathfont{PlusJakartaSans-Medium.ttf}[Path=../../../fonts/,range={up/num,up/{Latin,latin},\mathup}]
\usepackage{icomma} % virgule décimale sans espace en mode math
% Caractères mathématiques Unicode absents des polices de texte (Plus Jakarta,
% Archivo Black) : rendus par leur équivalent mathématique, en texte comme en
% formule — sinon ils s'affichent en carré vide (ex. « Arithmétique dans ℤ »).
\usepackage{newunicodechar}
\newunicodechar{ℝ}{\ensuremath{\mathbb{R}}}
\newunicodechar{ℤ}{\ensuremath{\mathbb{Z}}}
\newunicodechar{ℂ}{\ensuremath{\mathbb{C}}}
\newunicodechar{ℕ}{\ensuremath{\mathbb{N}}}
\newunicodechar{ℚ}{\ensuremath{\mathbb{Q}}}
\newunicodechar{𝔻}{\ensuremath{\mathbb{D}}}
\newunicodechar{α}{\ensuremath{\alpha}}
\newunicodechar{β}{\ensuremath{\beta}}
\newunicodechar{γ}{\ensuremath{\gamma}}
\newunicodechar{δ}{\ensuremath{\delta}}
\newunicodechar{ε}{\ensuremath{\varepsilon}}
\newunicodechar{θ}{\ensuremath{\theta}}
\newunicodechar{λ}{\ensuremath{\lambda}}
\newunicodechar{μ}{\ensuremath{\mu}}
\newunicodechar{σ}{\ensuremath{\sigma}}
\newunicodechar{τ}{\ensuremath{\tau}}
\newunicodechar{φ}{\ensuremath{\varphi}}
\newunicodechar{ψ}{\ensuremath{\psi}}
\newunicodechar{ω}{\ensuremath{\omega}}
\newunicodechar{Σ}{\ensuremath{\Sigma}}
% majuscules produites par \MakeUppercase dans les titres (α-aminés -> Α)
\newunicodechar{Α}{\ensuremath{\alpha}}
\newunicodechar{Β}{\ensuremath{\beta}}
\newunicodechar{Γ}{\ensuremath{\Gamma}}
\newunicodechar{Δ}{\ensuremath{\Delta}}
\newunicodechar{Ε}{\ensuremath{\varepsilon}}
\newunicodechar{Θ}{\ensuremath{\Theta}}
\newunicodechar{Λ}{\ensuremath{\Lambda}}
\newunicodechar{Μ}{\ensuremath{\mu}}
\newunicodechar{Τ}{\ensuremath{\tau}}
\newunicodechar{Φ}{\ensuremath{\Phi}}
\newunicodechar{Ψ}{\ensuremath{\Psi}}
\newunicodechar{Ω}{\ensuremath{\Omega}}
\newunicodechar{↦}{\ensuremath{\mapsto}}
\newunicodechar{∈}{\ensuremath{\in}}
\newunicodechar{∉}{\ensuremath{\notin}}
\newunicodechar{∧}{\ensuremath{\wedge}}
\newunicodechar{⇔}{\ensuremath{\Leftrightarrow}}
\newunicodechar{⟺}{\ensuremath{\Longleftrightarrow}}
\newunicodechar{⇒}{\ensuremath{\Rightarrow}}
\newunicodechar{⟹}{\ensuremath{\Longrightarrow}}
\newunicodechar{∘}{\ensuremath{\circ}}
\newunicodechar{∪}{\ensuremath{\cup}}
\newunicodechar{∩}{\ensuremath{\cap}}
\newunicodechar{≡}{\ensuremath{\equiv}}
\newunicodechar{⊂}{\ensuremath{\subset}}
\newunicodechar{⊥}{\ensuremath{\perp}}
\newunicodechar{∃}{\ensuremath{\exists}}
\newunicodechar{∀}{\ensuremath{\forall}}
\newunicodechar{∛}{\ensuremath{\sqrt[3]{\,}}}
\newunicodechar{∜}{\ensuremath{\sqrt[4]{\,}}}
\newunicodechar{⃗}{\ensuremath{{}^{\to}}}
\newunicodechar{ⁿ}{\ifmmode^{n}\else\textsuperscript{\ensuremath{n}}\fi}
\newunicodechar{ˣ}{\ifmmode^{x}\else\textsuperscript{\ensuremath{x}}\fi}
\newunicodechar{ᵇ}{\ifmmode^{b}\else\textsuperscript{\ensuremath{b}}\fi}
\newunicodechar{ʳ}{\ifmmode^{r}\else\textsuperscript{\ensuremath{r}}\fi}
\newunicodechar{ᵘ}{\ifmmode^{u}\else\textsuperscript{\ensuremath{u}}\fi}
\newunicodechar{ʸ}{\ifmmode^{y}\else\textsuperscript{\ensuremath{y}}\fi}
\newunicodechar{ᵃ}{\ifmmode^{a}\else\textsuperscript{\ensuremath{a}}\fi}
\newunicodechar{ᵉ}{\ifmmode^{e}\else\textsuperscript{\ensuremath{e}}\fi}
\newunicodechar{ᵏ}{\ifmmode^{k}\else\textsuperscript{\ensuremath{k}}\fi}
\newunicodechar{ᵖ}{\ifmmode^{p}\else\textsuperscript{\ensuremath{p}}\fi}
\newunicodechar{ᴺ}{\ifmmode^{N}\else\textsuperscript{\ensuremath{N}}\fi}
\newunicodechar{ᶜ}{\ifmmode^{c}\else\textsuperscript{\ensuremath{c}}\fi}
\newunicodechar{ʰ}{\ifmmode^{h}\else\textsuperscript{\ensuremath{h}}\fi}
\newunicodechar{ᵠ}{\ifmmode^{q}\else\textsuperscript{\ensuremath{q}}\fi}
\newunicodechar{ₙ}{\ifmmode_{n}\else\textsubscript{\ensuremath{n}}\fi}
\newunicodechar{ₐ}{\ifmmode_{a}\else\textsubscript{\ensuremath{a}}\fi}
\newunicodechar{ᵢ}{\ifmmode_{i}\else\textsubscript{\ensuremath{i}}\fi}
\newunicodechar{ᵣ}{\ifmmode_{r}\else\textsubscript{\ensuremath{r}}\fi}
\newunicodechar{ₖ}{\ifmmode_{k}\else\textsubscript{\ensuremath{k}}\fi}
\newunicodechar{ᵦ}{\ifmmode_{\beta}\else\textsubscript{\ensuremath{\beta}}\fi}
\newunicodechar{ₓ}{\ifmmode_{x}\else\textsubscript{\ensuremath{x}}\fi}
\newfontfamily\popdisplay{ArchivoBlack-Regular.ttf}[Path=../../../fonts/]
\newfontfamily\poplogo{SpaceGrotesk-Bold.ttf}[Path=../../../fonts/]
\newfontfamily\popsemi{PlusJakartaSans-SemiBold.ttf}[Path=../../../fonts/]
\newfontfamily\popxbold{PlusJakartaSans-ExtraBold.ttf}[Path=../../../fonts/]
\newfontfamily\popxboldls{PlusJakartaSans-ExtraBold.ttf}[Path=../../../fonts/,LetterSpace=3.0]

\setlist[enumerate]{leftmargin=1.7em,itemsep=5pt,topsep=4pt}
\setlist[itemize]{leftmargin=1.7em,itemsep=4pt,topsep=4pt,label={\color{poporange}\textbullet}}
\setlength{\parindent}{0pt}
\setlength{\parskip}{5pt}
\renewcommand{\arraystretch}{1.25}

% pgfplots : style Pop par défaut
\pgfplotsset{
  every axis/.append style={axis line style={popink,line width=1pt},tick style={popink},
    grid style={popline,line width=0.5pt},label style={font=\small},tick label style={font=\footnotesize}},
  popcurve/.style={poporange,line width=1.8pt,no markers,smooth},
}
% Même style disponible dans un \draw TikZ simple (hors pgfplots)
\tikzset{popcurve/.style={poporange,line width=1.8pt}}
\tikzset{popgrid/.style={popline,very thin}}
\colorlet{popcurve}{poporange} % le nom du style est parfois utilisé comme couleur

% ============================================================
% Pill / squiggle
% ============================================================
\newcommand{\poppill}[3]{%
  \tikz[baseline=-0.55ex]\node[draw=popink,line width=1.2pt,rounded corners=2.6pt,%
    fill=#2,inner xsep=6.5pt,inner ysep=3pt]{%
    \popxboldls\fontsize{7}{7}\selectfont\color{#3}\MakeUppercase{#1}};%
}
\newcommand{\popsquiggle}[1]{%
  \tikz[baseline=-0.4ex]{
    \draw[#1,line width=2.6pt,line cap=round]
      plot [smooth] coordinates {(0,0.09) (0.34,0.01) (0.68,0.21) (1.02,0.01) (1.36,0.10)};
  }%
}
% Mot-clé surligné (marqueur orange sous le texte)
\newcommand{\cle}[1]{{\popxbold\color{popdark}#1}}

% ============================================================
% En-tête / pied
% ============================================================
\pagestyle{fancy}
\fancyhf{}
\renewcommand{\headrulewidth}{0pt}
\renewcommand{\footrulewidth}{0pt}
\lhead{\raisebox{-0.15ex}{\poplogo\fontsize{13}{13}\selectfont\color{popink}OnBuch\color{poporange}+}}
\rhead{\poppill{\DOCMATIERE\ \textbullet\ \DOCNIVEAU\ \textbullet\ Leçon \DOCLECON}{white}{popink}}
\lfoot{\popxbold\fontsize{7.6}{7.6}\selectfont\color{popmuted}\MakeUppercase{Cours signé OnBuch+}\ \textbullet\ {\popsemi\DOCTITRE}}
\rfoot{\tikz[baseline=-0.6ex]\node[rounded corners=8.5pt,fill=popink,minimum height=17pt,minimum width=17pt,inner xsep=5pt,inner sep=0]{\popxbold\fontsize{8}{8}\selectfont\color{popcream}\thepage\,/\,\pageref*{LastPage}};}

% ============================================================
% Titres
% ============================================================
\newcommand{\chaptitre}[2]{%
  \begin{tcolorbox}[enhanced,colback=poporange,colframe=popink,boxrule=2.6pt,arc=11pt,
      drop shadow=popink,left=15pt,right=15pt,top=12pt,bottom=11pt,before skip=3pt,after skip=18pt]
    \poppill{#2}{popink}{popcream}\par\vspace{9pt}
    {\raggedright\hyphenpenalty=10000\exhyphenpenalty=10000\popdisplay\fontsize{22}{24}\selectfont\color{popcream}\MakeUppercase{#1}\par}
  \end{tcolorbox}
}
% Section numérotée : pastille orange avec le numéro + titre Archivo
\newcounter{popsec}
\newcommand{\coursec}[1]{%
  \stepcounter{popsec}\setcounter{popsub}{0}%
  \par\vspace{16pt}\noindent
  \begin{minipage}{\linewidth}
  \tikz[baseline=-0.7ex]\node[draw=popink,line width=1.6pt,fill=poporange,rounded corners=4pt,minimum size=22pt,inner sep=0]{\popdisplay\fontsize{12}{12}\selectfont\color{popcream}\thepopsec};%
  \hspace{8pt}{\popdisplay\fontsize{15}{17}\selectfont\color{popink}\MakeUppercase{#1}}\par
  \vspace{4pt}\noindent{\color{poporange}\rule{36pt}{3.6pt}}
  \end{minipage}\par\nopagebreak\vspace{8pt}
}
\newcounter{popsub}
\newcommand{\courssub}[1]{%
  \stepcounter{popsub}%
  \par\vspace{9pt}\noindent{\popxbold\fontsize{11.5}{13}\selectfont\color{popink}{\color{poporange}\thepopsec.\thepopsub}\hspace{6pt}#1}\par\nopagebreak\vspace{4pt}
}

% ============================================================
% Boîtes pédagogiques — même recette : fond blanc, bordure épaisse colorée,
% ombre dure encre, badge plein. Titre optionnel [#1].
% ============================================================
\tcbset{popbase/.style={breakable,enhanced,colback=white,boxrule=2.3pt,arc=9pt,
  shadow={3pt}{-3pt}{0pt}{popink},left=14pt,right=14pt,top=11pt,bottom=11pt,before skip=11pt,after skip=15pt,
  fontupper=\popsemi\fontsize{10.6}{14.5}\selectfont\color{popink},
  fontlower=\popsemi\fontsize{10.6}{14.5}\selectfont\color{popink}}}
% Coche dessinée (le glyphe ✓ n'existe pas dans les polices du document)
\newcommand{\popcheck}[1]{\tikz[baseline=-0.5ex]\draw[#1,line width=1.6pt,line cap=round,line join=round](-0.13,0.0)--(-0.04,-0.09)--(0.13,0.1);}
\providecommand{\coche}{\popcheck{popgreen}}
\providecommand{\resultat}[1]{\par\smallskip\noindent\fcolorbox{popgreen}{popgreenL}{\parbox{\dimexpr\linewidth-2\fboxsep-2\fboxrule\relax}{\textbf{Résultat.} #1}}\par\smallskip}
\newcommand{\popboxhead}[3]{\poppill{#1}{#2}{white}\ifx\relax#3\relax\else\hspace{6pt}{\popxbold\fontsize{10.6}{12}\selectfont\color{popink}#3}\fi\par\vspace{7pt}}

\newtcolorbox{aretenir}[1][]{popbase,colframe=popgreen,before upper={\popboxhead{À retenir}{popgreen}{#1}}}
\newtcolorbox{exemplebox}[1][]{popbase,colframe=poporange,before upper={\popboxhead{Exemple}{poporange}{#1}}}
\newtcolorbox{definition}[1][]{popbase,colframe=popblue,before upper={\popboxhead{Définition}{popblue}{#1}}}
\newtcolorbox{propriete}[1][]{popbase,colframe=poppurple,before upper={\popboxhead{Propriété}{poppurple}{#1}}}
\newtcolorbox{methode}[1][]{popbase,colframe=popgold,before upper={\popboxhead{Méthode}{popgold}{#1}}}
\newtcolorbox{attention}[1][]{popbase,colframe=poppink,before upper={\popboxhead{Attention}{poppink}{#1}}}
\newtcolorbox{experience}[1][]{popbase,colframe=popblue,colback=popblueL,before upper={\popboxhead{Expérience}{popblue}{#1}}}
% « Le savais-tu ? » : carte violette pleine (contexte, culture, Cameroun)
\newtcolorbox{savaistu}[1][]{popbase,colback=poppurple,colframe=popink,
  fontupper=\popsemi\fontsize{10.4}{14.2}\selectfont\color{white},
  before upper={\poppill{Le savais-tu\,?}{white}{poppurple}\ifx\relax#1\relax\else\hspace{6pt}{\popxbold\color{white}#1}\fi\par\vspace{7pt}}}
% Exercice résolu : énoncé en haut, \tcblower puis la solution sur fond crème
\newcounter{popexo}\newcounter{popexr}
\newtcolorbox{exoresolu}[1][]{popbase,colframe=poporange,colbacklower=popcream,
  segmentation style={popink,line width=1.2pt,dashed},
  before upper={\stepcounter{popexr}\popboxhead{Exercice résolu \thepopexr}{poporange}{#1}},
  before lower={\poppill{Solution}{popink}{popcream}\par\vspace{6pt}}}
% Exercice (non résolu en place) — la correction est dans la partie Corrigés
\newtcolorbox{exercice}[1][]{popbase,colframe=popink,
  before upper={\stepcounter{popexo}\popboxhead{Exercice \thepopexo}{popink}{#1}}}
% Corrigé
\newtcolorbox{corrige}[1][]{popbase,colframe=popgreen,colback=popgreenL,
  before upper={\popboxhead{Corrigé}{popgreen}{#1}}}
% Objectifs / prérequis / bilan
\newtcolorbox{objectifs}{popbase,colframe=popink,colback=poporangeL,
  before upper={\popboxhead{Objectifs de la leçon}{popink}{}}}
\newtcolorbox{prerequis}{popbase,colframe=popink,colback=popblueL,
  before upper={\popboxhead{Prérequis}{popblue}{}}}
\newtcolorbox{bilan}{popbase,colframe=popink,colback=white,boxrule=2.8pt,
  before upper={\popboxhead{Fiche bilan}{poporange}{L'essentiel à connaître pour l'examen}}}
% Figure encadrée avec légende
\newtcolorbox{popfigure}[1][]{enhanced,breakable=false,colback=white,colframe=popline,boxrule=1.4pt,arc=8pt,
  left=10pt,right=10pt,top=10pt,bottom=8pt,before skip=10pt,after skip=12pt,halign=center,
  lower separated=false,halign lower=center,
  fontlower=\popsemi\fontsize{9}{11.5}\selectfont\color{popmuted},
  #1}
\newcommand{\legende}[1]{\par\vspace{6pt}{\popsemi\fontsize{9}{11.5}\selectfont\color{popmuted}#1}}

% ============================================================
% Signature OnBuch+
% ============================================================
\newcommand{\poplogoinline}[1]{{\poplogo\fontsize{#1}{#1}\selectfont\color{popink}OnBuch\color{poporange}+}}
\newcommand{\signatureonbuch}{%
  \begin{tcolorbox}[enhanced,colback=popink,colframe=popink,boxrule=2.2pt,arc=10pt,
      drop shadow=poporange,left=14pt,right=14pt,top=10pt,bottom=10pt,before skip=0pt,after skip=0pt]
    \tikz[baseline=-0.7ex]\node[circle,fill=popgreen,draw=popcream,line width=1.3pt,minimum size=22pt,inner sep=0]{\popcheck{white}};%
    \hspace{9pt}\begin{minipage}[c]{0.62\linewidth}
      {\popxbold\fontsize{12}{13}\selectfont\color{popcream}Signé\ \textbullet\ {\poplogo OnBuch\color{poporange}+}}\par\vspace{2pt}
      {\raggedright\popsemi\fontsize{8}{10.5}\selectfont\color{popline}\MakeUppercase{Équipe pédagogique OnBuch+}\\\MakeUppercase{Conforme au programme officiel MINESEC}\par}
    \end{minipage}\hfill
    {\poplogo\fontsize{22}{22}\selectfont\color{popcream}OnBuch\color{poporange}+}
  \end{tcolorbox}%
}

% ============================================================
% Page de garde
% #1 titre · #2 matière · #3 niveau · #4 module · #5 n° leçon · #6 durée ·
% #7 description / objectifs (texte ou itemize)
% ============================================================
\newcommand{\pagegarde}[7]{%
  \thispagestyle{empty}
  \newgeometry{a4paper,margin=1.7cm,top=1.6cm,bottom=1.4cm}%
  \begin{tikzpicture}[remember picture,overlay]
    \fill[poporange!18] ([xshift=-3.2cm,yshift=-3.6cm]current page.north east) circle (4.6cm);
    \fill[poppurple!14] ([xshift=2.4cm,yshift=7.6cm]current page.south west) circle (3.8cm);
  \end{tikzpicture}%
  {\poplogo\fontsize{30}{30}\selectfont\color{popink}OnBuch\color{poporange}+}\hfill
  \raisebox{0.6ex}{\poppill{Cours signé \textbullet\ Premium}{popink}{popcream}}\par
  \vspace{1pt}\popsquiggle{poppurple}\par
  \vspace{8pt}
  \begin{tcolorbox}[enhanced,colback=poporange,colframe=popink,boxrule=2.8pt,arc=13pt,
      drop shadow=popink,left=18pt,right=18pt,top=12pt,bottom=13pt,after skip=10pt]
    \poppill{#2 \textbullet\ #3}{popink}{popcream}\hspace{5pt}\poppill{#4}{white}{popink}\par\vspace{8pt}
    {\popdisplay\fontsize{14}{14}\selectfont\color{popink}LEÇON #5}\par\vspace{4pt}
    {\raggedright\hyphenpenalty=10000\exhyphenpenalty=10000\popdisplay\fontsize{25}{27}\selectfont\color{popcream}\MakeUppercase{#1}\par}
    \vspace{8pt}
    {\popxbold\fontsize{9.5}{9.5}\selectfont\color{popink}\MakeUppercase{Durée conseillée : #6}}
  \end{tcolorbox}
  \begin{tcolorbox}[enhanced,colback=white,colframe=popink,boxrule=2.2pt,arc=10pt,
      drop shadow=popink,left=16pt,right=16pt,top=10pt,bottom=10pt,after skip=8pt]
    \poppill{Dans cette leçon}{poppurple}{white}\par\vspace{5pt}
    \popsemi\fontsize{9.3}{12.6}\selectfont\color{popink}#7
  \end{tcolorbox}
  \noindent\begin{minipage}[t]{0.487\linewidth}
    \begin{tcolorbox}[enhanced,colback=popgreen,colframe=popink,boxrule=2.2pt,arc=10pt,
        drop shadow=popink,left=11pt,right=11pt,top=10pt,bottom=10pt,height=3.1cm,valign=center]
      \tcbox[colback=white,colframe=popink,boxrule=1.3pt,arc=5pt,boxsep=0pt,nobeforeafter,
        left=3pt,right=3pt,top=3pt,bottom=3pt,valign=center]{\qrcode[height=1.9cm]{https://play.google.com/store/apps/details?id=com.luvvix.onbuch&hl=fr}}%
      \hspace{8pt}\begin{minipage}[c]{0.5\linewidth}
        {\popdisplay\fontsize{11}{12.5}\selectfont\color{white}\MakeUppercase{Télécharge OnBuch}}\par\vspace{4pt}
        {\raggedright\popsemi\fontsize{8.4}{11}\selectfont\color{white}Gratuit \textbullet\ Google Play\\Tuteur IA, annales, concours, résultats\par}
      \end{minipage}
    \end{tcolorbox}
  \end{minipage}\hfill
  \begin{minipage}[t]{0.487\linewidth}
    \begin{tcolorbox}[enhanced,colback=poppurple,colframe=popink,boxrule=2.2pt,arc=10pt,
        drop shadow=popink,left=11pt,right=11pt,top=10pt,bottom=10pt,height=3.1cm,valign=center]
      \tikz[baseline=-0.7ex]\node[circle,fill=white,draw=popink,line width=1.3pt,minimum size=1.6cm,inner sep=0]{\popdisplay\fontsize{24}{24}\selectfont\color{poppurple}+};%
      \hspace{8pt}\begin{minipage}[c]{0.6\linewidth}
        {\popdisplay\fontsize{11}{12.5}\selectfont\color{white}\MakeUppercase{Passe à OnBuch+}}\par\vspace{4pt}
        {\raggedright\popsemi\fontsize{8.4}{11}\selectfont\color{white}Tous les cours signés, Tuteur IA illimité, fiches PDF — onglet OnBuch+ de l'app\par}
      \end{minipage}
    \end{tcolorbox}
  \end{minipage}
  \vspace{8pt}
  \popcontenu
  \vfill
  \signatureonbuch
  \restoregeometry
  \newpage
}

% Rangée « Ce cours contient » (page de garde)
\newcommand{\poptile}[3]{% #1 couleur #2 gros texte #3 légende
  \begin{tcolorbox}[enhanced,colback=#1,colframe=popink,boxrule=2pt,arc=9pt,shadow={2.5pt}{-2.5pt}{0pt}{popink},
      left=6pt,right=6pt,top=9pt,bottom=9pt,halign=center,valign=center,height=1.75cm,width=\linewidth,nobeforeafter]
    {\popdisplay\fontsize{12.5}{13}\selectfont\color{white}\MakeUppercase{#2}}\par\vspace{3pt}
    {\centering\popsemi\fontsize{7.8}{9.5}\selectfont\color{white}#3\par}
  \end{tcolorbox}}
\newcommand{\popcontenu}{%
  \par\noindent{\popxboldls\fontsize{7.5}{7.5}\selectfont\color{popmuted}CE COURS SIGNÉ CONTIENT}\par\vspace{7pt}
  \noindent\begin{minipage}[t]{0.235\linewidth}\poptile{popblue}{Cours}{complet, illustré, pas à pas}\end{minipage}\hfill
  \begin{minipage}[t]{0.235\linewidth}\poptile{poporange}{Méthodes}{et exercices résolus}\end{minipage}\hfill
  \begin{minipage}[t]{0.235\linewidth}\poptile{poppink}{Exercices}{type examen + corrigés}\end{minipage}\hfill
  \begin{minipage}[t]{0.235\linewidth}\poptile{popgreen}{Fiche bilan}{l'essentiel à retenir}\end{minipage}\par}

% Sommaire
\newtcolorbox{sommaire}{enhanced,colback=white,colframe=popink,boxrule=2.2pt,arc=10pt,
  drop shadow=popink,left=18pt,right=18pt,top=13pt,bottom=13pt,before skip=6pt,after skip=20pt,
  before upper={\poppill{Sommaire}{poppurple}{white}\par\vspace{9pt}\popsemi\fontsize{10.6}{14.5}\selectfont\color{popink}}}

% Signature de fin de cours — poussée tout en bas de la dernière page
\newcommand{\findecours}{%
  \par\vfill\nopagebreak
  \begin{center}\popsquiggle{poporange}\end{center}\vspace{4pt}
  \signatureonbuch
}
\AtBeginDocument{\def\llbracket{\mathopen{[\mkern-3.5mu[}}\def\rrbracket{\mathclose{]\mkern-3.5mu]}}} % intervalles d'entiers (glyphe ⟦ absent de la police)
% Glyphes absents de Fira Math : redéfinis à partir de glyphes disponibles
\AtBeginDocument{%
  \def\setminus{\mathbin{\backslash}}\let\smallsetminus\setminus
  \def\vdots{\mathord{\vbox{\baselineskip3.2pt\lineskiplimit0pt\kern4pt\hbox{.}\hbox{.}\hbox{.}}}}%
  \def\ddots{\mathinner{\mkern1mu\raise6.5pt\hbox{.}\mkern2mu\raise3.6pt\hbox{.}\mkern2mu\raise0.7pt\hbox{.}\mkern1mu}}%
  \def\triangle{\mathord{\scalebox{0.9}{$\symup{\Delta}$}}}%
  \def\nabla{\mathord{\raisebox{\depth}{\scalebox{1}[-1]{$\symup{\Delta}$}}}}%
  \def\checkmark{\popcheck{popdark}}%
  \def\star{\mathbin{\ast}}\def\bigstar{\mathord{\ast}}%
  \def\diamond{\mathbin{\scalebox{0.6}{$\Diamond$}}}%
  \def\bigcup{\mathop{\mathchoice{\raisebox{-0.3em}{\scalebox{1.7}{$\cup$}}}{\scalebox{1.25}{$\cup$}}{\cup}{\cup}}\displaylimits}%
  \def\bigcap{\mathop{\mathchoice{\raisebox{-0.3em}{\scalebox{1.7}{$\cap$}}}{\scalebox{1.25}{$\cap$}}{\cap}{\cap}}\displaylimits}%
  \def\lesseqqgtr{\mathrel{\lessgtr}}\def\bigoplus{\mathop{\oplus}\displaylimits}%
}
\providecommand{\ssi}{si et seulement si} % abréviation fréquente
\AtBeginDocument{\providecommand{\wideparen}{\overparen}} % arc de cercle

%% FIGFIX-BEGIN v4 — correctifs globaux des figures (docs/onbuch-generation/tools/figfix)
\usepackage{adjustbox}\usepackage{etoolbox}
% toute figure TikZ/pgfplots plus large que la ligne est réduite à la largeur disponible
\BeforeBeginEnvironment{tikzpicture}{\begin{adjustbox}{max width=\linewidth}}
\AfterEndEnvironment{tikzpicture}{\end{adjustbox}}
% légendes pgfplots lisibles par-dessus les courbes
\pgfplotsset{every axis/.append style={legend style={fill=white,fill opacity=0.92,text opacity=1,draw=black!15,inner sep=3pt}}}
% étiquettes d'arêtes (syntaxe edge["..."]) avec fond blanc
\tikzset{every edge quotes/.append style={fill=white,inner sep=1.2pt}}
%% FIGFIX-END
\begin{document}

\pagegarde{La logique classique}{Philosophie}{1re Tech}{Philosophie – classe de Première technique}{N17}{6 séances (fiche de progression harmonisée)}{Cette leçon présente la logique classique comme science des lois formelles de la pensée correcte. De la définition du concept à l'analyse du jugement et du raisonnement, elle outille l'élève pour penser rigoureusement, argumenter et déjouer les sophismes de la vie quotidienne.
\vspace{4pt}
\begin{multicols}{2}\small
\begin{itemize}
\item À la fin de cette leçon, je sais définir la logique et expliquer son importance pour la pensée et la communication
\item À la fin de cette leçon, je sais définir le concept et distinguer sa compréhension de son extension
\item À la fin de cette leçon, je sais analyser la structure d'un jugement (sujet, copule, prédicat) et classer les jugements selon la quantité et la qualité
\item À la fin de cette leçon, je sais identifier et construire un raisonnement déductif (syllogisme) et vérifier sa validité
\item À la fin de cette leçon, je sais distinguer raisonnement déductif, inductif et par analogie
\item À la fin de cette leçon, je sais énoncer et appliquer les trois principes logiques de la pensée (identité, non-contradiction, tiers exclu)
\end{itemize}\end{multicols}}

\begin{sommaire}
\begin{multicols}{2}
\begin{enumerate}
\item Définition et importance de la logique
\item Le concept
\item Le jugement
\item Le raisonnement
\item Les principes logiques de la pensée
\item Application : penser juste dans la vie quotidienne
\item Activité d'intégration
\item Méthodes et exercices résolus
\item Exercices
\item Corrigés des exercices
\item Fiche bilan
\end{enumerate}
\end{multicols}
\end{sommaire}

\begin{objectifs}
\begin{itemize}
\item À la fin de cette leçon, je sais définir la logique et expliquer son importance pour la pensée et la communication
\item À la fin de cette leçon, je sais définir le concept et distinguer sa compréhension de son extension
\item À la fin de cette leçon, je sais analyser la structure d'un jugement (sujet, copule, prédicat) et classer les jugements selon la quantité et la qualité
\item À la fin de cette leçon, je sais identifier et construire un raisonnement déductif (syllogisme) et vérifier sa validité
\item À la fin de cette leçon, je sais distinguer raisonnement déductif, inductif et par analogie
\item À la fin de cette leçon, je sais énoncer et appliquer les trois principes logiques de la pensée (identité, non-contradiction, tiers exclu)
\item À la fin de cette leçon, je sais déceler les sophismes et paralogismes dans un discours courant (marché, politique, réseaux sociaux)
\item À la fin de cette leçon, je sais rédiger et justifier une démarche argumentative rigoureuse à partir d'une situation concrète
\end{itemize}
\end{objectifs}

\begin{prerequis}
\begin{itemize}
\item Notion d'argumentation et de thèse (français, classes de Seconde)
\item Distinction entre opinion et vérité (introduction à la philosophie, début d'année)
\item Notion de définition et de vocabulaire précis (Seconde)
\item Capacité à lire et analyser un texte argumentatif court
\end{itemize}
\end{prerequis}

\newpage

\coursec{Définition et importance de la logique}

\courssub{Accroche et problématique}

Au marché Mokolo à Yaoundé, un commerçant affirme avec assurance : « Tous les clients qui marchandent achètent moins cher. Or tu as marchandé, donc tu as payé moins cher. » Un élève de Première, un peu surpris, réplique : « Mais j'ai marchandé et j'ai payé le même prix ! » Qui a raison ? Le raisonnement du commerçant est-il valide ? Pour trancher, il faut des règles précises : c'est toute la logique classique, héritée d'Aristote, qui permet de distinguer un raisonnement correct d'un raisonnement fallacieux. Cette section pose les fondements : qu'est-ce que la logique ? Pourquoi est-elle indispensable, tant pour la rigueur scientifique que pour la vie quotidienne ?

\courssub{Étymologie et définition}

Le mot « logique » vient du grec ancien \emph{logikè} (technè), « l'art du discours », lui-même dérivé de \cle{logos}. Ce terme polysémique désigne à la fois la \emph{parole}, le \emph{langage}, la \emph{raison} et le \emph{discours rationnel structuré}. Chez les Grecs, le \emph{logos} est ce qui permet de mettre de l'ordre dans le chaos, de rendre compte du réel par un enchaînement cohérent.

\begin{definition}[La logique]
La \textbf{logique} est la \textbf{science des lois formelles de la pensée correcte}. Elle étudie les \textbf{conditions de validité du raisonnement}, c'est-à-dire les règles qui garantissent que, si les prémisses sont vraies, la conclusion ne peut pas être fausse. Elle s'intéresse à la \emph{forme} du raisonnement, indépendamment de son \emph{contenu} (de la matière traitée).
\end{definition}

\courssub{Logique formelle et logique matérielle}

Il est crucial de distinguer deux acceptions historiques et épistémologiques :

\begin{itemize}
    \item \textbf{La logique formelle (ou logique classique, aristotélicienne) :} Elle étudie les \emph{formes} de la pensée (concept, jugement, raisonnement) et les lois qui régissent leur enchaînement valide. Elle ne se préoccupe pas de la vérité matérielle des prémisses, mais de la correction de l'inférence. C'est la logique du « si... alors... ».
    \item \textbf{La logique matérielle :} Elle s'intéresse au \emph{contenu} de la pensée, aux conditions de vérité des propositions par rapport au réel. Elle relève davantage de l'épistémologie ou de la théorie de la connaissance.
\end{itemize}

\begin{aretenir}[Distinction clé]
La logique formelle répond à la question : « Ce raisonnement est-il \textbf{valide} ? » (la conclusion suit-elle nécessairement des prémisses ?).\\
La logique matérielle (ou la critique du contenu) répond à : « Ce raisonnement est-il \textbf{vrai} ? » (les prémisses sont-elles exactes ? la conclusion correspond-elle au réel ?).
\end{aretenir}

\courssub{Aristote, fondateur de la logique classique}

Aristote (384--322 av. J.-C.) est le premier à avoir systématisé la logique comme discipline autonome. Il n'a jamais utilisé le mot « logique » : il parlait d'\emph{analytique} (l'art d'analyser, de décomposer le raisonnement). Ses six traités logiques ont été regroupés par ses successeurs (les Péripatéticiens) sous le titre de \textbf{l'Organon} (« l'instrument »), car la logique est l'instrument commun à toutes les sciences.

\begin{savaistu}[Aristote et le mot « logique »]
Aristote n'a jamais écrit le mot « logique ». Il désignait sa discipline par les termes \emph{analytique} (analytique première et seconde) ou \emph{dialectique}. C'est le stoïcien Zénon de Cition (vers 300 av. J.-C.) qui a forgé l'adjectif \emph{logikè}, et les commentateurs alexandrins (Andronicos de Rhodes, 1er s. av. J.-C.) qui ont baptisé le corpus aristotélicien « l'Organon ». Le terme « logique » comme nom commun s'impose seulement au Moyen Âge (traduction latine \emph{logica}).
\end{savaistu}

\begin{popfigure}
\begin{tikzpicture}[
    node distance=1.5cm and 2cm,
    every node/.style={font=\small, align=center},
    box/.style={draw=popblue, thick, rounded corners, fill=popblueL, minimum width=2.5cm, minimum height=1cm},
    arrow/.style={-Stealth, thick, popdark}
]
% Nœuds
\node[box, align=center] (aristote){Aristote\\ (384--322 av. J.-C.)\\(L'Organon)};
\node[box, right=of aristote, align=center] (moyenage){Logique scolastique\\ (Moyen Âge)\\(Abélard, Occam,\\Thomas d'Aquin)};
\node[box, right=of moyenage, align=center] (classique){Logique classique\\ (XVIIe--XIXe s.)\\(Port-Royal, Kant,\\Hamilton)};
\node[box, below right=1.2cm and 2cm of classique, align=center] (moderne){Logique moderne /\ symbolique\\ (XIXe--XXe s.)\\(Frege, Russell,\\Wittgenstein, Gödel)};
\node[box, below left=1.2cm and 2cm of moderne, align=center] (informatique){Applications\\ (Informatique, IA,\\Preuve formelle,\\Linguistique)};
% Flèches
\draw[arrow] (aristote) -- (moyenage) node[midway, above, sloped, font=\tiny, align=center]{Commentaires,\\traductions};
\draw[arrow] (moyenage) -- (classique) node[midway, above, sloped, font=\tiny] {Réforme pédagogique};
\draw[arrow] (classique) -- (moderne) node[midway, right, font=\tiny] {Formalisation mathématique};
\draw[arrow] (moderne) -- (informatique) node[midway, below, font=\tiny, align=center]{Calculabilité,\\modèles formels};
% Complément Bac
\node[draw=poporange, thick, rounded corners, fill=poporangeL, minimum width=3cm, minimum height=1.2cm, below left=0.5cm and 1cm of aristote, font=\small, align=center] (complement) {\textbf{Complément Bac}\\Logique symbolique\\(Frege, Russell)\\\emph{Hors programme}\\\emph{1re Tech}};
\draw[arrow, poporange, dashed] (complement) -- (moderne);
\end{tikzpicture}
\legende{Arbre généalogique simplifié de la logique : de l'Organon d'Aristote à la logique symbolique moderne (complément Bac).}
\end{popfigure}

\courssub{Importance théorique : fondement des sciences}

Théoriquement, la logique est la \textbf{condition de possibilité de toute science rigoureuse}.
\begin{itemize}
    \item \textbf{Cohérence interne :} Elle garantit qu'un système de pensée ne se contredit pas. En mathématiques, un système axiomatique doit être \emph{consistant} (on ne peut pas démontrer à la fois $P$ et $\neg P$).
    \item \textbf{Déduction rigoureuse :} Elle permet de passer de principes premiers (axiomes) à des théorèmes par des étapes incontestables. Sans logique, pas de démonstration mathématique, pas de modélisation physique fiable.
    \item \textbf{Clarté conceptuelle :} Elle force à définir les termes, à distinguer l'ambiguïté de la précision, condition \emph{sine qua non} du progrès scientifique.
\end{itemize}

\begin{exemplebox}[Logique et mathématiques : le syllogisme d'Euclide]
La géométrie d'Euclide (\emph{Éléments}, vers -300) est le modèle historique de la déduction logique :
\begin{enumerate}
    \item \textbf{Axiomes / Postulats} (ex. : « Par deux points distincts passe une unique droite »).
    \item \textbf{Définitions} précises (cercle, triangle, parallèles).
    \item \textbf{Raisonnements syllogistiques} enchaînés pour prouver des théorèmes (ex. : somme des angles d'un triangle = 180°).
\end{enumerate}
La validité de la démonstration ne dépend pas de la figure tracée sur le sable, mais de la forme logique de l'enchaînement.
\end{exemplebox}

\courssub{Importance pratique : argumenter, débattre, se défendre}

Dans la vie quotidienne, au Cameroun comme ailleurs, la logique est un \textbf{outil de lucidité et de liberté}. Elle permet de :
\begin{itemize}
    \item \textbf{Structurer sa pensée} pour s'exprimer clairement (exposé oral, rédaction, lettre de motivation).
    \item \textbf{Évaluer les discours d'autrui} : discours politiques, publicités, rumeurs sur WhatsApp, « fake news », arguments fallacieux dans un conflit familial ou professionnel.
    \item \textbf{Déjouer les manipulations} : \emph{argument d'autorité} (« Le ministre l'a dit, donc c'est vrai »), \emph{appel à la peur} (« Si tu ne votes pas pour moi, le chaos arrive »), \emph{faux dilemme} (« Soit tu es avec nous, soit tu es contre le pays »), \emph{généralisation hâtive} (« Un Bamileke m'a arnaqué, donc tous les Bamileke sont des voleurs »).
    \item \textbf{Négocier et décider} : au marché, en entreprise, dans une association, distinguer un bon argument d'un sophisme évite de mauvais engagements.
\end{itemize}

\begin{exoresolu}[Analyse d'un argument publicitaire]
\textbf{Énoncé :} Une affiche à Douala proclame : « 9 femmes sur 10 préfèrent la lessive \emph{Blanc Éclat}. Choisissez le meilleur pour votre linge ! »
\textbf{Analyse logique :}
\begin{enumerate}
    \item \textbf{Conclusion implicite :} « Vous devez acheter \emph{Blanc Éclat}. »
    \item \textbf{Prémisse :} « 9 femmes sur 10 préfèrent \emph{Blanc Éclat}. »
    \item \textbf{Critique (validité) :} Le raisonnement est un \emph{argument de popularité} (\emph{ad populum}). La popularité ne prouve pas la qualité intrinsèque (efficacité, composition, prix). C'est un sophisme.
    \item \textbf{Critique (vérité) :} D'où vient le chiffre ? Échantillon ? Méthode ? Questionnement biaisé ? Sans source vérifiable, la prémisse est suspecte.
\end{enumerate}
\textbf{Conclusion :} L'argument ne tient pas logiquement. Le consommateur rationnel compare composition, prix, labels, tests indépendants.
\end{exoresolu}

\courssub{Distinction fondamentale : Vérité vs Validité}

C'est le point le plus délicat et le plus important pour l'examen. Un raisonnement peut être \textbf{valide} (forme correcte) mais aboutir à une conclusion \textbf{fausse} (prémisses fausses). Inversement, un raisonnement \textbf{invalide} (forme incorrecte) peut aboutir par hasard à une conclusion \textbf{vraie}.

\begin{propriete}[Validité et Vérité]
\begin{itemize}
    \item \textbf{Validité (forme) :} Un raisonnement est valide \textbf{si et seulement s'il} est impossible que les prémisses soient vraies et la conclusion fausse. La conclusion \emph{découle nécessairement} des prémisses.
    \item \textbf{Vérité (contenu) :} Une proposition est vraie si elle \textbf{correspond au réel} (adéquation intellectus et rei). Un raisonnement est \emph{solide} (ou \emph{sain}) s'il est \textbf{valide} \textbf{ET} s'il a des \textbf{prémisses vraies}.
\end{itemize}
\end{propriete}

\begin{attention}[Piège majeur : Confondre validité et vérité]
\textbf{Erreur fréquente :} Dire « Ce raisonnement est faux » pour signifier « Il est invalide ».\\
\textbf{Correction :} Un raisonnement n'est pas « vrai » ou « faux », il est \textbf{valide} ou \textbf{invalide}. Seules les \emph{propositions} (prémisses, conclusion) sont vraies ou fausses.\\
\textbf{Exemple type au Bac :} On vous donne un syllogisme valide avec une prémisse fausse. La question : « Ce raisonnement est-il correct ? » Réponse attendue : « Il est \textbf{valide} (forme correcte) mais \textbf{non solide} car la prémisse mineure/majeure est fausse. »
\end{attention}

\begin{popfigure}
\begin{center}
\begin{tabularx}{\linewidth}{|l|X|X|}
\hline
\rowcolor{popblueL}
\textbf{Critère} & \textbf{Validité (Forme)} & \textbf{Vérité (Contenu)} \\ \hline
\textbf{Question} & La conclusion suit-elle logiquement des prémisses ? & Les prémisses et la conclusion correspondent-elles au réel ? \\ \hline
\textbf{Dépend de} & Structure logique uniquement (règles d'inférence) & Faits, observations, définitions, contexte \\ \hline
\textbf{Exemple 1 (Cameroun)} & 
\textbf{Prémisse 1 :} Tous les élèves de l'ENIEG de Douala portent l'uniforme kaki. \\
\textbf{Prémisse 2 :} Ngo Bassa est élève à l'ENIEG de Douala. \\
\textbf{Conclusion :} Donc Ngo Bassa porte l'uniforme kaki. \\
$\rightarrow$ \textbf{Valide} (Modus Ponens / Barbara). & 
Si Ngo Bassa est en réalité au Lycée de New-Bell (prémisse 2 fausse), la conclusion est fausse dans le réel, mais le raisonnement \textbf{reste valide}. \\ \hline
\textbf{Exemple 2 (Cameroun)} & 
\textbf{Prémisse 1 :} Yaoundé est la capitale politique. \\
\textbf{Prémisse 2 :} Douala est la capitale économique. \\
\textbf{Conclusion :} Donc le Mont Cameroun est un volcan actif. \\
$\rightarrow$ \textbf{Invalide} (non sequitur : conclusion sans lien logique). & 
Les trois propositions sont \textbf{vraies} dans le réel, mais l'enchaînement est \textbf{absurde logiquement}. \\ \hline
\end{tabularx}
\end{center}
\legende{Tableau comparatif : Validité (forme) vs Vérité (contenu) avec exemples camerounais.}
\end{popfigure}

\begin{exemplebox}[Le cas du commerçant de Mokolo (reprise de l'accroche)]
\textbf{Raisonnement du commerçant :}
\begin{itemize}
    \item P1 (Majeure) : Tous les clients qui marchandent achètent moins cher. (Forme : $\forall x (M(x) \to A(x))$)
    \item P2 (Mineure) : Toi, tu as marchandé. ($M(toi)$)
    \item C (Conclusion) : Donc tu as payé moins cher. ($A(toi)$)
\end{itemize}
\textbf{Analyse :}
\begin{enumerate}
    \item \textbf{Validité :} La forme est \textbf{valide} (Modus Ponens / Syllogisme Barbara). Si P1 et P2 sont vraies, C est nécessairement vraie.
    \item \textbf{Vérité des prémisses :} P2 est vraie (l'élève a marchandé). P1 est \textbf{fausse} : « Tous les clients... » est une universalité fausse (contre-exemple : l'élève a marchandé et payé le même prix).
    \item \textbf{Statut :} Raisonnement \textbf{valide mais non solide} (prémisse majeure fausse).
    \item \textbf{Réplique de l'élève :} Il fournit un \textbf{contre-exemple} à la prémisse majeure, prouvant sa fausseté. Il ne conteste pas la logique, il conteste le fait.
\end{enumerate}
\end{exemplebox}

\courssub{Méthode : Évaluer un discours (démarche en 3 étapes)}

\begin{methode}[Analyser et critiquer un raisonnement]
\begin{enumerate}
    \item \textbf{Identifier la conclusion :} Que veut-on me faire admettre ? (Cherchez les marqueurs : « donc », « par conséquent », « c'est pourquoi », « il faut que... »).
    \item \textbf{Relever les prémisses (explicites et implicites) :} Sur quels arguments repose la conclusion ? (Marqueurs : « car », « parce que », « puisque », « vu que... »). Attention aux \emph{enthymèmes} (prémisses cachées, « évidentes » pour l'orateur).
    \item \textbf{Vérifier dans l'ordre :}
        \begin{itemize}
            \item \textbf{Validité :} La forme est-elle correcte ? (Règles du syllogisme, tableaux de vérité, déduction naturelle). Si invalide $\to$ \textbf{Rejet immédiat} (sophisme formel).
            \item \textbf{Vérité des prémisses :} Si valide, les prémisses sont-elles exactes ? (Vérification factuelle, sources, définitions). Si prémisse fausse $\to$ \textbf{Raisonnement non solide}.
        \end{itemize}
\end{enumerate}
\textbf{Résultat :} Un raisonnement \textbf{solide} = Valide + Prémisses vraies $\to$ Conclusion \textbf{nécessairement vraie}.
\end{methode}

\begin{exoresolu}[Application de la méthode : Une rumeur WhatsApp]
\textbf{Message reçu :} « URGENT : Le Ministère de la Santé interdit la vente de \emph{foléré} (bissap) dans les buvettes. Cause : colorant cancérigène détecté. Partagez à tous vos contacts ! »
\textbf{Analyse :}
\begin{enumerate}
    \item \textbf{Conclusion :} « Il ne faut plus boire de foléré en buvette / Le Ministère l'a interdit. »
    \item \textbf{Prémisses :} P1 : « Le Ministère a émis un communiqué d'interdiction. » P2 (implicite) : « Ce communiqué est authentique et actuel. » P3 (implicite) : « Le colorant incriminé est présent dans tout le foléré vendu. »
    \item \textbf{Validité :} Si P1, P2, P3 vraies $\to$ Conclusion valide (Modus Ponens).
    \item \textbf{Vérité :} Recherche sur le site officiel \texttt{minsante.cm} ou page Facebook vérifiée $\to$ Aucun communiqué. P1 fausse. P3 fausse (généralisation).
    \item \textbf{Verdict :} \textbf{Fake news} (prémisses fausses). Ne pas partager. Appliquer la méthode : \emph{Source ? Date ? Preuve ?}
\end{enumerate}
\end{exoresolu}

\courssub{Synthèse de la section}

\begin{aretenir}[Ce qu'il faut retenir de la section 1]
\begin{itemize}
    \item \textbf{Logique} = Science des lois formelles de la pensée correcte (validité), née du \emph{logos} grec (raison/discours).
    \item \textbf{Aristote} = Fondateur (l'Organon) ; \textbf{Logique moderne} (Frege, Russell) = Formalisation mathématique (complément Bac).
    \item \textbf{Logique formelle} (forme/valide) $\neq$ \textbf{Logique matérielle} (contenu/vrai).
    \item \textbf{Importance :} Théorique (fondement sciences/maths) + Pratique (argumentation, détection manipulations, vie citoyenne).
    \item \textbf{Distinction cruciale :} \textbf{Validité} (structure correcte) $\neq$ \textbf{Vérité} (accord au réel). Un raisonnement valide peut avoir une conclusion fausse (prémisses fausses). Un raisonnement invalide peut avoir une conclusion vraie (par hasard).
    \item \textbf{Méthode d'évaluation :} 1) Conclusion, 2) Prémisses, 3) Validité puis Vérité.
\end{itemize}
\end{aretenir}

\coursec{Le concept}

Après avoir précisé ce qu'est la logique et pourquoi elle est indispensable pour structurer notre pensée et nos arguments, nous devons maintenant examiner l'outil premier de cette pensée : \textbf{le concept}. C'est la « brique » élémentaire avec laquelle nous construisons nos jugements et nos raisonnements. Sans maîtrise des concepts, point de logique rigoureuse possible.

\courssub{Définition et nature du concept}

\begin{definition}[Le concept]
Le \cle{concept} (aussi appelé \cle{idée} ou \cle{notion}) est la \textbf{représentation mentale générale et abstraite} d'un objet ou d'une classe d'objets. Il saisit ce qui est commun et essentiel à plusieurs individus pour les réunir en une même classe. Il s'exprime dans le langage par un \cle{terme} (un mot ou une locution).
\end{definition}

Contrairement à l'image sensible (qui est singulière, concrète et liée à la perception immédiate d'un objet particulier : \emph{ce} manguier devant chez moi, \emph{ce} taxi jaune), le concept porte sur l'universel. L'image est toujours \emph{ceci} ou \emph{cela} ; le concept est \emph{l'homme}, \emph{l'arbre}, \emph{la justice}, \emph{le triangle}.

\begin{exemplebox}[Image vs Concept]
\begin{itemize}
    \item \textbf{Image :} Je vois le proviseur de mon lycée, M. Mbarga, avec sa cravate bleue, sa taille grande, sa voix grave. Cette représentation est riche, colorée, singulière.
    \item \textbf{Concept :} « Proviseur ». Ce concept ne retient que les caractères \emph{essentiels} : « personne physique », « chef d'établissement », « autorité administrative et pédagogique ». Il s'applique à M. Mbarga, mais aussi à Mme Ngo à Yaoundé, à M. Kouam à Bafoussam, et à tous les proviseurs passés, présents et futurs.
\end{itemize}
\end{exemplebox}

\courssub{Les deux dimensions du concept : compréhension et extension}

Tout concept possède deux dimensions inséparables qui en déterminent la portée logique.

\begin{propriete}[Compréhension et Extension]
\begin{itemize}
    \item La \cle{compréhension} (ou intension) d'un concept est l'\textbf{ensemble des caractères essentiels} (attributs, propriétés) qui définissent ce concept. C'est le « contenu » de la définition. \emph{Exemple :} La compréhension de « triangle » = « figure plane » + « trois côtés » + « trois angles ».
    \item L'\cle{extension} d'un concept est l'\textbf{ensemble des objets réels ou possibles} auxquels ce concept s'applique (les individus qui tombent sous le concept). C'est la « portée » du concept. \emph{Exemple :} L'extension de « triangle » = tous les triangles existants ou pensables (équilatéraux, isocèles, rectangles, etc.).
\end{itemize}
\end{propriete}

\begin{aretenir}[Lien entre les deux]
La compréhension détermine l'extension : ce sont les caractères définis dans la compréhension qui permettent de reconnaître les objets appartenant à l'extension. Réciproquement, connaître l'extension aide à préciser la compréhension.
\end{aretenir}

\courssub{La loi de réciprocité (ou loi inverse)}

Il existe une relation inverse stricte entre la compréhension et l'extension d'un concept : \textbf{plus la compréhension est riche (beaucoup de caractères), plus l'extension est restreinte (peu d'objets), et vice-versa.}

\begin{propriete}[Loi de réciprocité]
\[
\text{Compréhension} \uparrow \quad \Longleftrightarrow \quad \text{Extension} \downarrow
\]
\[
\text{Compréhension} \downarrow \quad \Longleftrightarrow \quad \text{Extension} \uparrow
\]
\end{propriete}

\begin{popfigure}
\begin{tikzpicture}[scale=0.9, every node/.style={font=\small}]
% Styles
\tikzset{
    venn/.style={draw=popblue, thick, fill=popblue!10, ellipse},
    label/.style={font=\bfseries\small, text=popdark},
    count/.style={font=\small, text=popmuted, anchor=north},
}
% Ellipses emboîtées (centrées approximativement)
% 1. Homme (le plus grand)
\draw[venn, fill=popblue!5] (0,0) ellipse ({4.5cm} and {3.5cm});
\node[label] at (0, 2.8) {HOMME};
\node[count] at (0, 2.2) {Extension : $\sim$ 8 milliards};
% 2. Africain
\draw[venn, fill=popgreen!10, draw=popgreen] (0,-0.2) ellipse ({3.2cm} and {2.4cm});
\node[label, text=popgreen!70!black] at (0, 1.6) {AFRICAIN};
\node[count] at (0, 1.0) {Extension : $\sim$ 1,4 milliard};
% 3. Camerounais
\draw[venn, fill=poporange!10, draw=poporange] (0,-0.3) ellipse ({2.0cm} and {1.5cm});
\node[label, text=poporange!80!black] at (0, 0.7) {CAMEROUNAIS};
\node[count] at (0, 0.2) {Extension : $\sim$ 28 millions};
% 4. Lycéen camerounais
\draw[venn, fill=poppurple!10, draw=poppurple] (0,-0.35) ellipse ({1.1cm} and {0.8cm});
\node[label, text=poppurple!80!black] at (0, 0.1) {LYCÉEN CAMEROUNAIS};
\node[count] at (0, -0.3) {Extension : $\sim$ 1,5 million};
% Flèches explicatives à droite
\draw[popgray, thick, <->] (5.5, 2.5) -- (5.5, -2.5);
\node[align=center, text=popdark] at (7.2, 0) {Extension\\DIMINUE $\downarrow$};
\draw[popgray, thick, <->] (5.5, -2.8) -- (5.5, -3.5); % dummy for alignment
\node[align=center, text=popdark] at (7.2, -3.2) {Compréhension\\AUGMENTE $\uparrow$};
% Légende des caractères ajoutés (Compréhension)
\node[align=left, anchor=west, text=popdark] at (5.8, 2.0) {\textbf{+ Caractères :}};
\node[align=left, anchor=west, font=\small] at (5.8, 1.5) {• Être vivant};
\node[align=left, anchor=west, font=\small] at (5.8, 1.1) {• Rationalité};
\node[align=left, anchor=west, font=\small, text=popgreen!70!black] at (5.8, 0.6) {+ Nationalité africaine};
\node[align=left, anchor=west, font=\small, text=poporange!80!black] at (5.8, 0.1) {+ Nationalité camerounaise};
\node[align=left, anchor=west, font=\small, text=poppurple!80!black] at (5.8, -0.4) {+ Statut d'élève au secondaire};
\end{tikzpicture}
\legende{Diagramme de Venn illustrant la loi de réciprocité : en descendant dans la hiérarchie des concepts (Homme $\to$ Africain $\to$ Camerounais $\to$ Lycéen camerounais), la compréhension s'enrichit (flèche verte montante) tandis que l'extension se rétrécit (flèche rouge descendante).}
\end{popfigure}

\begin{exemplebox}[Illustration concrète : L'élève]
Prenons le concept \cle{élève} dans un lycée de Douala (ex. : Lycée de New-Bell, effectif 1 200 élèves).
\begin{itemize}
    \item \textbf{Concept « Élève » :} Compréhension = « Personne inscrite dans un établissement d'enseignement ». Extension = 1 200 individus.
    \item \textbf{Concept « Élève de Terminale » :} Compréhension = « Élève » + « Niveau Terminale ». Extension = 300 individus (environ).
    \item \textbf{Concept « Élève de Terminale C » :} Compréhension = « Élève de Terminale » + « Série C (Maths-Physique) ». Extension = 45 individus.
    \item \textbf{Concept « Élève de Terminale C, délégué de classe » :} Compréhension = « Élève de Terminale C » + « Élu délégué ». Extension = 1 individu (concept \textbf{singulier}).
\end{itemize}
À chaque ajout de caractère (augmentation de la compréhension), l'extension diminue drastiquement.
\end{exemplebox}

\courssub{Classification des concepts}

On distingue les concepts selon leur extension et leur nature.

\begin{definition}[Types de concepts selon l'extension]
\begin{itemize}
    \item \textbf{Concept singulier :} Extension = 1 seul objet. \emph{Ex. :} « Le Mont Cameroun », « Le Président de la République en exercice », « Mon acte de naissance n° 12345 ».
    \item \textbf{Concept général :} Extension > 1 (pluralité d'objets). \emph{Ex. :} « Montagne », « Président », « Acte de naissance ».
    \item \textbf{Concept vide (ou nul) :} Extension = 0 (aucun objet réel ne lui correspond). \emph{Ex. :} « Licorne », « Cercle carré », « Homme immortel » (au sens biologique).
\end{itemize}
\end{definition}

\begin{definition}[Types de concepts selon la nature]
\begin{itemize}
    \item \textbf{Concept concret :} Désigne un être ou un objet substantiel, existant par lui-même. \emph{Ex. :} « Table », « Élève », « Fleuve Wouri ».
    \item \textbf{Concept abstrait :} Désigne une qualité, une relation, un état, une action, considérée à part de son sujet. \emph{Ex. :} « Dureté », « Fraternité », « Indépendance », « Corruption », « Développement ».
\end{itemize}
\end{definition}

\begin{definition}[Termes univoques et équivoques]
\begin{itemize}
    \item \textbf{Terme univoque :} Un seul concept pour un terme (sens unique dans un contexte donné). \emph{Ex. :} « Triangle » (en géométrie), « Photosynthèse » (en biologie).
    \item \textbf{Terme équivoque :} Plusieurs concepts distincts pour un même terme (polysémie). Le sens dépend du contexte. \emph{Ex. :} « Banque » (siège financier / bord de rivière / série de batteries), « Vol » (action de voler / envol d'oiseau / vol d'avion), « Manche » (partie vêtement / manche d'outil / manche de tennis).
\end{itemize}
\end{definition}

\begin{attention}[Piège : L'équivoque dans l'argumentation]
Utiliser un terme équivoque sans en préciser le sens entraîne le \textbf{sophisme d'équivoque} (ex. : « La loi est dure, mais c'est la loi. Or, la loi de la gravité est une loi. Donc la loi de la gravité est dure. » $\to$ confusion entre « loi juridique » et « loi physique »). \textbf{Toujours définir ses termes} au début d'une dissertation ou d'un commentaire.
\end{attention}

\courssub{Les opérations sur le concept : Définition et Division}

La logique classique identifie deux opérations fondamentales pour clarifier et organiser les concepts : la \cle{définition} (qui explicite la compréhension) et la \cle{division} (qui parcourt l'extension).

\paragraph{1. La définition : expliciter la compréhension}

\begin{methode}[Règle de la définition par genre prochain et différence spécifique (Aristote / Porphyre)]
Une bonne définition (essentielle) suit deux étapes :
\begin{enumerate}
    \item Trouver le \cle{genre prochain} : la classe la plus proche, la plus large, qui contient l'objet à définir (le « groupe »).
    \item Trouver la \cle{différence spécifique} : le caractère propre, distinctif, qui sépare l'objet de tous les autres membres du genre (la « signature »).
\end{enumerate}
\textbf{Formule :} \textit{Concept à définir = Genre prochain + Différence spécifique.}
\end{methode}

\begin{exemplebox}[Exemples de définitions canoniques]
\begin{itemize}
    \item \textbf{L'homme} = \textit{Animal} (genre) + \textit{Rationnel} (différence spécifique).
    \item \textbf{Le triangle} = \textit{Figure plane} (genre) + \textit{À trois côtés} (différence spécifique).
    \item \textbf{Le lycée} = \textit{Établissement d'enseignement} (genre) + \textit{Du second cycle} (différence spécifique).
    \item \textbf{La corruption} = \textit{Détournement} (genre) + \textit{De fonds ou biens publics par un dépositaire de l'autorité} (différence spécifique).
\end{itemize}
\end{exemplebox}

\begin{propriete}[Règles de la bonne définition (Canons)]
Pour être valide, une définition doit respecter quatre règles :
\begin{enumerate}
    \item \textbf{Ni trop large, ni trop étroite :} Elle doit s'appliquer à \emph{tous} les objets du concept (extension exacte) et à \emph{eux seuls}.
        \begin{itemize}
            \item \emph{Trop large :} « L'homme est un animal qui marche. » (Les poules marchent aussi).
            \item \emph{Trop étroite :} « Le triangle est une figure à trois côtés égaux. » (Exclut les triangles non équilatéraux).
        \end{itemize}
    \item \textbf{Non circulaire :} Le défini ne doit pas réapparaître dans la définition (pas de \textit{definitio circulus}).
        \\ \emph{Exemple fautif :} « Un menteur est quelqu'un qui ment. » / « La corruption, c'est quand on est corrompu. »
    \item \textbf{Affirmative :} Privilégier ce que la chose \emph{est} plutôt que ce qu'elle \emph{n'est pas} (sauf pour les concepts négatifs par nature comme « aveugle », « vide », « injustice »).
    \item \textbf{Claire et intelligible :} Utiliser des termes plus connus que le concept défini.
\end{enumerate}
\end{propriete}


\coursec{Le jugement}

Dans la section précédente, nous avons étudié le \cle{concept}, cette représentation générale et abstraite que l'esprit se forme des choses : « homme », « pays », « justice », « élève ». Mais un concept, pris isolément, ne dit encore rien. Penser « élève » ou « sérieux » séparément, ce n'est pas encore penser véritablement : c'est seulement disposer des matériaux de la pensée. Pour que la pensée devienne effective, il faut \emph{relier} ces concepts entre eux, dire quelque chose de quelque chose : « les élèves sont sérieux ». Cette opération de liaison, c'est précisément le \cle{jugement}, deuxième opération fondamentale de l'esprit après la conception et avant le raisonnement. C'est lui que nous étudions maintenant.

\courssub{Définition et structure du jugement}

\begin{definition}[Le jugement]
Le \cle{jugement} est l'opération de l'esprit par laquelle nous \emph{unissons} ou \emph{séparons} deux concepts, en \textbf{affirmant} ou en \textbf{niant} un prédicat d'un sujet. Extérieurement, le jugement s'exprime par une \cle{proposition}, c'est-à-dire un énoncé susceptible d'être vrai ou faux.
\end{definition}

L'intuition est simple : quand je dis « le Cameroun est un pays bilingue », mon esprit prend le concept « Cameroun » et le concept « pays bilingue », et il déclare que le second convient au premier. Il \emph{juge} : il prononce un verdict, comme un juge qui tranche. À l'inverse, quand je dis « l'huile n'est pas soluble dans l'eau », mon esprit \emph{sépare} deux concepts : il refuse d'attribuer la solubilité dans l'eau à l'huile. Affirmer et nier sont donc les deux gestes fondamentaux du jugement.

\begin{propriete}[Structure du jugement]
Tout jugement comporte trois éléments :
\begin{enumerate}
\item le \cle{sujet} (noté $S$) : le concept dont on parle, celui dont on affirme ou nie quelque chose ;
\item le \cle{prédicat} (noté $P$) : le concept que l'on attribue ou que l'on refuse au sujet ;
\item la \cle{copule} : le verbe « être » (ou sa négation « n'est pas ») qui unit ou sépare le sujet et le prédicat.
\end{enumerate}
Schéma général : $S$ \; \textbf{est / n'est pas} \; $P$.
\end{propriete}

\begin{popfigure}
\begin{center}
\begin{tikzpicture}[font=\small]
\node[draw=popblue, fill=popblueL, rounded corners, minimum width=3.2cm, minimum height=1cm, align=center] (S) at (0,0){\textbf{Sujet ($S$)}\\ « Le Cameroun »};
\node[draw=popgold, fill=popgoldL, rounded corners, minimum width=2.6cm, minimum height=1cm, align=center] (C) at (5,0){\textbf{Copule}\\ « est »};
\node[draw=popgreen, fill=popgreenL, rounded corners, minimum width=3.4cm, minimum height=1cm, align=center] (P) at (10,0){\textbf{Prédicat ($P$)}\\ « un pays bilingue »};
\draw[-{Stealth}, thick, popdark] (S) -- (C);
\draw[-{Stealth}, thick, popdark] (C) -- (P);
\node[below=0.35cm of C, text=popmuted, align=center] {affirmation : union de $S$ et $P$};
\end{tikzpicture}
\end{center}
\legende{Structure du jugement analysée sur la phrase « Le Cameroun est un pays bilingue » : le sujet (ce dont on parle), la copule (le lien « est ») et le prédicat (ce qu'on en dit).}
\end{popfigure}

\begin{attention}[La copule n'est pas le verbe « être » au sens ordinaire]
Dans « l'élève \emph{travaille} », il n'y a pas de verbe « être » apparent. La logique classique ramène néanmoins tout jugement à la forme canonique : « l'élève \emph{est} travaillant ». De même, « le taxi roule vite » devient « le taxi \emph{est} roulant vite ». Cette reformulation permet de repérer clairement le sujet et le prédicat. Autre piège : ne pas confondre la proposition (énoncé du jugement) avec une question, un ordre ou un souhait, qui n'affirment rien et ne sont ni vrais ni faux.
\end{attention}

\courssub{Classification des jugements selon la qualité et la quantité}

Pour classer les jugements, la logique classique utilise deux critères croisés.

\textbf{Selon la qualité} (la manière dont la copule lie $S$ et $P$), un jugement est :
\begin{itemize}
\item \cle{affirmatif} : il \emph{unit} le prédicat au sujet (« les élèves sont attentifs ») ;
\item \cle{négatif} : il \emph{sépare} le prédicat du sujet (« les élèves ne sont pas attentifs »).
\end{itemize}

\textbf{Selon la quantité} (l'étendue du sujet), un jugement est :
\begin{itemize}
\item \cle{universel} : le prédicat est affirmé ou nié de \emph{tous} les individus désignés par le sujet (« \emph{tous} les élèves sont présents ») ;
\item \cle{particulier} : le prédicat n'est affirmé ou nié que d'\emph{une partie} des individus (« \emph{certains} élèves sont présents »).
\end{itemize}

En croisant ces deux critères, on obtient les quatre types fondamentaux de propositions, que la tradition logique, depuis \textbf{Aristote} (Antiquité grecque, fondateur de la logique), désigne par les lettres \textbf{A}, \textbf{E}, \textbf{I}, \textbf{O} (voyelles tirées des mots latins \emph{affirmo} : j'affirme, et \emph{nego} : je nie).

\begin{aretenir}[Les quatre types de propositions aristotéliciennes]
\begin{center}
\begin{tabularx}{\linewidth}{|c|l|X|l|}
\hline
\rowcolor{popblueL} \textbf{Type} & \textbf{Nom} & \textbf{Forme} & \textbf{Exemple} \\
\hline
\textbf{A} & Universelle affirmative & Tout $S$ est $P$ & Tous les élèves sont ponctuels. \\
\hline
\textbf{E} & Universelle négative & Nul $S$ n'est $P$ & Aucun élève sérieux ne triche. \\
\hline
\textbf{I} & Particulière affirmative & Quelque $S$ est $P$ & Certains élèves sont sportifs. \\
\hline
\textbf{O} & Particulière négative & Quelque $S$ n'est pas $P$ & Certains élèves ne sont pas sportifs. \\
\hline
\end{tabularx}
\end{center}
\end{aretenir}

\begin{methode}[Analyser une proposition en quatre étapes]
\begin{enumerate}
\item \textbf{Reformuler} si nécessaire la phrase sous la forme canonique « $S$ est / n'est pas $P$ ».
\item \textbf{Repérer} le sujet $S$, la copule et le prédicat $P$.
\item \textbf{Déterminer la qualité} : la copule affirme-t-elle (affirmatif) ou nie-t-elle (négatif) ?
\item \textbf{Déterminer la quantité} : le sujet est-il pris dans toute son extension (« tous », « aucun ») ou partiellement (« certains », « quelques ») ? Conclure en classant la proposition en A, E, I ou O.
\end{enumerate}
\end{methode}

\begin{exoresolu}[Classer des propositions de la vie courante]
Énoncé. Classer en A, E, I ou O les propositions suivantes : (1) « Tous les commerçants du marché Mokolo paient leurs impôts » ; (2) « Aucun baobab ne pousse sur les rives de l'estuaire du Wouri » ; (3) « Certains lycéens de Douala pratiquent le football » ; (4) « Certains élèves de Première technique ne révisent pas leurs leçons ».
\tcblower
Solution détaillée. (1) La copule affirme (« paient » = « sont payants ») et le sujet est pris universellement (« tous ») : c'est une proposition de type \textbf{A}. (2) La copule nie et le sujet est universel (« aucun ») : type \textbf{E}. (3) La copule affirme et le sujet est particulier (« certains ») : type \textbf{I}. (4) La copule nie (« ne révisent pas ») et le sujet est particulier : type \textbf{O}. On remarque que (1) et (4) portent sur des sujets voisins mais diffèrent à la fois par la quantité et par la qualité : ce sont des propositions de nature logique totalement différente.
\end{exoresolu}

\courssub{Le carré logique d'opposition}

Les quatre types de propositions, lorsqu'ils portent sur les \emph{mêmes} sujet et prédicat, entretiennent entre eux des relations précises que la tradition représente par le \cle{carré logique} (ou carré d'opposition).

\begin{popfigure}
\begin{center}
\begin{tikzpicture}[font=\small, every node/.style={align=center}]
\node[draw=popblue, fill=popblueL, rounded corners, minimum width=3.6cm, minimum height=1.2cm, align=center] (A) at (0,4){\textbf{A} : universelle affirmative\\ Tout $S$ est $P$};
\node[draw=poporange, fill=poporangeL, rounded corners, minimum width=3.6cm, minimum height=1.2cm, align=center] (E) at (9,4){\textbf{E} : universelle négative\\ Nul $S$ n'est $P$};
\node[draw=popgreen, fill=popgreenL, rounded corners, minimum width=3.6cm, minimum height=1.2cm, align=center] (I) at (0,0){\textbf{I} : particulière affirmative\\ Quelque $S$ est $P$};
\node[draw=poppurple, fill=poppurpleL, rounded corners, minimum width=3.6cm, minimum height=1.2cm, align=center] (O) at (9,0){\textbf{O} : particulière négative\\ Quelque $S$ n'est pas $P$};
\draw[{Stealth}-{Stealth}, very thick, popdark] (A) -- node[above]{contrariété} (E);
\draw[{Stealth}-{Stealth}, very thick, popdark] (I) -- node[above]{subcontrariété} (O);
\draw[-{Stealth}, very thick, popblue] (A) -- node[left, align=center]{subalternation\\ (A $\Rightarrow$ I)} (I);
\draw[-{Stealth}, very thick, poporange] (E) -- node[right, align=center]{subalternation\\ (E $\Rightarrow$ O)} (O);
\draw[{Stealth}-{Stealth}, very thick, poppink] (A) -- node[pos=0.5, sloped, above]{contradiction} (O);
\draw[{Stealth}-{Stealth}, very thick, poppink] (E) -- node[pos=0.5, sloped, above]{contradiction} (I);
\end{tikzpicture}
\end{center}
\legende{Le carré logique d'opposition : A et E sont contraires ; I et O sont subcontraires ; A et O, E et I sont contradictoires ; I découle de A et O découle de E (subalternation).}
\end{popfigure}

Définissons chacune de ces quatre relations.

\begin{definition}[Les quatre relations d'opposition]
\begin{itemize}
\item \textbf{Contrariété} (A--E) : deux jugements \emph{contraires} ne peuvent pas être vrais ensemble, mais ils peuvent être faux tous les deux. « Tous les élèves sont présents » et « aucun élève n'est présent » ne peuvent être vrais simultanément ; pourtant, si seulement la moitié de la classe est là, les deux sont faux.
\item \textbf{Subcontrariété} (I--O) : deux jugements \emph{subcontraires} ne peuvent pas être faux ensemble, mais ils peuvent être vrais tous les deux. « Certains élèves sont sportifs » et « certains élèves ne sont pas sportifs » peuvent être vrais en même temps.
\item \textbf{Contradiction} (A--O et E--I) : deux jugements \emph{contradictoires} ne peuvent être ni vrais ensemble, ni faux ensemble : si l'un est vrai, l'autre est nécessairement faux, et réciproquement.
\item \textbf{Subalternation} (A--I et E--O) : le jugement particulier (\emph{subalterne}) découle du jugement universel : si « tous les élèves sont présents » est vrai, alors « certains élèves sont présents » est vrai aussi. En revanche, la vérité du particulier n'entraîne pas celle de l'universel.
\end{itemize}
\end{definition}

\begin{propriete}[Principe des contradictoires]
Deux jugements contradictoires (A et O, ou bien E et I) ne peuvent être \textbf{ni vrais ensemble ni faux ensemble} : l'un est vrai exactement quand l'autre est faux. C'est la relation la plus forte du carré logique, et c'est elle qui fonde l'art de réfuter : pour détruire une thèse universelle, il suffit d'établir un seul cas contraire.
\end{propriete}

\begin{exemplebox}[La contradiction en action]
Considérons les deux propositions : « \emph{Aucun} élève sérieux ne triche » (type E) et « \emph{Quelques} élèves sérieux trichent » (type I). Elles sont contradictoires : si le surveillant de l'examen du Probatoire découvre un seul élève sérieux en train de tricher, la proposition E s'effondre immédiatement, et la proposition I devient vraie. Un seul contre-exemple suffit à renverser une universelle : voilà pourquoi, dans un débat, celui qui affirme « tous » ou « aucun » est toujours plus vulnérable que celui qui affirme « certains ».
\end{exemplebox}

\begin{savaistu}[Un outil vieux de plus de mille ans]
Le carré d'opposition, hérité d'Aristote et systématisé par les logiciens latins, a été utilisé pendant tout le Moyen Âge dans les universités européennes pour tester la cohérence des débats théologiques et juridiques. Les étudiants y apprenaient à repérer instantanément si deux thèses adverses étaient contraires ou contradictoires, afin de savoir si elles pouvaient être conciliées ou si l'une devait nécessairement être abandonnée. Cet outil est encore enseigné aujourd'hui dans les facultés de droit et de philosophie du monde entier.
\end{savaistu}

\courssub{Jugements de fait et jugements de valeur}

Une dernière distinction, capitale pour la vie intellectuelle et civique, concerne le \emph{contenu} du jugement.

\begin{definition}[Jugement de fait, jugement de valeur]
Un \cle{jugement de fait} constate une réalité observable et vérifiable : « Yaoundé est la capitale politique du Cameroun », « il a plu hier à Bafoussam ». Il est vrai ou faux, et on peut trancher par l'observation, la mesure ou le témoignage vérifiable. Un \cle{jugement de valeur} apprécie une réalité à partir d'une norme (le bien, le beau, le juste) : « ce film est magnifique », « la corruption est une injustice ». Il exprime une évaluation qui engage celui qui juge et qui peut être discutée.
\end{definition}

La distinction n'est pas toujours facile à opérer, car un même énoncé peut mêler les deux. Dire « cet élève a obtenu 15/20 » est un jugement de fait ; dire « cet élève est excellent » est un jugement de valeur qui s'appuie sur ce fait. Dans la vie courante, et particulièrement dans les débats publics, les confusions abondent.

\begin{attention}[Le piège de la généralisation abusive]
L'erreur la plus fréquente consiste à transformer abusivement un jugement \emph{particulier} en jugement \emph{universel}. De l'observation « \emph{certains} chauffeurs de taxi roulent vite » (I, légitime si on l'a constaté), on glisse à « \emph{tous} les chauffeurs roulent vite » (A, injustifié). De même, l'affirmation « \emph{tous} les politiciens sont corrompus » est un jugement universel (type A) présenté comme un fait, alors qu'il s'agit d'un jugement de valeur discutable : il suffirait d'un seul politicien honnête pour le réfuter, puisque son contradictoire « certains politiciens ne sont pas corrompus » (O) serait alors vrai. Penser logiquement, c'est donc aussi résister aux raccourcis de la rumeur, des préjugés et des réseaux sociaux.
\end{attention}

\begin{exoresolu}[Distinguer fait et valeur, particulier et universel]
Énoncé. Pour chacune des affirmations suivantes, dire s'il s'agit d'un jugement de fait ou de valeur, et préciser sa quantité : (1) « Douala est la capitale économique du Cameroun » ; (2) « Tous les jeunes d'aujourd'hui sont paresseux » ; (3) « Certains plats camerounais sont très épicés » ; (4) « Le ndolé est le meilleur plat du monde ».
\tcblower
Solution détaillée. (1) Jugement de \textbf{fait}, universel et singulier à la fois (il porte sur une ville unique) : vérifiable par la constitution économique du pays. (2) Jugement de \textbf{valeur} à prétention universelle (type A) : « paresseux » est une appréciation, et la généralisation à « tous » est abusive ; un seul jeune travailleur suffit à la réfuter. (3) Jugement de \textbf{fait} particulier (type I), vérifiable par l'expérience gustative ; sa forme prudente (« certains ») le rend difficilement réfutable. (4) Jugement de \textbf{valeur} : « le meilleur » exprime un goût personnel, non un fait mesurable ; on peut le défendre par des arguments, mais non le démontrer comme un fait.
\end{exoresolu}

\begin{aretenir}[L'essentiel de la section]
Le jugement est l'opération qui affirme ou nie un prédicat d'un sujet ; il s'exprime par une proposition de structure $S$ -- copule -- $P$. Selon la qualité, il est affirmatif ou négatif ; selon la quantité, universel ou particulier. Le croisement donne quatre types : A, E, I, O, reliés dans le carré logique par la contrariété, la subcontrariété, la contradiction et la subalternation. Enfin, il faut distinguer les jugements de fait (vérifiables) des jugements de valeur (appréciatifs), et se méfier des généralisations abusives. Ces acquis nous préparent à l'étape suivante : enchaîner plusieurs jugements pour produire un \emph{raisonnement}.
\end{aretenir}

\begin{exercice}[Pour s'entraîner]
1. Analyser (sujet, copule, prédicat, type A/E/I/O) les propositions : « Aucun sommet de l'Ouest camerounais n'est enneigé » ; « Certains proverbes sont des leçons de sagesse ». 2. Construire le contradictoire de « Tous les élèves de la classe ont réussi l'évaluation ». 3. Dire si les propositions « Tous les commerçants sont fermés le lundi » et « Aucun commerçant n'est fermé le lundi » peuvent être vraies ensemble, et si elles peuvent être fausses ensemble. Justifier à l'aide du carré logique.
\end{exercice}

\begin{corrige}[Exercice : pour s'entraîner]
1. « Aucun sommet de l'Ouest camerounais n'est enneigé » : $S$ = « sommet de l'Ouest camerounais », copule négative, $P$ = « enneigé » ; universelle négative, type \textbf{E}. « Certains proverbes sont des leçons de sagesse » : $S$ = « proverbes », copule affirmative, $P$ = « leçons de sagesse » ; particulière affirmative, type \textbf{I}. 2. Le contradictoire d'une universelle affirmative (A) est la particulière négative (O) : « Certains élèves de la classe n'ont pas réussi l'évaluation ». 3. Il s'agit de deux propositions contraires (A et E) : elles \emph{ne peuvent pas} être vraies ensemble, mais elles \emph{peuvent} être fausses toutes les deux, par exemple si certains commerçants sont fermés le lundi et d'autres ouverts.
\end{corrige}

\coursec{Le raisonnement}

Dans la section précédente, nous avons étudié le \cle{jugement}, cette opération par laquelle l'esprit affirme ou nie quelque chose de quelque chose : « le ciel est couvert », « Amina est admise ». Mais l'esprit ne s'arrête pas à énoncer des jugements isolés : il les \emph{enchaîne}. Lorsqu'un élève de Première technique constate que « tous les élèves qui ont la moyenne passent en classe supérieure » et que « lui-même a la moyenne », il en tire immédiatement : « donc je passe en classe supérieure ». Cette opération, qui fait naître un jugement nouveau à partir de jugements déjà posés, s'appelle le \cle{raisonnement}. C'est le sommet de l'activité logique : le concept fournit les matériaux, le jugement les assemble, le raisonnement les met en mouvement.

\courssub{Définition du raisonnement}

\begin{definition}[Le raisonnement]
Le \cle{raisonnement} est l'opération de l'esprit par laquelle on tire un jugement nouveau, appelé \cle{conclusion}, à partir d'un ou plusieurs jugements déjà posés, appelés \cle{prémisses}. Il est donc un enchaînement ordonné de jugements, reliés par des mots comme « donc », « par conséquent », « ainsi ».
\end{definition}

Prenons une observation de la vie courante. Un commerçant de Douala remarque : « chaque fois qu'il pleut fort, la rue est inondée » (première prémisse) ; « or aujourd'hui il pleut fort » (deuxième prémisse) ; il conclut : « donc la rue sera inondée ». La conclusion n'est pas une simple répétition des prémisses : c'est un jugement \emph{nouveau}, que l'esprit n'avait pas encore formulé, mais qui était en quelque sorte \emph{contenu} dans les prémisses.

Tout raisonnement possède donc une structure minimale : des \textbf{prémisses} (le point de départ), une \textbf{conclusion} (le point d'arrivée) et un \textbf{lien logique} qui les unit. La question centrale de la logique est alors : ce lien est-il \emph{solide} ? Autrement dit, le raisonnement est-il \cle{valide} ? On distingue, selon la nature de ce lien, trois grands types de raisonnement : la déduction, l'induction et l'analogie.

\courssub{Le raisonnement déductif : du général au particulier}

\begin{definition}[La déduction]
Le raisonnement \cle{déductif} est celui qui va du \textbf{général au particulier} : il applique une loi ou une vérité générale à un cas particulier. La conclusion est déjà \emph{contenue} dans les prémisses ; la déduction se contente de l'expliciter. C'est pourquoi, si les prémisses sont vraies, la conclusion s'impose avec une \cle{nécessité logique} : il est impossible de l'accepter fausse sans se contredire.
\end{definition}

L'intuition est simple : si je sais que \emph{tous} les fruits du manguier de la cour sont mûrs, alors \emph{ce} fruit-ci, que je viens de cueillir sur ce manguier, est nécessairement mûr. Je n'ai rien découvert de nouveau sur le monde ; j'ai seulement tiré du général ce qui s'y trouvait déjà. C'est le propre de la déduction : elle ne nous apprend rien que les prémisses ne contiennent, mais elle nous \emph{garantit} la vérité de la conclusion. Aristote, dans l'\emph{Organon}, fut le premier à en étudier la forme la plus célèbre : le syllogisme.

\courssub{Le syllogisme catégorique}

\begin{definition}[Le syllogisme]
Le \cle{syllogisme} catégorique est un raisonnement déductif composé de \textbf{deux prémisses} et d'\textbf{une conclusion}, mettant en jeu exactement \textbf{trois termes} :
\begin{itemize}
\item le \cle{grand terme}, attribut de la conclusion (il figure dans la \textbf{majeure}, la première prémisse) ;
\item le \cle{petit terme}, sujet de la conclusion (il figure dans la \textbf{mineure}, la deuxième prémisse) ;
\item le \cle{moyen terme}, qui figure dans les deux prémisses mais \emph{jamais} dans la conclusion : c'est lui qui fait le lien.
\end{itemize}
\end{definition}

L'exemple canonique, célèbre depuis Aristote, est le suivant :

\begin{exemplebox}[Le syllogisme de Socrate]
\textbf{Majeure} : Tous les hommes sont mortels.\\
\textbf{Mineure} : Or Socrate est un homme.\\
\textbf{Conclusion} : Donc Socrate est mortel.

\medskip
Analyse des termes : le \emph{moyen terme} est « homme » (il relie les deux prémisses et disparaît de la conclusion) ; le \emph{petit terme} est « Socrate » ; le \emph{grand terme} est « mortel ».
\end{exemplebox}

\begin{propriete}[Validité et vérité]
Dans un syllogisme \textbf{valide}, si les prémisses sont vraies, alors la conclusion est \emph{nécessairement} vraie. La validité concerne la \textbf{forme} du raisonnement (sa structure), non le contenu : un syllogisme peut être valide avec des prémisses fausses (la conclusion n'est alors pas garantie), mais un syllogisme \emph{invalide} ne garantit jamais sa conclusion, même si celle-ci est vraie par hasard.
\end{propriete}

\begin{popfigure}
\begin{center}
\begin{tikzpicture}[scale=1.0]
% Cercle des mortels
\draw[thick, popblue, fill=popblueL] (0,0) circle (2.6);
\node[popblue, font=\bfseries] at (-1.3,2.1) {Mortels};
% Cercle des hommes
\draw[thick, poporange, fill=poporangeL] (0.6,-0.3) circle (1.5);
\node[popdark, font=\bfseries] at (0.6,0.75) {Hommes};
% Socrate
\fill[popink] (0.6,-0.5) circle (2.2pt);
\node[popink, font=\bfseries, right] at (0.72,-0.55) {Socrate};
\end{tikzpicture}
\legende{Diagramme de Venn vérifiant le syllogisme : l'ensemble des hommes est inclus dans celui des mortels (majeure) ; Socrate appartient à l'ensemble des hommes (mineure) ; il appartient donc nécessairement à l'ensemble des mortels (conclusion).}
\end{center}
\end{popfigure}

Le diagramme rend la nécessité visible : une fois que le cercle « hommes » est entièrement dessiné à l'intérieur du cercle « mortels », tout point placé dans « hommes » se trouve \emph{forcément} dans « mortels ». La conclusion n'ajoute rien : elle était déjà là, enfermée dans les prémisses.

\courssub{Les règles de validité du syllogisme}

Pour qu'un syllogisme soit valide, il ne suffit pas qu'il ait l'air correct : il doit obéir à des règles précises.

\begin{aretenir}[Principales règles de validité]
\begin{enumerate}
\item Le syllogisme ne comporte que \textbf{trois termes}, chacun gardant le même sens du début à la fin (pas d'équivoque).
\item Le \textbf{moyen terme ne doit jamais figurer dans la conclusion} : son rôle est de servir de pont, puis de disparaître.
\item Le moyen terme doit être pris \textbf{au moins une fois universellement} (dans toute son extension : « \emph{tous} les hommes », « \emph{aucun} menteur ») : c'est la condition de la \emph{distribution} du moyen terme.
\item On ne peut tirer \textbf{aucune conclusion de deux prémisses négatives} (deux refus ne relient rien).
\item On ne peut tirer \textbf{aucune conclusion de deux prémisses particulières} (« certains... certains... donc... » ne conclut rien de nécessaire).
\item \textbf{Un terme non distribué dans les prémisses ne peut pas l'être dans la conclusion} (pas de procès illicite du majeur ni du mineur).
\item Si une prémisse est négative, la conclusion doit l'être ; si la conclusion est négative, une prémisse au moins doit l'être.
\end{enumerate}
\end{aretenir}

\begin{methode}[Vérifier la validité d'un syllogisme]
\begin{enumerate}
\item \textbf{Identifier les trois termes} : repérer la conclusion, en extraire le petit terme (sujet) et le grand terme (attribut) ; le terme restant est le moyen terme.
\item \textbf{Vérifier le rôle du moyen terme} : il doit apparaître dans les deux prémisses, jamais dans la conclusion, et être pris au moins une fois universellement (règle 3).
\item \textbf{Contrôler la distribution des termes extrêmes} : si le petit terme (sujet de la conclusion) ou le grand terme (attribut de la conclusion) sont pris universellement dans la conclusion, ils doivent l'être aussi dans la prémisse où ils figurent (règle 6).
\item \textbf{Contrôler les règles de forme} : pas de deux prémisses négatives (règle 4), pas de deux prémisses particulières (règle 5), pas d'équivoque (règle 1), cohérence quantité/qualité (règle 7).
\item \textbf{Tester avec un diagramme de Venn} : dessiner les ensembles correspondant aux prémisses, puis vérifier si la conclusion y est déjà contenue.
\end{enumerate}
\end{methode}

\begin{attention}[Le piège du moyen terme non distribué]
Erreur très fréquente : « Tous les B sont des C ; or A est un C ; donc A est un B ». Ce raisonnement est \textbf{invalide}, car le moyen terme « C » n'est jamais pris universellement. Exemple : « Tous les lions sont des mammifères ; or mon chat est un mammifère ; donc mon chat est un lion. » La conclusion est absurde alors que les prémisses sont vraies : la forme est fautive. Être dans le grand ensemble « mammifères » ne suffit pas à être dans le petit ensemble « lions ».
\end{attention}

\begin{attention}[Le piège du procès illicite (majeur ou mineur)]
Autre erreur classique : « Aucun homme n'est parfait ; or tous les saints sont des hommes ; donc aucun saint n'est parfait. » (Valide). Mais : « Aucun homme n'est parfait ; or Socrate est un homme ; donc Socrate n'est pas parfait. » (Valide). Contre-exemple invalide : « Tous les hommes sont mortels ; or Socrate est un homme ; donc tous les mortels sont des hommes. » Ici le grand terme « mortel » est distribué dans la conclusion (sujet d'une universelle) mais non distribué dans la majeure (attribut d'une affirmative). Le syllogisme est invalide (procès illicite du majeur).
\end{attention}

Appliquons maintenant la méthode à un exemple camerounais.

\begin{exoresolu}[Le cas d'Amina au concours]
\textbf{Énoncé.} Un camarade affirme : « Tous les élèves admissibles ont la moyenne ; or Amina a la moyenne au concours ; donc Amina est admissible. » Ce raisonnement est-il valide ?

\tcblower
\textbf{Solution.} Identifions les termes : petit terme = « Amina », grand terme = « admissible », moyen terme = « avoir la moyenne ». Le moyen terme figure bien dans les deux prémisses et absent de la conclusion : jusque-là, tout va bien. Mais est-il pris au moins une fois universellement ? Dans la majeure, « ont la moyenne » est l'attribut de « tous les élèves admissibles » : on ne parle pas de \emph{tous} ceux qui ont la moyenne, seulement des admissibles. Dans la mineure, on parle seulement d'Amina. Le moyen terme n'est donc pris \emph{nulle part} universellement. En effet, des élèves peuvent avoir la moyenne sans être admissibles (si la barre d'admissibilité est fixée à 12/20, par exemple). Le raisonnement a exactement la forme fautive « Tous les B sont des C ; A est un C ; donc A est un B ». \textbf{Conclusion : le syllogisme est invalide}, même si, par chance, Amina est réellement admissible.
\end{exoresolu}

\courssub{Le raisonnement inductif : du particulier au général}

La déduction part du général. Mais d'où vient le général lui-même ? Comment savons-nous que « tous les hommes sont mortels » ? Parce que nous avons observé une multitude de cas particuliers. C'est l'autre grand mouvement de la pensée : l'\cle{induction}.

\begin{definition}[L'induction]
Le raisonnement \cle{inductif} va du \textbf{particulier au général} : à partir de l'observation répétée de cas semblables, on formule une loi générale. Contrairement à la déduction, sa conclusion n'est pas nécessaire mais seulement \cle{probable} : elle dépasse ce que les prémisses contiennent. C'est le raisonnement propre des \textbf{sciences expérimentales}.
\end{definition}

\begin{exemplebox}[Une induction chiffrée]
Une enquête porte sur 200 ménages de Bafoussam : 180 d'entre eux, soit 90\,\%, utilisent le mobile money pour leurs paiements. On en conclut : « la grande majorité des ménages de Bafoussam utilise le mobile money ». Cette conclusion est \emph{vraisemblable}, mais non certaine : les 200 ménages enquêtés ne représentent qu'un échantillon ; peut-être les quartiers non enquêtés utilisent-ils moins ce service. L'induction gagne en force quand le nombre d'observations augmente et quand l'échantillon est varié, mais elle n'atteint jamais la nécessité absolue de la déduction.
\end{exemplebox}

Francis Bacon, puis John Stuart Mill, ont fait de l'induction la méthode par excellence de la science : observer, accumuler les faits, puis généraliser prudemment. Le savant camerounais qui teste une variété de maïs sur plusieurs parcelles de l'Ouest avant de la recommander à tous les paysans de la région raisonne inductivement.

\courssub{Le raisonnement par analogie : du particulier au particulier}

\begin{definition}[L'analogie]
Le raisonnement par \cle{analogie} va du \textbf{particulier au particulier} : sachant que deux choses se ressemblent sur certains points, on conclut qu'elles se ressemblent aussi sur un autre point. Sa conclusion est encore plus fragile que celle de l'induction : simple \cle{présomption}.
\end{definition}

Exemple : « Le lycée de la ville voisine a installé un club informatique et ses résultats au Probatoire ont progressé ; notre lycée lui ressemble ; donc un club informatique fera progresser nos résultats. » Ce raisonnement donne une \emph{hypothèse} intéressante, utile pour orienter une décision, mais il ne prouve rien : la ressemblance entre les deux lycées est peut-être superficielle (effectifs, encadrement, motivation différents). L'analogie est précieuse pour \emph{inventer} des hypothèses ; elle est dangereuse quand on la prend pour une preuve.

\begin{popfigure}
\begin{center}
\begin{tikzpicture}[scale=0.95]
% Déduction : entonnoir vers le bas
\draw[thick, popblue, fill=popblueL] (-5.2,2.2) -- (-1.4,2.2) -- (-3.3,0.4) -- cycle;
\node[popblue, font=\bfseries] at (-3.3,2.75) {D\'eduction};
\node[align=center, font=\small] at (-3.3,1.7) {G\'en\'eral\\ \emph{(tous les hommes...)}};
\node[align=center, font=\small] at (-3.3,0.0) {Particulier\\ \emph{(Socrate...)}};
\node[align=center, font=\small\itshape, popblue] at (-3.3,-0.8) {Conclusion n\'ecessaire};
% Induction : entonnoir vers le haut
\draw[thick, poporange, fill=poporangeL] (1.4,0.4) -- (5.2,0.4) -- (3.3,2.2) -- cycle;
\node[poporange, font=\bfseries] at (3.3,2.75) {Induction};
\node[align=center, font=\small] at (3.3,0.0) {Particuliers\\ \emph{(cas observ\'es)}};
\node[align=center, font=\small] at (3.3,1.55) {G\'en\'eral\\ \emph{(loi probable)}};
\node[align=center, font=\small\itshape, poporange] at (3.3,-0.8) {Conclusion probable};
\end{tikzpicture}
\legende{Les deux grands mouvements de la pensée : la déduction resserre le général vers le cas particulier avec nécessité ; l'induction s'élève des cas particuliers vers une loi générale, avec probabilité.}
\end{center}
\end{popfigure}

\begin{center}
\begin{tabularx}{\linewidth}{|l|X|X|X|}
\hline
\rowcolor{popblueL}
\textbf{Critère} & \textbf{Déduction} & \textbf{Induction} & \textbf{Analogie} \\
\hline
Point de départ & Vérité générale & Cas particuliers observés & Un cas particulier \\
\hline
Point d'arrivée & Cas particulier & Loi générale & Autre cas particulier \\
\hline
Force de la conclusion & Nécessaire (certaine si les prémisses sont vraies) & Probable & Simple présomption \\
\hline
Domaine privilégié & Mathématiques, logique & Sciences expérimentales & Hypothèses, vie pratique \\
\hline
\end{tabularx}
\end{center}

\courssub{Sophismes et paralogismes : les erreurs de raisonnement}

Un \cle{paralogisme} est un raisonnement faux commis de bonne foi ; un \cle{sophisme} est un raisonnement faux présenté volontairement pour tromper ou pour gagner un débat. Les deux violent les règles de la logique. En voici trois formes fréquentes, faciles à repérer dans les discussions quotidiennes et sur les réseaux sociaux.

\begin{itemize}
\item \textbf{Le faux syllogisme} : il imite la forme du syllogisme mais viole une règle de validité. Exemple : « Tous les candidats sérieux révisent la nuit ; or ce candidat révise la nuit ; donc il est sérieux. » (Moyen terme non distribué : beaucoup révisent la nuit sans méthode ni sérieux.)
\item \textbf{La pétition de principe} : on suppose déjà vraie, dans les prémisses, la chose qu'on prétend démontrer. Exemple : « Ce journal dit la vérité, car tout ce qu'il publie est vrai. » Le cercle est parfait... et vide.
\item \textbf{La généralisation hâtive} : induction abusive à partir de trop peu de cas. Exemple lu sur un réseau social : « Deux jeunes de mon quartier ont échoué au Probatoire en suivant des cours du soir ; donc les cours du soir ne servent à rien. » Deux cas ne fondent pas une loi générale.
\end{itemize}

\begin{attention}[Distinguer validité et vérité]
Ne jamais confondre « le raisonnement est valide » et « la conclusion est vraie ». Un raisonnement peut être valide avec des prémisses fausses (« Tous les Camerounais parlent chinois ; or Amina est camerounaise ; donc Amina parle chinois » : forme correcte, prémisse fausse, conclusion fausse). Inversement, une conclusion vraie peut sortir d'un raisonnement invalide. La logique juge d'abord la \emph{forme}.
\end{attention}

\begin{exercice}[À vous de raisonner]
\begin{enumerate}
\item « Aucun tricheur ne mérite le respect ; or certains élèves sont des tricheurs ; donc certains élèves ne méritent pas le respect. » Identifiez les trois termes et vérifiez la validité.
\item « J'ai vu trois taxis jaunes rouler vite ce matin ; donc tous les taxis jaunes roulent vite. » Quel type de raisonnement ? Quelle erreur ?
\item « Mon grand frère a réussi en travaillant seul ; je lui ressemble ; donc je réussirai en travaillant seul. » Quel type de raisonnement ? Quelle est sa force ?
\end{enumerate}
\end{exercice}

\begin{corrige}[Exercice]
\begin{enumerate}
\item Petit terme : « certains élèves » ; grand terme : « mériter le respect » ; moyen terme : « tricheur », pris universellement dans la majeure (« aucun ») et absent de la conclusion. Une seule prémisse négative, une seule particulière : les règles sont respectées. Le syllogisme est \textbf{valide} (mode Ferison, première figure).
\item Raisonnement \textbf{inductif} (du particulier au général), mais fautif : c'est une \textbf{généralisation hâtive} fondée sur trois cas seulement, observés un seul matin.
\item Raisonnement par \textbf{analogie} (du particulier au particulier par ressemblance). Sa conclusion n'est qu'une présomption : la ressemblance entre les deux frères ne garantit pas que la même méthode de travail convienne aux deux.
\end{enumerate}
\end{corrige}

\begin{aretenir}[L'essentiel du raisonnement]
Le raisonnement tire une conclusion de prémisses. La \textbf{déduction} va du général au particulier avec nécessité ; sa forme modèle est le \textbf{syllogisme} (deux prémisses, trois termes), soumis à des règles de validité strictes (distribution du moyen terme, pas de procès illicite, cohérence quantité/qualité). L'\textbf{induction} va du particulier au général avec probabilité ; l'\textbf{analogie} va du particulier au particulier par ressemblance. Sophismes et paralogismes (faux syllogisme, pétition de principe, généralisation hâtive) sont des raisonnements qui enfreignent ces règles : savoir les démasquer, c'est se protéger des manipulations dans les débats et sur les réseaux sociaux.
\end{aretenir}

\coursec{Les principes logiques de la pensée}

Après avoir analysé la structure du raisonnement — ses formes (déduction, induction), ses règles de validité et ses pièges (sophismes) —, nous devons nous interroger sur le socle même qui rend possible toute pensée cohérente. Pourquoi un raisonnement est-il valide ? Qu'est-ce qui garantit qu'une affirmation conserve son sens d'un bout à l'autre d'un discours ? La réponse se trouve dans les \textbf{principes logiques de la pensée} : des lois premières, évidentes par elles-mêmes, que l'esprit ne peut violer sans cesser de penser. Aristote, dans sa \emph{Métaphysique}, les présente comme les conditions \emph{sine qua non} de tout discours rationnel. Les maîtriser, c'est s'outiller pour penser juste, argumenter sans se contredire et détecter les incohérences dans les discours qui nous entourent — au tribunal, dans un contrat, dans un débat public ou dans une copie d'examen.

\courssub{La notion de principe : loi première et condition de la pensée}

\begin{definition}[Principe logique]
Un \textbf{principe logique} est une loi fondamentale, première et indémontrable, qui régit l'opération de l'esprit. Il ne se démontre pas : il s'impose comme une évidence immédiate (une \emph{intuition intellectuelle}) et sert de fondement à toute démonstration. Violenter un principe logique, c'est tomber dans l'absurde, c'est-à-dire produire un énoncé qui ne signifie plus rien.
\end{definition}

\begin{aretenir}[Trois principes classiques + un complément]
La logique classique (aristotélicienne) reconnaît trois principes suprêmes :
\begin{enumerate}
    \item Le \textbf{principe d'identité} : \og A est A \fg.
    \item Le \textbf{principe de non-contradiction} : \og A n'est pas non-A \fg.
    \item Le \textbf{principe du tiers exclu} : \og Ou A, ou non-A, pas de troisième \fg.
\end{enumerate}
Le \textbf{principe de raison suffisante} (Leibniz) est souvent adjoint comme complément à l'examen du Probatoire / Baccalauréat technique.
\end{aretenir}

Ces principes ne sont pas des \emph{observations} sur le monde (comme \og l'eau bout à 100°C \fg), mais des \emph{conditions} pour que nous puissions parler du monde. Ils concernent la \textbf{forme} de la pensée, non son \textbf{contenu} matériel. C'est pourquoi ils sont dits \emph{formels} et \emph{a priori}.

\courssub{Le principe d'identité : la stabilité du sens}

\begin{definition}[Principe d'identité]
\textbf{Tout ce qui est, est ce qu'il est.} Sous sa forme symbolique : \textbf{$A = A$} ou \textbf{$A$ est $A$}. Ce principe exige que, dans un même discours, un terme conserve une signification identique. Si le mot \og banque \fg\ désigne un établissement financier au début d'une phrase, il ne peut désigner un bord de rivière à la fin sans créer une équivoque.
\end{definition}

\begin{exemplebox}[Identité et équivoque]
\begin{itemize}
    \item \textbf{Correct :} \og Le triangle a trois côtés. Donc ce triangle-ci a trois côtés. \fg\ (Le concept \og triangle \fg\ garde son sens.)
    \item \textbf{Violation (équivoque) :} \og La loi est dure, mais c'est la loi. \fg\ (Ici \og loi \fg\ change de sens : loi morale / loi juridique. C'est un sophisme d'équivoque.)
\end{itemize}
\end{exemplebox}

\begin{attention}[Piège : confondre identité et permanence]
Le principe d'identité ne dit pas qu'une chose ne change \emph{jamais} dans le temps (un arbre grandit, une entreprise évolue). Il dit que \textbf{au moment où l'on pense ou l'on parle}, le terme utilisé doit désigner un objet déterminé, sans glissement de sens. \og L'entreprise X fait 10 millions de bénéfices \fg\ (en 2023) et \og L'entreprise X fait 2 millions de pertes \fg\ (en 2024) ne se contredisent pas : le sujet \og entreprise X \fg\ est identique, le prédicat porte sur des temps différents.
\end{attention}

\begin{experience}[Tester l'identité dans un contrat]
\textbf{Matériel :} Extrait d'un contrat de location (bail) où le terme \og charges \fg\ apparaît trois fois.
\textbf{Protocole :} Relever chaque occurrence et vérifier si la définition (liste des postes : eau, électricité, ordures, garde) est stable.
\textbf{Observation :} Si l'article 3 définit \og charges \fg\ comme \og eau + électricité \fg\ et l'article 12 y ajoute \og frais de syndic \fg\ sans avertissement, le principe d'identité est violé : le signataire croit signer pour A, il signe pour A+B.
\textbf{Interprétation :} La sécurité juridique repose sur l'identité rigoureuse des termes. Tout glissement est une faille exploitable.
\end{experience}

\courssub{Le principe de non-contradiction : l'interdiction de l'absurde}

\begin{definition}[Principe de non-contradiction]
\textbf{Il est impossible que la même chose soit et ne soit pas en même temps et sous le même rapport.} (Aristote, \emph{Métaphysique}, Γ, 3). Symboliquement : \textbf{$\neg(A \land \neg A)$}. On ne peut affirmer et nier simultanément le même attribut du même sujet, \emph{au même moment} et \emph{sous le même aspect}.
\end{definition}

\begin{savaistu}[Aristote et le \og plus certain des principes \fg]
Aristote écrit : \og Le principe de non-contradiction est le plus certain de tous les principes [...] car quiconque veut faire une démonstration doit l'admettre implicitement. \fg\ (Métaphysique, 1005b). Pour lui, le nier revient à ne rien dire du tout, à se réduire au silence ou au babillage.
\end{savaistu}

\begin{exemplebox}[Témoignage au commissariat — Douala, 15 mai, 14 h]
Un témoin déclare successivement :
\begin{enumerate}
    \item \og J'étais à Douala, au marché Congo, le 15 mai à 14 h précises. \fg
    \item \og Je n'étais pas à Douala le 15 mai à 14 h ; j'étais à Yaoundé. \fg
\end{enumerate}
Ces deux affirmations ne peuvent être \textbf{toutes deux vraies}. Elles violent le principe de non-contradiction : même sujet (le témoin), même attribut (être à Douala), même moment (15 mai 14 h), même rapport (présence physique). L'enquêteur n'a pas besoin d'aller sur le terrain pour savoir qu'au moins une est fausse : la contradiction \emph{logique} suffit à discréditer le témoignage.
\end{exemplebox}

\begin{attention}[Erreur fréquente : \og C'est vrai dans un certain sens \fg]
Dire \og Il fait chaud et il ne fait pas chaud \fg\ n'est une contradiction \textbf{que si} le rapport (lieu, moment, échelle, ressenti) est identique.
\begin{itemize}
    \item \og Il fait chaud à Maroua (40°C) et il ne fait pas chaud à Ebolowa (22°C) \fg\ : \textbf{pas de contradiction} (lieux différents).
    \item \og Cette eau est chaude (pour du thé) et cette eau n'est pas chaude (pour un bain) \fg\ : \textbf{pas de contradiction} (rapports différents).
    \item \og Il fait 30°C ici à 14 h et il ne fait pas 30°C ici à 14 h \fg\ : \textbf{contradiction formelle}.
\end{itemize}
La nuance exige de \textbf{préciser le rapport} (lieu, temps, critère), pas de tolérer l'incohérence.
\end{attention}

\courssub{Le principe du tiers exclu : l'alternative radicale}

\begin{definition}[Principe du tiers exclu]
\textbf{Entre A et non-A, il n'y a pas de troisième possibilité.} Toute proposition est \textbf{nécessairement vraie ou fausse} (valeur de vérité bivalente). Symboliquement : \textbf{$A \lor \neg A$}. Il n'existe pas d'état \og entre les deux \fg\ pour la vérité d'un énoncé assertorique (qui affirme quelque chose).
\end{definition}

\begin{popfigure}
\begin{tikzpicture}[scale=1.1]
    % Fond : deux demi-plans
    \fill[popblue!15] (-4,-2) rectangle (0,2);
    \fill[poporange!15] (0,-2) rectangle (4,2);
    % Ligne de séparation épaisse
    \draw[popdark, line width=2pt] (0,-2) -- (0,2);
    % Flèches indicatrices
    \draw[-{Stealth[length=3mm]}, popblue, thick] (-3,1.3) -- (-0.5,1.3) node[midway, above] {$A$};
    \draw[-{Stealth[length=3mm]}, poporange, thick] (0.5,1.3) -- (3,1.3) node[midway, above] {$\neg A$};
    % Étiquettes
    \node[popblue, font=\bfseries\large] at (-2,0) {VRAI};
    \node[poporange, font=\bfseries\large] at (2,0) {FAUX};
    \node[popdark, font=\bfseries, rotate=90] at (0, -2.5) {Ligne de partage exclue};
    % Exemple concret
    \node[align=center, font=\small, text width=3.5cm] at (-2,-2.7) {\og Camair-Co est la compagnie nationale \fg\ (VRAI)};
    \node[align=center, font=\small, text width=3.5cm] at (2,-2.7) {\og 2 + 2 = 5 \fg\ (FAUX)};
    % Pas de zone grise
    \draw[popmuted, dashed, thick] (-4,-2) rectangle (4,2);
\end{tikzpicture}
\legende{Le principe du tiers exclu : une frontière nette, sans zone grise. Toute proposition assertorique tombe d'un côté ou de l'autre.}
\end{popfigure}

\begin{exemplebox}[Tiers exclu et vie courante]
\begin{itemize}
    \item \og Paul a réussi le Probatoire. \fg\ $\rightarrow$ Vrai \textbf{ou} Faux. Pas \og un peu réussi \fg.
    \item \og Ce médicament guérit le paludisme. \fg\ $\rightarrow$ Vrai \textbf{ou} Faux (selon les essais cliniques). Pas \og ça marche un peu \fg\ comme valeur de vérité.
    \item \textbf{Attention :} Le tiers exclu s'applique aux \emph{propositions} (énoncés fermés), pas aux \emph{réalités graduées} (température, taille, opinion). \og Il fait tiède \fg\ n'est pas une proposition bivalente sans seuil défini.
\end{itemize}
\end{exemplebox}

\begin{attention}[Piège : confondre \og tiers exclu \fg\ et \og faux dilemme \fg]
Le principe du tiers exclu est une \textbf{loi logique} ($A$ ou $\neg A$). Le \textbf{faux dilemme} est un \textbf{sophisme} qui présente artificiellement deux alternatives alors qu'il en existe d'autres (\og Soit tu es pour nous, soit tu es contre nous \fg\ — on peut être neutre, critique, indifférent). Le tiers exclu s'applique à la \emph{valeur de vérité} d'une même proposition ; le faux dilemme manipule le \emph{contenu} des choix.
\end{attention}

\courssub{Complément Bac : le principe de raison suffisante (Leibniz)}

\begin{definition}[Principe de raison suffisante]
\textbf{Rien n'est sans raison pourquoi il est plutôt que non.} (Leibniz, \emph{Monadologie}, § 32). Tout fait, toute vérité, toute existence a une raison suffisante qui l'explique — même si nous ne la connaissons pas encore. Ce principe fonde l'exigence de \textbf{justification} : on ne peut accepter une affirmation comme vraie sans en donner le fondement (cause, motif, preuve).
\end{definition}

\begin{aretenir}[Différence avec les trois principes aristotéliciens]
\begin{itemize}
    \item Identité, non-contradiction, tiers exclu = \textbf{principes formels} (structure de la pensée).
    \item Raison suffisante = \textbf{principe matériel / ontologique} (lien entre la pensée et la réalité, exigence d'explication).
\end{itemize}
Au Probatoire / Baccalauréat technique, on peut vous demander de \textbf{citer} Leibniz et de \textbf{distinguier} ce principe des trois premiers.
\end{aretenir}

\begin{exemplebox}[Raison suffisante dans un exposé technique]
Un élève soutient : \og La résistance de ce béton est de 35 MPa. \fg
\begin{itemize}
    \item \textbf{Sans raison suffisante :} \og C'est comme ça. \fg\ (Inacceptable.)
    \item \textbf{Avec raison suffisante :} \og Selon la norme NC 234, l'essai de compression à 28 jours sur 3 éprouvettes donne 34,8 ; 35,2 ; 35,1 MPa. La moyenne est 35 MPa. \fg\ (Raison suffisante = protocole normé + mesures.)
\end{itemize}
\end{exemplebox}

\courssub{Tableau récapitulatif des trois principes classiques}

\begin{popfigure}
\begin{tikzpicture}
    % Tableau TikZ récapitulatif
    \def\colA{3.2cm}
    \def\colB{4.5cm}
    \def\colC{5.5cm}
    \def\rowH{1.1cm}
    % En-tête
    \fill[popblue] (0,0) rectangle (\colA+\colB+\colC, \rowH);
    \node[white, font=\bfseries\large, align=center] at (\colA/2, \rowH/2) {PRINCIPE};
    \node[white, font=\bfseries\large, align=center] at (\colA+\colB/2, \rowH/2) {FORMULE \& SIGNIFICATION};
    \node[white, font=\bfseries\large, align=center] at (\colA+\colB+\colC/2, \rowH/2) {EXEMPLE CAMEROUNAIS};
    % Ligne 1 : Identité
    \fill[popblue!10] (0,-\rowH) rectangle (\colA+\colB+\colC, -2*\rowH);
    \draw[popline] (0,-\rowH) -- (\colA+\colB+\colC,-\rowH);
    \draw[popline] (\colA,0) -- (\colA,-3*\rowH);
    \draw[popline] (\colA+\colB,0) -- (\colA+\colB,-3*\rowH);
    \node[font=\bfseries, align=center] at (\colA/2, -1.5*\rowH) {Identité};
    \node[align=left, text width=\colB-0.4cm] at (\colA+\colB/2, -1.5*\rowH) {$A = A$\\Tout être est ce qu'il est.\\Stabilité du sens.};
    \node[align=left, text width=\colC-0.4cm] at (\colA+\colB+\colC/2, -1.5*\rowH) {Dans un contrat de bail à Yaoundé, \og locataire \fg\ désigne la même personne de l'art. 1 à l'art. 20.};
    % Ligne 2 : Non-contradiction
    \fill[poporange!10] (0,-2*\rowH) rectangle (\colA+\colB+\colC, -3*\rowH);
    \draw[popline] (0,-2*\rowH) -- (\colA+\colB+\colC,-2*\rowH);
    \node[font=\bfseries, align=center] at (\colA/2, -2.5*\rowH) {Non-\\contradiction};
    \node[align=left, text width=\colB-0.4cm] at (\colA+\colB/2, -2.5*\rowH) {$\neg(A \land \neg A)$\\Impossible d'être et de ne pas être\\en même temps et sous même rapport.};
    \node[align=left, text width=\colC-0.4cm] at (\colA+\colB+\colC/2, -2.5*\rowH) {Un conducteur ne peut pas dire : \og J'ai passé le feu rouge \fg\ et \og Je n'ai pas passé le feu rouge \fg\ pour le même feu, même instant.};
    % Ligne 3 : Tiers exclu
    \fill[popgreen!10] (0,-3*\rowH) rectangle (\colA+\colB+\colC, -4*\rowH);
    \draw[popline] (0,-3*\rowH) -- (\colA+\colB+\colC,-3*\rowH);
    \node[font=\bfseries, align=center] at (\colA/2, -3.5*\rowH) {Tiers exclu};
    \node[align=left, text width=\colB-0.4cm] at (\colA+\colB/2, -3.5*\rowH) {$A \lor \neg A$\\Toute proposition est VRAIE ou FAUSSE.\\Pas de troisième valeur de vérité.};
    \node[align=left, text width=\colC-0.4cm] at (\colA+\colB+\colC/2, -3.5*\rowH) {\og Le Cameroun a remporté la CAN 2017 \fg\ est VRAI. \og Le Cameroun a remporté la CAN 2023 \fg\ est FAUX. Pas \og à moitié vrai \fg.};
    % Bordure extérieure
    \draw[popdark, line width=1.2pt] (0,0) rectangle (\colA+\colB+\colC, -4*\rowH);
\end{tikzpicture}
\legende{Les trois principes logiques classiques : formule, signification et illustration par des situations concrètes camerounaises.}
\end{popfigure}

\courssub{Applications concrètes : penser juste au quotidien}

\begin{methode}[Détecter une incohérence dans un discours (politique, médiatique, administratif)]
\begin{enumerate}
    \item \textbf{Isoler les propositions} : repérer les affirmations nettes (sujet + prédicat + temps + lieu).
    \item \textbf{Fixer le sujet, l'attribut, le moment et le rapport} : \og Le ministre a dit que le taux de chômage est de 4\% en 2023 \fg\ vs \og Le ministre a dit qu'il est de 12\% en 2023 \fg.
    \item \textbf{Tester l'identité} : les termes \og taux de chômage \fg, \og 2023 \fg, \og Cameroun \fg\ ont-ils le même sens ?
    \item \textbf{Tester la non-contradiction} : les deux valeurs (4\% et 12\%) peuvent-elles être vraies ensemble ? Non $\rightarrow$ contradiction.
    \item \textbf{Tester le tiers exclu} : la proposition \og Le taux est 4\% \fg\ est vraie ou fausse. Si fausse, la vraie valeur est une autre (non-4\%).
    \item \textbf{Exiger la raison suffisante} : demander la source (INS, BIT), la méthodologie (enquête, définition du chômage).
\end{enumerate}
\end{methode}

\begin{exoresolu}[Analyse d'un extrait de discours politique]
\textbf{Énoncé :} Un candidat déclare : \og Je n'ai jamais détourné d'argent public. Mes adversaires mentent. Mais si l'on prouve que j'ai utilisé des fonds de la commune pour ma campagne, c'était pour le bien du peuple. \fg
\textbf{Question :} Identifiez les violations des principes logiques.

\tcblower

\textbf{Solution détaillée :}
\begin{enumerate}
    \item \textbf{Violation de la non-contradiction :} \og Je n'ai jamais détourné \fg\ ($A$) et \og J'ai utilisé des fonds de la commune pour ma campagne \fg\ ($\neg A$, car détournement = usage privé/illicite de fonds publics). Même sujet (le candidat), mêmes fonds (communaux), même période (campagne). Les deux affirmations ne peuvent être vraies ensemble.
    \item \textbf{Violation de l'identité (équivoque) :} Le terme \og détournement \fg\ est implicitement redéfini : au début, c'est un crime ; à la fin, c'est \og pour le bien du peuple \fg\ (justification morale). Le sens glisse.
    \item \textbf{Tiers exclu :} La proposition \og Il a détourné \fg\ est vraie ou fausse. Le candidat tente d'introduire un \og tiers \fg\ : \og c'est vrai mais justifié \fg. Logiquement, l'acte est soit un détournement (vrai), soit ce n'en est pas un (faux). La justification morale ne change pas la qualification juridique.
    \item \textbf{Raison suffisante absente :} Aucune preuve, aucun document comptable n'est produit pour étayer \og c'était pour le bien du peuple \fg.
\end{enumerate}
\end{exoresolu}

\courssub{Portée et limites : fondements de la rationalité, non garantie de vérité}

\begin{propriete}[Portée des principes logiques]
\begin{itemize}
    \item \textbf{Condition de la cohérence :} Sans eux, tout discours s'effondre (on pourrait tout dire et son contraire).
    \item \textbf{Outil de validation formelle :} Ils permettent de vérifier la \textbf{validité} d'un raisonnement (la conclusion suit-elle des prémisses ?).
    \item \textbf{Base de la communication :} Le contrat, la loi, le témoignage, l'exposé scientifique supposent l'identité des termes et l'exclusion de la contradiction.
\end{itemize}
\end{propriete}

\begin{propriete}[Limites essentielles — à connaître pour le Bac]
\begin{itemize}
    \item \textbf{Ils ne garantissent pas la vérité matérielle.} Un raisonnement peut être \textbf{valide} (respecter les principes) mais \textbf{faux} (prémisses fausses).
    \item \textbf{Exemple :} \og Tous les éléphants volent. Dumbo est un éléphant. Donc Dumbo vole. \fg\ $\rightarrow$ Valide (structure correcte, non-contradiction respectée), mais faux (prémisse majeure fausse).
    \item \textbf{Ils ne s'appliquent qu'aux propositions assertoriques claires.} Les questions, ordres, vœux, phrases ambiguës, métaphores, poèmes échappent au vrai/faux bivalent.
    \item \textbf{Logiques non classiques :} En logique floue, multivalente, quantique, le tiers exclu est assoupli. Mais au programme de Première technique, on reste dans le cadre \textbf{classique (bivalent)}.
\end{itemize}
\end{propriete}

\begin{aretenir}[Synthèse pour l'épreuve]
\begin{itemize}
    \item Les \textbf{trois principes aristotéliciens} (identité, non-contradiction, tiers exclu) sont les \textbf{lois formelles} de la pensée cohérente.
    \item Le \textbf{principe de raison suffisante} (Leibniz) est l'exigence de \textbf{justification} de toute vérité.
    \item \textbf{Maîtriser ces principes} = savoir \textbf{structurer} sa pensée, \textbf{repérer} les contradictions, \textbf{exiger} des preuves, \textbf{refuser} les faux dilemmes.
    \item \textbf{Attention :} Validité logique $\neq$ Vérité réelle. Un syllogisme parfait peut partir de prémisses fausses.
\end{itemize}
\end{aretenir}

\begin{exercice}[Entraînement — Détection de violations]
\textbf{Consignes :} Pour chaque énoncé, précisez quel(s) principe(s) est/sont violé(s) et justifiez en une phrase.
\begin{enumerate}
    \item \og Le nouveau programme est excellent. Mais il est nul. \fg\ (même locuteur, même instant, même programme).
    \item \og Soit tu acceptes ce marché tel quel, soit tu perds le client. \fg\ (négociation commerciale).
    \item \og L'eau bout à 100°C. Donc l'eau bout à 100°C au sommet du Mont Cameroun. \fg
    \item \og Je ne sais pas si le dossier est complet ou incomplet. \fg\ (archiviste devant un dossier).
\end{enumerate}
\end{exercice}

\begin{corrige}[Exercice n°1]
\begin{enumerate}
    \item \textbf{Non-contradiction} (et identité) : \og excellent \fg\ et \og nul \fg\ sont contradictoires pour le même objet, même moment, même rapport (qualité globale).
    \item \textbf{Faux dilemme} (sophisme), \textbf{pas} violation du tiers exclu : le tiers exclu s'applique à une proposition (\og Ce marché est acceptable \fg\ : V/F), pas à un choix d'action où d'autres options existent (négocier, refuser, reporter).
    \item \textbf{Identité / Raison suffisante} : \og eau \fg\ même terme, mais conditions changées (pression atmosphérique). La prémisse manque de précision (pression normale). Violation de la raison suffisante (condition omise) et glissement d'identité (eau à 1 atm vs eau à 0,8 atm).
    \item \textbf{Aucune violation} : L'ignorance épistémique (\og je ne sais pas \fg{}) n'est pas une troisième valeur de vérité. La proposition \og Le dossier est complet \fg\ EST vraie ou fausse dans la réalité ; l'archiviste ignore laquelle. Le tiers exclu tient sur la proposition, pas sur la connaissance du sujet.
\end{enumerate}
\end{corrige}

\coursec{Application : penser juste dans la vie quotidienne}

Après avoir établi les \emph{principes logiques de la pensée} — identité, non-contradiction, tiers exclu, raison suffisante — qui sont les lois souveraines de tout discours rationnel, nous devons maintenant descendre de la théorie à la pratique. La logique n'est pas une gymnastique stérile réservée aux manuels ; c'est une \cle{boîte à outils} pour naviguer dans le flot incessant d'informations, d'arguments et de manipulations qui caractérise notre vie quotidienne au Cameroun comme ailleurs. Cette section finale a pour but de vous rendre \emph{autonome} : savoir déconstruire un discours publicitaire, repérer un sophisme dans un meeting politique, structurer votre propre pensée pour convaincre sans tromper, et faire de la rigueur logique un acte de citoyenneté responsable.

\courssub{Synthèse du parcours : les trois opérations de l'esprit}

La logique classique, depuis Aristote jusqu'aux Port-Royal, distingue trois \cle{opérations de l'esprit} qui s'enchaînent nécessairement. On ne peut raisonner sans juger, et on ne peut juger sans concevoir des concepts. C'est une \emph{hiérarchie de complexité} :

\begin{enumerate}
    \item \textbf{Concevoir (l'appréhension) :} L'esprit saisit l'essence d'une chose et forme un \cle{concept}. Le concept n'est ni vrai ni faux, il est \emph{clair} ou \emph{obscur}, \emph{distinct} ou \emph{confus}. Il s'exprime par un \cle{terme} (mot ou groupe de mots) : « huile moteur », « citoyen », « justice ».
    \item \textbf{Juger :} L'esprit unit ou sépare deux concepts par un acte d'affirmation ou de négation. C'est le \cle{jugement}. Il porte en lui la \cle{valeur de vérité} (vrai ou faux). Il s'exprime par une \cle{proposition} (phrase complète) : « Cette huile moteur est certifiée ISO », « Tout citoyen a des devoirs ».
    \item \textbf{Raisonner (le discours) :} L'esprit relie plusieurs jugements (prémisses) pour en déduire un nouveau (conclusion). C'est le \cle{raisonnement} (déduction, induction, analogie). Il s'exprime par une \cle{argumentation} (texte structuré) : « Toute huile certifiée ISO résiste à 150°C. Cette huile est certifiée ISO. Donc elle résiste à 150°C. »
\end{enumerate}

\begin{popfigure}
\begin{tikzpicture}[
    node distance=1.5cm and 2cm,
    op/.style={draw=popblue, thick, fill=popblueL, rounded corners, minimum width=3cm, minimum height=1.2cm, align=center, font=\bfseries},
    expr/.style={draw=poporange, thick, fill=poporangeL, rounded corners, minimum width=3.5cm, minimum height=1cm, align=center, font=\small},
    arrow/.style={-Stealth, thick, popdark},
    link/.style={draw=popgreen!50!black, thick, dashed, -Stealth}
]
% Nœuds opérations
\node[op, align=center] (c1){1. CONCEVOIR\\ (Appréhension)};
\node[op, right=of c1] (c2) {2. JUGER};
\node[op, right=of c2, align=center] (c3){3. RAISONNER\\ (Discours)};
% Nœuds expressions
\node[expr, below=of c1, align=center] (e1){TERME\\ \textit{ex: ``citoyen''}};
\node[expr, below=of c2, align=center] (e2){PROPOSITION\\ \textit{ex: ``Le citoyen vote''}};
\node[expr, below=of c3, align=center] (e3){ARGUMENTATION\\ \textit{ex: Syllogisme}};
% Flèches verticales
\draw[arrow] (c1.south) -- (e1.north);
\draw[arrow] (c2.south) -- (e2.north);
\draw[arrow] (c3.south) -- (e3.north);
% Flèches horizontales (enchaînement)
\draw[arrow] (c1.east) -- node[above, font=\small] {forme} (c2.west);
\draw[arrow] (c2.east) -- node[above, font=\small] {structure} (c3.west);
% Flèches expressions
\draw[link] (e1.east) -- (e2.west);
\draw[link] (e2.east) -- (e3.west);
% Annotations latérales
\node[left=0.5cm of c1, align=right, font=\small, text width=2.5cm] (a1) {Opération\\simple\\\textit{(ni V ni F)}};
\node[left=0.5cm of c2, align=right, font=\small, text width=2.5cm] (a2) {Acte de\\vérité\\ \textit{(V ou F)}};
\node[left=0.5cm of c3, align=right, font=\small, text width=2.5cm] (a3) {Enchaînement\\ déductif/\\inductif};
\draw[popmuted, decorate, decoration={brace, amplitude=5pt}] (a1.north west) -- (a1.south west);
\draw[popmuted, decorate, decoration={brace, amplitude=5pt}] (a2.north west) -- (a2.south west);
\draw[popmuted, decorate, decoration={brace, amplitude=5pt}] (a3.north west) -- (a3.south west);
\end{tikzpicture}
\legende{Les trois opérations de l'esprit et leurs expressions linguistiques. La flèche pleine indique la dépendance ontologique ; la flèche pointillée le lien linguistique.}
\end{popfigure}

\begin{aretenir}[Les trois opérations de l'esprit]
\begin{itemize}
    \item \textbf{Concept $\rightarrow$ Terme :} Unité de pensée, pas de valeur de vérité. Exigence : \cle{clarté} et \cle{distinction} (définition par genre et différence spécifique).
    \item \textbf{Jugement $\rightarrow$ Proposition :} Synthèse de deux concepts (Sujet - Copule - Attribut). Exigence : \cle{véracité} (adéquation au réel) et respect des \cle{principes logiques}.
    \item \textbf{Raisonnement $\rightarrow$ Argumentation :} Enchaînement de propositions (Prémisses $\rightarrow$ Conclusion). Exigences : \cle{validité} (forme correcte) et \cle{vérité des prémisses} (solidité = validité + vérité).
\end{itemize}
\end{aretenir}

\courssub{Analyse guidée de discours authentiques : la grille de lecture logique}

Face à un discours (publicité, allocution politique, post Facebook, conversation familiale), le citoyen logique ne se laisse pas emporter par l'émotion ou l'autorité de l'émetteur. Il applique une \cle{grille d'analyse} systématique. Cette grille opère à trois niveaux correspondant aux trois opérations de l'esprit.

\begin{popfigure}
\begin{center}
\begin{tabularx}{\linewidth}{|>{\bfseries}l|X|X|}
\hline
\rowcolor{popblueL}
\textbf{Niveau d'analyse} & \textbf{Questions opératoires} & \textbf{Critères de validation / Pièges à repérer} \\
\hline
\textbf{1. Niveau Conceptuel (Termes)} & Les termes clés sont-ils définis ? Sont-ils univoques ou équivoques ? La division/classement est-elle exhaustive et exclusive ? & \textbf{Clarté / Précision.} \textit{Pièges :} Amphibologie (ambiguïté), équivoque (glissement de sens), concept « valise » (fourre-tout vide). \\
\hline
\textbf{2. Niveau Judicatif (Propositions)} & Les affirmations sont-elles vérifiables ? Distingue-t-on fait / opinion / valeur ? Y a-t-il des jugements universels abusifs (``Tous...'', ``Jamais...'') ? & \textbf{Véracité / Nuance.} \textit{Pièges :} Généralisation hâtive, confusion fait/opinion, énoncés non falsifiables. \\
\hline
\textbf{3. Niveau Raisonnement (Argumentation)} & Quelle structure ? (Syllogisme, induction, analogie, cause/effet). La forme est-elle valide ? Les prémisses sont-elles vraies ? La conclusion suit-elle nécessairement ? & \textbf{Validité / Solidité.} \textit{Pièges :} Sophismes (formels ou informels), paralogismes, \textit{non sequitur}, pétition de principe, homme de paille, argument d'autorité, \textit{ad hominem}. \\
\hline
\textbf{4. Respect des Principes} & Le discours respecte-t-il l'identité (A=A) ? Évite-t-il la contradiction (A et non-A) ? Assume-t-il le tiers exclu ? Cherche-t-il la raison suffisante ? & \textbf{Cohérence / Rationalité.} \textit{Pièges :} Double langage, contradictions internes, assertions gratuites (``C'est comme ça''). \\
\hline
\end{tabularx}
\end{center}
\legende{Grille d'analyse logique d'un discours (vierge à reproduire pour vos exercices).}
\end{popfigure}

\begin{exemplebox}[Analyse d'une publicité : « 9 chauffeurs sur 10 choisissent l'Huile X »]
\textbf{Discours :} Une affiche géante au carrefour Warda (Yaoundé) : « 9 chauffeurs sur 10 choisissent l'Huile X pour leur moteur. Faites comme les pros ! »

\textbf{Analyse par la grille :}
\begin{enumerate}
    \item \textbf{Concepts :} « Chauffeur » (taxi ? poids lourd ? particulier ?), « choisir » (libre choix ou contrainte contrat flotte ?), « pro » (critère de compétence ?). \textit{Flou conceptuel.}
    \item \textbf{Jugements :} « 9 sur 10 choisissent ». Source ? Échantillon ? Méthode ? « Faites comme les pros » : jugement de valeur implicite (le choix du pro = bon choix).
    \item \textbf{Raisonnement :} Structure inductive (statistique) + Argument d'analogie / conformisme social (\textit{argumentum ad populum} déguisé).
    \begin{itemize}
        \item \textit{Validité inductive :} Faible. Sans protocole (échantillon représentatif, taille, indépendance), la statistique est un \textbf{chiffre ornemental}.
        \item \textit{Force persuasive :} Forte psychologiquement (preuve sociale), nulle logiquement.
    \end{itemize}
    \item \textbf{Principes :} Principe de raison suffisante violé : aucune raison technique (viscosité, additifs, normes API/ACEA) n'est donnée pour justifier le choix.
\end{enumerate}
\textbf{Conclusion logique :} Ce n'est pas un argument, c'est une \cle{technique de persuasion} (marketing). Le consommateur rationnel demandera la fiche technique, pas le sondage.
\end{exemplebox}

\begin{exemplebox}[Autres discours à analyser en classe]
\begin{itemize}
    \item \textbf{Discours électoral :} « Mon adversaire a géré la mairie pendant 5 ans, les routes sont dégradées. Donc il est incompétent. Votez pour moi, je construirai des routes. » \textit{Analyse :} Lien de causalité unique simplifié (post hoc ergo propter hoc) ; promesse non conditionnée (budget, études) ; homme de paille possible si le bilan est plus nuancé.
    \item \textbf{Débat familial :} « Tu ne m'as pas appelé pour la fête, donc tu ne m'aimes plus. » \textit{Analyse :} Faux dilemme (appel = amour ?) ; confusion signe/réalité ; raisonnement hypotético-déductif inversé (conséquent affirmé).
    \item \textbf{Publication Facebook :} « Partagez si vous êtes contre la vie chère ! 1 partage = 1 FCFA pour les pauvres. » \textit{Analyse :} Mécanisme viral (chaîne) ; fausse causalité (partage $\neq$ don) ; appel à l'émotion (pitié) ; absence de traçabilité (violation raison suffisante).
\end{itemize}
\end{exemplebox}

\courssub{Rédiger une démarche argumentative rigoureuse}

Inversement, vous devez être capable de \emph{produire} un discours logiquement solide, que ce soit pour une dissertation au Probatoire, une lettre de motivation, une réclamation administrative ou une prise de parole en réunion. La méthode en cinq étapes ci-dessous est votre standard.

\begin{methode}[Rédiger une démarche argumentative (standard Bac / Vie pro)]
\begin{enumerate}
    \item \textbf{Définir les termes (Opération 1 : Concevoir).} Posez les définitions précises des concepts clés du sujet. Éliminez l'équivoque. \textit{Ex: « Par "vie chère", j'entends l'augmentation générale et durable du niveau des prix... »}
    \item \textbf{Énoncer clairement la thèse (Opération 2 : Juger).} Formulez votre conclusion visée sous forme de proposition affirmative ou négative, nuancée si nécessaire. \textit{Ex: « La lutte contre la vie chère nécessite une action structurelle sur l'offre locale, pas seulement des subventions conjoncturelles. »}
    \item \textbf{Enchaîner prémisses et conclusion (Opération 3 : Raisonner).} Construisez votre argumentation comme une chaîne : Prémisse 1 $\rightarrow$ Prémisse 2 $\rightarrow$ ... $\rightarrow$ Conclusion. Utilisez des connecteurs logiques (\textit{donc, car, si... alors, or, par conséquent}). Vérifiez chaque maillon : est-ce un syllogisme valide ? Une induction forte ? Une analogie pertinente ?
    \item \textbf{Vérifier validité et vérité (Contrôle qualité).}
    \begin{itemize}
        \item \textit{Forme :} La conclusion suit-elle nécessairement des prémisses ? (Test du contre-exemple : puis-je imaginer prémisses vraies et conclusion fausse ?).
        \item \textit{Fond :} Les prémisses sont-elles vraies, vérifiées, sourcées ? Sont-elles acceptables par l'adversaire (prémisses communes) ?
    \end{itemize}
    \item \textbf{Conclure et ouvrir.} Répétez la thèse démontrée. Signalez les limites ou les conditions d'application. Proposez une perspective concrète.
\end{enumerate}
\end{methode}

\begin{attention}[Piège mortel : L'argument d'autorité mal placé]
\textbf{Erreur fréquente :} « Le Ministre / Le Chef / Mon grand-père / La star X l'a dit, donc c'est vrai. »
\begin{itemize}
    \item \textbf{Logique :} L'argument d'autorité (\textit{argumentum ad verecundiam}) n'est \textbf{jamais} une démonstration en soi. Il ne prouve que : « X \emph{affirme} que P ». Il ne prouve pas P.
    \item \textbf{Usage légitime :} Seule une \textbf{autorité compétente, indépendante, consensuelle dans son domaine, citant des faits vérifiables} peut servir de \emph{prémisse inductive forte} (ex: « L'OMS affirme que le vaccin X est sûr à 95\% » $\rightarrow$ prémisse forte pour « Je me fais vacciner »). Mais même alors, on cite la \emph{source primaire} (l'étude), pas l'institution.
    \item \textbf{Au Bac / Concours :} Une copie qui justifie une thèse par « Victor Hugo a dit que... » ou « Selon mon professeur... » sans analyse propre perd les points de « rigueur logique » et « esprit critique ».
\end{itemize}
\end{attention}

\courssub{TP : Construire un tableau de classification et un diagramme de Venn}

La logique se \emph{manipule}. Deux outils visuels indispensables pour valider vos concepts et vos syllogismes : la \cle{division dichotomique} (tableau de classification) et le \cle{diagramme de Venn} (représentation ensembliste des propositions catégoriques).

\begin{experience}[TP Logique : De la classification au syllogisme validé]
\textbf{Objectif :} Classer un concept complexe par division dichotomique et valider visuellement un syllogisme catégorique (Mode AAA-1 : Barbara).

\textbf{Matériel :} Feuille A4, crayon, règle, compas (ou logiciel de géométrie dynamique).

\textbf{Étape 1 : Division dichotomique d'un concept (Tableau de classification)}
\begin{enumerate}
    \item Choisissez un genre large : \textbf{« Véhicules terrestres à moteur »}.
    \item Appliquez la \cle{règle de division dichotomique} : une différence spécifique à la fois, exhaustive et exclusive (A / non-A).
    \item \textbf{Différence 1 :} Usage principal $\rightarrow$ \textbf{Transport de personnes} / \textbf{Transport de marchandises} (exclusif et exhaustif pour l'usage principal).
    \item \textbf{Différence 2 (sous Transport de personnes) :} Capacité $\rightarrow$ \textbf{$\le$ 9 places (VP, Minibus)} / \textbf{$>$ 9 places (Autobus, Autocar)}.
    \item \textbf{Différence 3 (sous $\le$ 9 places) :} Énergie $\rightarrow$ \textbf{Thermique} / \textbf{Non thermique (Électrique, Hybride, Hydrogène)}.
\end{enumerate}
\textbf{Produit attendu :} Un tableau arborescent (arbre de Porphyre) ou un tableau à colonnes \textit{Genre / Espèce / Différence spécifique} respectant l'exhaustivité et l'exclusion mutuelle à chaque nœud.

\textbf{Étape 2 : Diagramme de Venn pour valider un syllogisme (Mode AAA-1 : Barbara)}
\textbf{Syllogisme :}
\begin{itemize}
    \item \textit{Majeure :} Tous les \textbf{Hommes} sont \textbf{Mortels}. (A : Tout M est P)
    \item \textit{Mineure :} Tous les \textbf{Philosophes} sont des \textbf{Hommes}. (A : Tout S est M)
    \item \textit{Conclusion :} Tous les \textbf{Philosophes} sont \textbf{Mortels}. (A : Tout S est P)
\end{itemize}
\textbf{Protocole Venn (3 cercles S, M, P) :}
\begin{enumerate}
    \item Dessinez 3 cercles sécants (8 zones). Légendez : S = Philosophes, M = Hommes, P = Mortels.
    \item \textit{Majeure « Tout M est P » :} La zone M \emph{hors} P est \textbf{vidée} (hachurée). $\rightarrow$ Hachurez la partie du cercle M qui ne chevauche pas P.
    \item \textit{Mineure « Tout S est M » :} La zone S \emph{hors} M est \textbf{vidée}. $\rightarrow$ Hachurez la partie du cercle S qui ne chevauche pas M.
    \item \textit{Lecture conclusion :} Regardez la zone S \emph{hors} P. Est-elle vidée par les hachures précédentes ? \textbf{OUI.} La conclusion « Tout S est P » est \textbf{nécessairement vraie} (valide).
\end{enumerate}

\begin{popfigure}
\begin{tikzpicture}[scale=1.3]
% Définition des centres et rayon
\def\R{1.6}
\def\Sx{-1.0}\def\Sy{0.0}
\def\Px{1.0}\def\Py{0.0}
\def\Mx{0.0}\def\My{1.3}
% --- Hachures (Zones VIDES) ---
% 1. Majeure : Tout M est P  => Zone M \ P est vide (hachurée)
\begin{scope}
    \clip (\Mx,\My) circle (\R);          % Restreint au cercle M
    \fill[white] (\Px,\Py) circle (\R);    % Efface l'intersection M∩P (peint en blanc par-dessus)
    \fill[pattern=north east lines, pattern color=popdark] (\Mx,\My) circle (\R); % Hachure tout M (donc M\P reste hachuré)
\end{scope}
% 2. Mineure : Tout S est M  => Zone S \ M est vide (hachurée différemment)
\begin{scope}
    \clip (\Sx,\Sy) circle (\R);          % Restreint au cercle S
    \fill[white] (\Mx,\My) circle (\R);    % Efface l'intersection S∩M
    \fill[pattern=north west lines, pattern color=popdark] (\Sx,\Sy) circle (\R); % Hachure tout S (donc S\M reste hachuré)
\end{scope}
% --- Cercles (Contours) ---
\draw[thick, popblue] (\Sx,\Sy) circle (\R) node[left=2.0cm, align=center]{\textbf{S}\\(Philosophes)};
\draw[thick, poporange] (\Px,\Py) circle (\R) node[right=2.0cm, align=center]{\textbf{P}\\(Mortels)};
\draw[thick, popgreen] (\Mx,\My) circle (\R) node[above=2.0cm, align=center]{\textbf{M}\\(Hommes)};
% --- Marqueur de conclusion : Zone S \ P ---
% Elle est couverte par : hachure S\M (nord-ouest) U hachure M\P (nord-est).
% On place un symbole de validation au centre de la zone S\P (en bas à gauche du cercle S)
\node[font=\bfseries\small, popgreen, fill=popcream, rounded corners, inner sep=2pt] at (\Sx-0.9,\Sy-0.9) {S \textbackslash P = $\varnothing$ \checkmark};
% Flèche explicative
\draw[popcurve, ->, thick] (\Sx-2.5,\Sy-2.0) -- (\Sx-1.2,\Sy-1.1) node[right, font=\small, align=left] {Conclusion valide :\\la zone S non P\\est entièrement hachurée};
\end{tikzpicture}
\legende{Diagramme de Venn validant le syllogisme \textit{Barbara} (AAA-1). Zones hachurées = vides. La zone S$\setminus$P (bas gauche du cercle bleu) est couverte par l'union des hachures des prémisses : la conclusion « Tout S est P » est forcée.}
\end{popfigure}
\end{experience}

\begin{exoresolu}[Exercice d'application : Analyser et corriger un raisonnement fallacieux]
\textbf{Énoncé :} Un élève de Première T déclare : « Puisque le prof de philo a dit que "la logique est l'art de ne pas se tromper", et que je me suis trompé sur l'exercice 3, c'est que je n'ai pas utilisé la logique. Donc, pour ne plus me tromper, il suffit que j'apprenne la logique par cœur. »

\textbf{Analyse logique :}
\begin{enumerate}
    \item \textbf{Concepts :} « Logique » (définition restrictive : art de ne pas se tromper ?), « Se tromper » (erreur de calcul ? de méthode ? d'inattention ?), « Apprendre par cœur » (mémorisation vs compréhension).
    \item \textbf{Structure apparente (Syllogisme) :}
    \begin{itemize}
        \item Majeure : User de logique $\rightarrow$ Ne pas se tromper. (Si L alors non-E)
        \item Mineure : Je me suis trompé. (E)
        \item Conclusion intermédiaire : Je n'ai pas usé de logique. (non L) $\leftarrow$ \textbf{Valide (Modus Tollens)}.
    \end{itemize}
    \item \textbf{Glissement sophistique (Enthymème caché) :}
    \begin{itemize}
        \item Prémisse cachée : « Apprendre la logique par cœur » = « User de la logique ».
        \item Conclusion finale : Apprendre par cœur $\rightarrow$ Ne plus se tromper.
    \end{itemize}
    \item \textbf{Critique :}
    \begin{itemize}
        \item \textit{Équivoque sur « Logique » :} La logique est une \emph{méthode} (opératoire), pas une \emph{somme de savoirs} (déclaratif). L'apprendre par cœur ne garantit pas son \emph{application} correcte (transfert, vigilance, esprit critique).
        \item \textit{Cause unique / Solution magique :} Se tromper peut venir de la fatigue, de la lecture trop rapide, du manque de connaissances du domaine (maths, physique), pas seulement de l'absence de logique.
        \item \textit{Non sequitur final :} La conclusion « il suffit que... » est un saut quantitatif non justifié (condition nécessaire $\neq$ condition suffisante).
    \end{itemize}
\end{enumerate}
\textbf{Correction reformulée :} « La logique est un outil de rigueur. Mon erreur à l'exercice 3 signale un défaut de rigueur \emph{à ce moment-là}. Pour progresser, je dois \emph{m'exercer} aux méthodes logiques (identifier prémisses, vérifier validité, définir termes) sur des cas variés, et non seulement mémoriser les définitions. »
\end{exoresolu}

\courssub{La logique, arme citoyenne contre la rumeur et la désinformation}

Au-delà de la réussite scolaire (Probatoire, concours ENAM/IRIC), la logique est une \cle{vertu civique}. Dans une société saturée d'infox (\textit{fake news}), de rumeurs WhatsApp, de vidéos sorties de contexte et de discours de haine ethnique ou religieux, le citoyen qui ne sait pas \emph{distinguer un argument d'une incantation}, un \emph{fait d'une opinion}, une \emph{corrélation d'une causalité}, est une proie facile pour les manipulateurs de tout bord.

\begin{savaistu}[Logique et Concours Administratifs Camerounais (ENAM, IRIC, ENSET, Police...)]
Les épreuves de \textbf{Culture Générale}, de \textbf{Français} (résumé, essai), et surtout les \textbf{Tests Psychotechniques / Aptitude} des grands concours (ENAM cycle A/B, IRIC, Police, Douanes, ENSET) comportent systématiquement des items de \textbf{raisonnement logique} :
\begin{itemize}
    \item \textbf{Syllogismes :} « Tous les A sont B. Certains B sont C. Peut-on conclure "Certains A sont C" ? » (Réponse : Non, figure invalide).
    \item \textbf{Suites logiques / Analogies :} « Avocat : Plaider :: Juge : ? » (Juger / Rendre justice).
    \item \textbf{Analyse de textes :} Repérer l'argument fallacieux, l'intrus, la thèse principale.
    \item \textbf{Épreuve de dissertation :} La note de « Plan / Structure / Logique du développement » pèse souvent 30 à 40\% de la note. Un plan illogique (sauts, répétitions, contradictions) = échec assuré.
\end{itemize}
\textbf{Conseil :} Entraînez-vous quotidiennement avec des exercices de logique formelle (sites comme \textit{Logique.fr}, \textit{Khan Academy}, annales concours). C'est un « muscle » qui se développe.
\end{savaistu}

\begin{definition}[Sophisme vs Paralogisme : L'intention fait la différence]
\begin{itemize}
    \item \textbf{Le Sophisme :} Raisonnement \textbf{invalide ou faux} présenté \textbf{avec l'intention de tromper}, de persuader indûment, de manipuler. C'est une \emph{faute morale} intellectuelle. \textit{Ex:} Un politique qui sait que ses chiffres sont faux mais les utilise car « ça marche auprès de la base ».
    \item \textbf{Le Paralogisme :} Raisonnement \textbf{invalide ou faux} commis \textbf{involontairement}, par ignorance, précipitation, confusion, biais cognitif. C'est une \emph{erreur technique}. \textit{Ex:} Un élève qui affirme « Tous les serpents sont dangereux, cet animal est dangereux, donc c'est un serpent » (affirmation du conséquent) sans voir l'erreur.
\end{itemize}
\textbf{Devoir du philosophe / citoyen :} Dénoncer le sophisme (combat éthique), corriger le paralogisme (acte pédagogique), et surtout \emph{ne pas en commettre soi-même}.
\end{definition}

\courssub{Exercices d'entraînement}

\begin{exercice}[Analyse de discours « Vie chère »]
On lit sur un groupe WhatsApp « Parents d'élèves Quartier X » :
« \textit{Le prix du riz a augmenté de 500 FCFA le sac en un mois. C'est la faute du nouveau Directeur de la SONARA qui a augmenté le carburant la semaine dernière. Le gouvernement ne fait rien. Si on ne manifeste pas samedi, on mourra de faim. Partagez massivement ! } »

\textbf{Consignes :}
\begin{enumerate}
    \item Identifiez les \textbf{concepts} flous ou non définis.
    \item Repérez les \textbf{jugements} : faits vérifiables / opinions / peurs / menaces.
    \item Reconstituez le \textbf{raisonnement} implicite (prémisses $\rightarrow$ conclusion). Nommez le type de raisonnement.
    \item Identifiez \textbf{au moins 3 sophismes / paralogismes} précis (nommez-les : \textit{post hoc, homme de paille, faux dilemme, appel à la peur, généralisation...}).
    \item Proposez une \textbf{reformulation logique et citoyenne} de ce message (méthode en 5 étapes).
\end{enumerate}
\end{exercice}


\begin{corrige}[Exercice n°1 : Analyse de discours « Vie chère »]
\begin{enumerate}
    \item \textbf{Concepts :} « Faute » (responsabilité causale unique ?), « Gouvernement » (entité monolithique ?), « Manifester » (seule issue ?), « Mourir de faim » (littéral ou hyperbolique ?).
    \item \textbf{Jugements :}
    \begin{itemize}
        \item \textit{Fait (à vérifier) :} « Prix riz +500 FCFA », « Directeur SONARA nommé », « Carburant augmenté ».
        \item \textit{Opinion/Causalité abusive :} « C'est LA faute DU Directeur » (monocausalité).
        \item \textit{Prédiction catastrophiste non fondée :} « On mourra de faim ».
        \item \textit{Injonction/Menace :} « Partagez massivement ».
    \end{itemize}
    \item \textbf{Raisonnement (Enthymème reconstruit) :}
    \begin{itemize}
        \item P1 : Le carburant a augmenté (semaine dernière).
        \item P2 : Le riz a augmenté (mois dernier). \textit{[Inversion chronologique masquée]}
        \item P3 (cachée) : L'augmentation du carburant cause \emph{immédiatement et uniquement} l'augmentation du riz.
        \item P4 (cachée) : Le gouvernement (via SONARA) est le seul responsable et peut tout résoudre instantanément.
        \item P5 (cachée) : Manifester samedi est la \emph{seule} action efficace pour baisser les prix.
        \item C : Il faut manifester samedi (sinon mort).
    \end{itemize}
    Type : \textbf{Argumentation pratique par la peur} (\textit{argumentum ad metum}) + \textbf{Causalité unique} + \textbf{Faux dilemme}.
    \item \textbf{Sophismes / Paralogismes identifiés :}
    \begin{enumerate}
        \item \textbf{Post hoc ergo propter hoc / Inversion chronologique :} Le riz a augmenté \emph{avant} le carburant (un mois vs une semaine), mais le texte lie l'effet à la cause récente.
        \item \textbf{Cause unique (Réductionnisme) :} Prix riz = f(carburant uniquement). Oubli : cours mondial, taux de change, spéculations, coûts transport, douanes, marges distributeurs.
        \item \textbf{Faux dilemme (Dilemme fallacieux) :} « Manifester OU mourir de faim ». Exclusion des autres leviers : dialogue, contrôle prix, production locale, aides ciblées, grève consommation.
        \item \textbf{Appel à la peur / À l'émotion (Ad metum / Ad passiones) :} « Mourir de faim », « Partagez massivement » (viralité émotionnelle).
        \item \textbf{Généralisation abusive / Homme de paille :} « Le gouvernement ne fait RIEN » (absolu, invérifiable, caricature).
    \end{enumerate}
    \item \textbf{Reformulation logique (Méthode 5 étapes) :}
    \begin{enumerate}
        \item \textit{Définitions :} « Vie chère » = érosion pouvoir d'achat. « Responsabilité » = part causale partagée entre acteurs étatiques, marché global, opérateurs locaux.
        \item \textit{Thèse :} « L'augmentation du riz est multifactorielle ; la hausse carburant l'aggrave. Une réponse efficace combine mesures d'urgence (subventions ciblées, contrôle) et structurelles (production locale, stockage). La manifestation est un droit, mais non une baguette magique. »
        \item \textit{Argumentation :}
        \begin{itemize}
            \item Prémisse 1 (Factuelle) : Données statistiques MINCOMMERCE / INS sur évolution prix riz/carburant 12 mois.
            \item Prémisse 2 (Économique) : Structure des coûts riz importé (Fret 30\%, Douane 15\%, Carburant local 10\%, Marge 20\%... sources : études GICAM/BCEAO).
            \item Prémisse 3 (Politique) : Mesures gouvernementales actuelles (subvention carburant, suspension droits douane riz ?) + limites (budget, ciblage).
            \item Prémisse 4 (Civique) : Manifestation = signal d'alerte démocratique. Efficacité conditionnée à organisation, revendications chiffrées, non-violence.
        \end{itemize}
        \item \textit{Vérification :} Pas de contradiction. Prémisses sourcées. Conclusion nuancée (pas de « suffit que »).
        \item \textit{Conclusion/Action :} « Exigeons la publication des structures de prix, le renforcement des stocks de sécurité SONAC, le soutien aux riziculteurs Nord (SEMRY/UNVDA), et organisons une marche déclarée avec cahier de charges précis. »
    \end{enumerate}
\end{enumerate}
\end{corrige}

\begin{exercice}[Construction : Tableau de classification + Venn]
\begin{enumerate}
    \item Construisez un \textbf{tableau de classification dichotomique} complet (jusqu'aux espèces infimes) pour le concept : \textbf{« Actes juridiques unilatéraux »}. (Indices : Volonté / Non volonté ; Effets de droit / Effets de fait ; Ex: Testament, Reconnaissance d'enfant, Perte de chose, Trouvaille).
    \item Validez par \textbf{diagramme de Venn} le syllogisme suivant :
    \begin{itemize}
        \item Aucun \textbf{Chat} n'est \textbf{Chien}. (E)
        \item Certains \textbf{Mammifères} sont \textbf{Chiens}. (I)
        \item Donc, Certains \textbf{Mammifères} ne sont pas \textbf{Chats}. (O)
    \end{itemize}
    \item Ce syllogisme est-il \textbf{valide} ? Est-il \textbf{solide} (prémisses vraies dans la réalité) ? Précisez le terme moyen réel.
\end{enumerate}
\end{exercice}


\begin{corrige}[Exercice n°2 : Classification + Venn]
\begin{enumerate}
    \item \textbf{Tableau : Actes juridiques unilatéraux (Division dichotomique)}
    \begin{center}
    \begin{tabular}{|l|l|l|}
    \hline
    \textbf{Genre} & \textbf{Différence Spécifique} & \textbf{Espèces} \\
    \hline
    Actes juridiques unilatéraux & \textit{Selon la source de l'obligation} & \\
    \hline
    & 1. \textbf{Volontaires} (Émanant d'une volonté) & $\rightarrow$ Actes juridiques stricto sensu : Testament, Promesse de récompense, Reconnaissance de dette, Mandat, Procuration. \\
    \cline{2-3}
    & 2. \textbf{Non volontaires} (Faits juridiques) & $\rightarrow$ Quasi-contrats : Gestion d'affaires, Paiement de l'indu, Enrichissement sans cause. \\
    & & $\rightarrow$ Délits / Quasi-délits : Responsabilité civile (Art 1382-1386 Code Civil). \\
    & & $\rightarrow$ Faits naturels / Actes matériels : Trouvaille, Perte, Prescription acquisitive, Accession. \\
    \hline
    \end{tabular}
    \end{center}
    \begin{attention}[Nuance juridique importante]
    En droit strict, les « actes juridiques » sont \emph{volontaires} par définition. Les items 2 (Quasi-contrats, Délits, Faits naturels) sont des \textbf{faits juridiques}, non des actes. La division correcte du genre \textbf{« Sources d'obligations »} ou \textbf{« Faits générateurs de droit »} serait : Actes juridiques (volontaires) / Faits juridiques (non volontaires). L'exercice ci-dessus classe sous le terme large « Actes juridiques unilatéraux » pour l'exercice logique, mais la distinction conceptuelle est cruciale.
    \end{attention}

    \item \textbf{Diagramme de Venn (Analyse du syllogisme proposé)}
    \textbf{Termes réels :} S = Mammifères, M = Chats, P = Chiens.
    \textbf{Prémisses :}
    \begin{itemize}
        \item Majeure (E) : Aucun M n'est P $\rightarrow$ Zone M $\cap$ P = \textbf{VIDE} (Hachurer l'intersection Chats/Chiens).
        \item Mineure (I) : Certains S sont P $\rightarrow$ Zone S $\cap$ P $\neq$ \textbf{VIDE} (Placer une \textbf{croix $\times$} dans l'intersection Mammifères/Chiens, \emph{nécessairement en dehors} de la zone Chats car M$\cap$P est vide).
    \end{itemize}
    \textbf{Conclusion (O) :} Certains S ne sont pas M $\rightarrow$ Zone S $\cap$ $\overline{\text{M}}$ $\neq$ VIDE ?
    \textbf{Vérification :} La croix $\times$ (Mammifères qui sont Chiens) est \emph{nécessairement} en dehors du cercle Chats (car Chats $\cap$ Chiens = vide). Donc cette croix se trouve bien dans la zone « Mammifères non Chats ». \textbf{La zone n'est pas vide. Conclusion VALIDE.}

    \begin{popfigure}
    \begin{tikzpicture}[scale=1.2]
    % Centres: S(-1,0), P(1,0), M(0,1.3) Rayon 1.5
    \def\R{1.5}
    \def\Sx{-1.0}\def\Sy{0.0}
    \def\Px{1.0}\def\Py{0.0}
    \def\Mx{0.0}\def\My{1.3}
    % Cercles (fond léger pour lisibilité)
    \fill[popblue!20] (\Sx,\Sy) circle (\R);
    \fill[poporange!20] (\Px,\Py) circle (\R);
    \fill[popgreen!20] (\Mx,\My) circle (\R);
    % Hachure Majeure : Aucun M n'est P => M ∩ P = VIDE
    \begin{scope}
        \clip (\Mx,\My) circle (\R);
        \clip (\Px,\Py) circle (\R); % Intersection M∩P
        \fill[pattern=north east lines, pattern color=popdark] (\Mx,\My) circle (\R); % Hachure juste l'intersection
    \end{scope}
    % Croix Mineure : Certains S sont P => S ∩ P ≠ VIDE (et hors M)
    % Coordonnées approx du centre de S∩P hors M : (0.2, -0.2)
    \node[font=\Large\bfseries, popdark] at (0.2, -0.2) {$\times$};
    \node[font=\tiny, popdark, anchor=north] at (0.2, -0.6) {S $\cap$ P (Chiens)};
    % Contours
    \draw[thick, popblue] (\Sx,\Sy) circle (\R) node[left=1.8cm, align=center]{\textbf{S}\\(Mammifères)};
    \draw[thick, poporange] (\Px,\Py) circle (\R) node[right=1.8cm, align=center]{\textbf{P}\\(Chiens)};
    \draw[thick, popgreen] (\Mx,\My) circle (\R) node[above=1.8cm, align=center]{\textbf{M}\\(Chats)};
    % Légende zones
    \node[font=\tiny, popdark, fill=popcream, rounded corners, inner sep=1pt] at (-1.8, 1.2) {S $\cap$ $\overline{\text{M}}$ $\cap$ $\overline{\text{P}}$};
    \node[font=\tiny, popdark, fill=popcream, rounded corners, inner sep=1pt] at (1.8, 1.2) {P $\cap$ $\overline{\text{S}}$ $\cap$ $\overline{\text{M}}$};
    \node[font=\bfseries\tiny, popgreen] at (0, 0.1) {M $\cap$ P = $\varnothing$};
    \end{tikzpicture}
    \legende{Diagramme de Venn pour le syllogisme : Aucun Chat n'est Chien / Certains Mammifères sont Chiens $\rightarrow$ Certains Mammifères ne sont pas Chats. La croix (Chiens) est hors Chats, donc dans Mammifères non Chats. Valide.}
    \end{popfigure}

    \item \textbf{Validité :} \textbf{OUI}. La forme est valide (raisonnement par disjonction des cas).
    \textbf{Attention :} Ce n'est \textbf{PAS} un Festino (EIO-2). Dans Festino, le terme moyen (M) est le sujet de la majeure et le prédicat de la mineure. Ici, le terme commun aux deux prémisses est \textbf{P (Chien)}, qui est le \textbf{vrai terme moyen}. La structure est : Aucun M n'est P / Certains S sont P $\rightarrow$ Certains S ne sont pas M. C'est un syllogisme valide non standard (Figure 3 avec conversion de la majeure, ou Figure 4 selon l'ordre).
    \textbf{Solidité :} \textbf{OUI}.
    \begin{itemize}
        \item « Aucun Chat n'est Chien » : VRAI (Biologie, espèces distinctes).
        \item « Certains Mammifères sont Chiens » : VRAI (Le chien est un mammifère).
        \item Conclusion « Certains Mammifères ne sont pas Chats » : VRAI (Le chien en est un exemple).
    \end{itemize}
    $\rightarrow$ Syllogisme \textbf{valide ET solide} (démonstration parfaite).
\end{enumerate}
\end{corrige}

\newpage

\coursec{Activité d'intégration}

\coursec{La logique classique}

\courssub{Objectifs de l'activité}
\begin{itemize}
    \item Mobiliser les notions de \cle{concept}, \cle{jugement}, \cle{raisonnement} et les \cle{principes logiques} (identité, non-contradiction, tiers exclu, raison suffisante) pour analyser des discours authentiques.
    \item Identifier les erreurs logiques fréquentes : généralisation hâtive, concepts flous, syllogismes invalides, contradictions.
    \item Construire une \cle{grille d'analyse logique} réutilisable pour évaluer la rigueur de tout discours médiatique, publicitaire ou conversationnel.
    \item Rédiger un \emph{verdict logique} argumenté, structuré et justifié, conforme aux exigences de la dissertation et du commentaire au Probatoire.
\end{itemize}

\courssub{Situation}
\textbf{Contexte :} Dans le cadre de la semaine de l'éducation aux médias à l'École Technique de Douala-Bassa, la classe de Première A4 organise un « Tribunal de la Logique ». Vous êtes jurés logiciens. Votre mission : examiner quatre pièces à conviction issues de la sphère publique camerounaise récente, rendre un verdict motivé pour chacune, et élaborer collectivement une grille d'évaluation universelle.

\vspace{0.5em}
\noindent\textbf{Pièce 1 — Extrait de discours de campagne (Rallye à Bamenda, mars 2025).}
\begin{quote}
\emph{« Mes chers frères et sœurs du Nord-Ouest, depuis dix ans, \textbf{chaque fois} que le parti au pouvoir organise un meeting, il pleut des cordes sur la ville. La semaine dernière encore, le meeting du 12 mars a été noyé sous un déluge. C'est la \textbf{preuve irréfutable} que le ciel lui-même rejette leur \textbf{gestion}. Par conséquent, \textbf{tout candidat} issu de ce parti apportera la \textbf{malédiction} sur notre région. Ne votez pas pour eux, \textbf{sauf si vous voulez que la famine s'installe durablement} chez nous. »}
\end{quote}

\vspace{0.5em}
\noindent\textbf{Pièce 2 — Spot publicitaire radio (Radio Balafon, Douala, prime time, avril 2025).}
\begin{quote}
\emph{[Son d'ambiance : bruit de moteur qui tousse, puis ronronnement doux] \\
Voix off : « Monsieur Talla, chauffeur de taxi au carrefour Warda, perdait 15\,000 FCFA par jour en carburant avec son vieux moteur. Depuis qu'il a mis l'additif \textbf{MoteurPro+}, il \textbf{économise 40\%} et sa voiture « grimpe » le Mont Fébé sans tousser. \textbf{MoteurPro+}, c'est la \textbf{garantie d'un moteur neuf} pour 5\,000 FCFA le flacon. \textbf{Approuvé par 9 chauffeurs sur 10} interrogés à la station Total de Bépanda. N'attendez pas la panne, adoptez MoteurPro+ ! »}
\end{quote}

\vspace{0.5em}
\noindent\textbf{Pièce 3 — Conversation au marché Mokolo (Yaoundé, samedi matin, mai 2025).}
\begin{quote}
\emph{— « Patron, ce téléphone \textit{Tecno Phantom X} est trop cher, 350\,000 F ! » \\
— « Mon frère, \textbf{tous les téléphones chers sont solides}. Ce modèle coûte 350\,000 F, \textbf{donc c'est un téléphone cher}. \textbf{Donc il est solide}. Prends-le, tu ne regretteras pas. » \\
— « Et la garantie ? » \\
— « \textbf{Tous les téléphones solides ont une bonne garantie}. Celui-ci est solide, \textbf{donc il a une bonne garantie}. C'est la logique, mon ami ! »}
\end{quote}

\vspace{0.5em}
\noindent\textbf{Pièce 4 — Publication virale WhatsApp (Groupe « Parents d'élèves Lycée de Nkolbisson », juin 2025).}
\begin{quote}
\emph{« \textbf{URGENT : PARTAGEZ MASSIVEMENT} \\
Le Ministère de l'Éducation Secondaire vient d'interdire \textbf{tous les téléphones portables} dans \textbf{tous} les lycées du Cameroun à partir de la rentrée de septembre. \textbf{C'est officiel}, le communiqué a été signé hier par le MINESEC lui-même. \\
\textbf{MAIS ATTENTION} : Une source bien introduite au ministère nous confirme qu'en réalité, \textbf{aucune interdiction n'est prévue} et que les élèves pourront garder leurs téléphones pour « faciliter la recherche pédagogique ». \\
Le gouvernement nous ment encore une fois. \textbf{Soit on interdit tout, soit on autorise tout}. Il n'y a pas de milieu. Réveillez-vous parents ! \#StopMensongesMINESEC \#Rentrée2025 »}
\end{quote}

\courssub{Consignes}
Travail en groupes de 4 à 5 élèves. Durée : 2 séances (2h). Chaque groupe rend un rapport unique.

\begin{enumerate}
    \item \textbf{Analyse conceptuelle (Pièces 1 \& 2).} Relevez 3 concepts clés mal définis ou ambigus par pièce (ex. : « gestion », « malédiction », « économie », « approuvé »). Pour chacun, proposez une \emph{définition par genre prochain et différence spécifique} rigoureuse.
    \item \textbf{Analyse des jugements (Pièces 1 \& 2).} Identifiez les jugements principaux (quantité : universel/particulier ; qualité : affirmatif/négatif) et notez leur forme canonique (A, E, I, O). Repérez les passages abusifs du particulier au universel (généralisation hâtive).
    \item \textbf{Analyse du raisonnement (Pièce 3).} 
    \begin{enumerate}
        \item Reconstituez les deux syllogismes explicites (major, mineure, conclusion).
        \item Schématisez-les chacun par un \emph{diagramme de Venn} (trois cercles).
        \item Déterminez leur \cle{validité} formelle (la conclusion suit-elle nécessairement des prémisses ?) et leur \cle{validité matérielle} (les prémisses sont-elles vraies ?).
    \end{enumerate}
    \item \textbf{Analyse des principes (Pièce 4).} Montrez précisément comment le texte viole le \cle{principe de non-contradiction}. Identifiez aussi la violation du \cle{principe du tiers exclu} (faux dilemme) et du \cle{principe de raison suffisante} (affirmation gratuite « source bien introduite »).
    \item \textbf{Rédaction du verdict.} Chaque groupe rédige un « Verdict Logique » (environ 250 mots) pour l'ensemble du dossier, structuré ainsi : \emph{Qualification des infractions logiques -- Gravité pour le débat public -- Recommandations}.
    \item \textbf{Construction de la grille.} En mise en commun (plénière), la classe élabore une \emph{Grille d'analyse logique des discours} (5 critères, 3 niveaux : Conforme / Défectueux / Fallacieux) réutilisable pour tout document médiatique.
\end{enumerate}

\courssub{Ressources à mobiliser}
\begin{definition}[Concept]
Un concept est la représentation intellectuelle des caractères communs et essentiels d'un ensemble d'objets. \textbf{Définition par genre prochain et différence spécifique} : \emph{Genre} (classe plus large) + \emph{Différence spécifique} (caractère distinctif). Ex. : \emph{Triangle} = \textit{Genre} : Polygone ; \textit{Diff. sp.} : à trois côtés.
\end{definition}

\begin{definition}[Jugement]
Acte par lequel l'esprit affirme ou nie quelque chose de quelque chose. \textbf{Formes canoniques (Carré d'Aristote)} :
\begin{center}
\begin{tabular}{|c|c|c|c|}\hline
Symbole & Quantité & Qualité & Énoncé type \\ \hline
\textbf{A} & Universel & Affirmatif & Tout S est P \\ \hline
\textbf{E} & Universel & Négatif & Nul S n'est P \\ \hline
\textbf{I} & Particulier & Affirmatif & Certains S sont P \\ \hline
\textbf{O} & Particulier & Négatif & Certains S ne sont pas P \\ \hline
\end{tabular}
\end{center}
\end{definition}

\begin{definition}[Raisonnement / Syllogisme]
Enchaînement de jugements où une conclusion découle nécessairement de prémisses. \textbf{Syllogisme catégorique} : 3 termes (Majeur P, Mineur S, Moyen M), 3 propositions (Majeure, Mineure, Conclusion). \textbf{Règles de validité} : 1) 3 termes seulement. 2) Le moyen doit être universel au moins une fois. 3) Si un terme est universel dans la conclusion, il doit l'être dans la prémisse. 4) Deux prémisses négatives $\Rightarrow$ pas de conclusion. 5) Prémisse négative $\Rightarrow$ conclusion négative.
\end{definition}

\begin{propriete}[Principes logiques fondamentaux]
\begin{itemize}
    \item \textbf{Identité :} $A$ est $A$. Un concept garde le même sens dans le raisonnement.
    \item \textbf{Non-contradiction :} Impossible que $A$ soit $B$ et non-$B$ en même temps et sous le même rapport. $\neg(P \land \neg P)$.
    \item \textbf{Tiers exclu :} $A$ est $B$ ou non-$B$, pas de troisième possibilité. $P \lor \neg P$.
    \item \textbf{Raison suffisante :} Rien n'est sans raison. Toute affirmation vraie doit pouvoir être justifiée.
\end{itemize}
\end{propriete}

\begin{methode}[Diagramme de Venn (3 ensembles S, M, P)]
\begin{enumerate}
    \item Dessiner trois cercles chevauchants (S, M, P). Les 8 zones représentent les combinaisons d'appartenance.
    \item \textbf{Hachurer} les zones \textbf{vides} (jugements universels : A $\rightarrow$ hachurer S hors P ; E $\rightarrow$ hachurer S dans P).
    \item \textbf{Placer une croix ($\times$)} dans une zone \textbf{non vide} (jugements particuliers : I $\rightarrow$ $\times$ dans S$\cap$P ; O $\rightarrow$ $\times$ dans S hors P). Si zone multiple possible, $\times$ sur le trait.
    \item \textbf{Lire la conclusion :} Si la conclusion est \emph{représentée} (zone hachurée ou croix présente) $\rightarrow$ \textbf{Valide}. Sinon $\rightarrow$ \textbf{Invalide}.
\end{enumerate}
\end{methode}

\courssub{Démarche et corrigé commenté}
\begin{exoresolu}{Corrigé guidé du Tribunal de la Logique}
\textbf{1. Analyse conceptuelle (Pièces 1 \& 2)}

\vspace{0.3em}
\noindent\textbf{Pièce 1 (Discours) :}
\begin{itemize}
    \item \cle{Gestion} : \textit{Genre} : Action d'administrer ; \textit{Diff. sp.} : mise en œuvre de politiques publiques par une équipe au pouvoir. (Le discours glisse vers « personne du candidat »).
    \item \cle{Malédiction} : \textit{Genre} : Calamité ; \textit{Diff. sp.} : fléau d'origine surnaturelle perçu comme punition divine. (Usage métaphorique non défini, glissement sémantique).
    \item \cle{Famine} : \textit{Genre} : Pénurie alimentaire ; \textit{Diff. sp.} : absence durable de nourriture entraînant mortalité massive. (Prédiction apocalyptique non étayée).
\end{itemize}

\noindent\textbf{Pièce 2 (Publicité) :}
\begin{itemize}
    \item \cle{Économie (40\%)} : \textit{Genre} : Réduction de dépense ; \textit{Diff. sp.} : diminution proportionnelle du coût carburant par rapport à une référence. (Base de calcul absente : 40\% de quoi ? 15\,000 F ?).
    \item \cle{Approuvé (9/10)} : \textit{Genre} : Validation ; \textit{Diff. sp.} : avis favorable exprimé par un échantillon. (Échantillon non représentatif : 10 chauffeurs d'une seule station).
    \item \cle{Moteur neuf} : \textit{Genre} : État mécanique ; \textit{Diff. sp.} : performance équivalente à un moteur d'usine. (Métaphore publicitaire, fausse équivalence : additif $\neq$ pièce neuve).
\end{itemize}

\vspace{0.5em}
\noindent\textbf{2. Analyse des jugements (Formes canoniques)}

\vspace{0.3em}
\noindent\textbf{Pièce 1 :}
\begin{itemize}
    \item « \emph{Chaque fois que le parti au pouvoir organise un meeting, il pleut} » $\rightarrow$ \textbf{A} (Tout meeting du parti est un événement pluvieux). \textit{Particulier (10 ans) $\rightarrow$ Universel (toujours)} = \textbf{Généralisation hâtive}.
    \item « \emph{Tout candidat issu de ce parti apportera la malédiction} » $\rightarrow$ \textbf{A} (Tout candidat du parti est porteur de malédiction). \textit{Saut du météorologique au moral/surnaturel}.
    \item « \emph{Ne votez pas... sauf si vous voulez la famine} » $\rightarrow$ \textbf{E} (Nul électeur responsable ne vote pour eux) + menace.
\end{itemize}

\noindent\textbf{Pièce 2 :}
\begin{itemize}
    \item « \emph{Monsieur Talla économise 40\%} » $\rightarrow$ \textbf{I} (Certains utilisateurs économisent). \textit{Passage abusif : « MoteurPro+ \textbf{c'est} la garantie »} $\rightarrow$ \textbf{A} (Tout flacon MoteurPro+ est garantie d'économie). \textbf{Généralisation à partir d'un cas unique (M. Talla)}.
    \item « \emph{Approuvé par 9 chauffeurs sur 10} » $\rightarrow$ \textbf{I} (Certains chauffeurs approuvent). Présenté comme quasi-\textbf{A} (Presque tous).
\end{itemize}

\vspace{0.5em}
\noindent\textbf{3. Analyse du raisonnement (Pièce 3 — Marché Mokolo)}

\vspace{0.3em}
\noindent\textbf{Syllogisme 1 (Solidité) :}
\begin{itemize}
    \item \textbf{Majeure (A) :} \emph{Tous les téléphones chers (M) sont solides (P).} [Tout M est P]
    \item \textbf{Mineure (A) :} \emph{Ce Tecno (S) est un téléphone cher (M).} [Tout S est M]
    \item \textbf{Conclusion (A) :} \emph{Donc ce Tecno (S) est solide (P).} [Tout S est P]
\end{itemize}
\textbf{Figure 1 (Barbara) : M-P, S-M $\vdash$ S-P.}
\textbf{Diagramme de Venn :}
\begin{center}
\begin{tikzpicture}[scale=0.9, every node/.style={font=\small}]
    \def\R{1.6}
    \def\dx{1.1}
    \def\dy{0.95}
    \coordinate (S) at (0,0);
    \coordinate (M) at (\dx,0);
    \coordinate (P) at (0.55,\dy);
    % Cercles
    \draw (S) circle (\R) node[below left=-2pt] {S (Ce Tecno)};
    \draw (M) circle (\R) node[below right=-2pt] {M (Tél. chers)};
    \draw (P) circle (\R) node[above=-2pt] {P (Solides)};
    % Hachures : Majeure A (Tout M est P) => M \ P vide
    \begin{scope}
        \clip (M) circle (\R);
        \clip (P) circle (\R);
        % On veut hachurer M \ P. On clippe M, on "soustraite" P via even odd rule ou on hachure tout M puis on blanchit P.
        % Méthode simple : hachurer M, puis remplir P en blanc opaque par dessus (mais masque les traits).
        % Méthode pattern : clip M, clip complement de P ? Pas trivial en TikZ pur sans library.
        % Solution pédagogique : Hachurer la zone M \cap non-P manuellement via clip.
        \clip (P) circle (\R); % On est DANS P
        % On ne fait rien, on veut l'extérieur de P.
    \end{scope}
    % Utilisation de even odd rule pour hachurer M \ P
    \begin{scope}[even odd rule]
        \clip (M) circle (\R) (P) circle (\R);
        \fill[pattern=north east lines, pattern color=popdark] (M) circle (\R);
    \end{scope}
    % Mineure A (Tout S est M) => S \ M vide
    \begin{scope}[even odd rule]
        \clip (S) circle (\R) (M) circle (\R);
        \fill[pattern=north west lines, pattern color=popdark] (S) circle (\R);
    \end{scope}
    % Légende
    \node[popdark, align=left] at (-2.5, -2.0) {Hachures NE : M $\setminus$ P (Maj A)};
    \node[popdark, align=left] at (-2.5, -2.5) {Hachures NO : S $\setminus$ M (Min A)};
    \node[popgreen, align=left] at (2.5, 1.8) {S $\subset$ M $\subset$ P $\Rightarrow$ S $\setminus$ P hachuré : \textbf{Valide}};
\end{tikzpicture}
\legende{Syllogisme 1 (Barbara) : Valide formellement. La zone S hors P est entièrement hachurée $\rightarrow$ Tout S est P.}
\end{center}
\textbf{Validité formelle :} \textbf{VALIDE} (Mode Barbara, Figure 1).
\textbf{Validité matérielle :} \textbf{DOUTEUSE}. La majeure « Tous les téléphones chers sont solides » est fausse (prix $\neq$ robustesse ; marketing, marque, rareté font le prix). Conclusion vraie par hasard possible, mais non garantie par les prémisses.

\vspace{0.3em}
\noindent\textbf{Syllogisme 2 (Garantie) :}
\begin{itemize}
    \item \textbf{Majeure (A) :} \emph{Tous les téléphones solides (M) ont une bonne garantie (P).} [Tout M est P]
    \item \textbf{Mineure (A) :} \emph{Ce Tecno (S) est solide (M).} [Tout S est M] (Issue de la conclusion précédente).
    \item \textbf{Conclusion (A) :} \emph{Donc ce Tecno (S) a une bonne garantie (P).} [Tout S est P]
\end{itemize}
\textbf{Figure 1 (Barbara) : M-P, S-M $\vdash$ S-P.}
\textbf{Validité formelle :} \textbf{VALIDE} (Enchaînement de deux Barbara : Sorites).
\textbf{Validité matérielle :} \textbf{FAUSSE}. La mineure « Ce Tecno est solide » n'est pas prouvée (prémisse majeure du 1er syllogisme fausse). La majeure « Tous les solides ont bonne garantie » est fausse (ex. : marques robustes mais SAV inexistant au Cameroun). \textbf{Erreur : « Affirmation du conséquent » implicite si on raisonne à l'envers, mais ici structure formelle valide, contenu faux.}

\vspace{0.5em}
\noindent\textbf{4. Analyse des principes (Pièce 4 — WhatsApp)}

\vspace{0.3em}
\noindent\textbf{Violation du Principe de Non-Contradiction (PNC) :}
Le texte affirme simultanément :
$P_1$ : « Le MINESEC \textbf{a interdit} tous les téléphones » (Communiqué officiel signé).
$P_2$ : « \textbf{Aucune interdiction n'est prévue} » (Source bien introduite).
$P_1 \equiv \neg P_2$. Le texte affirme $P_1 \land P_2$ (ou les présente comme deux vérités concurrentes « C'est officiel MAIS en réalité... »).
\textbf{Formulation :} $(\text{Interdiction}) \land \neg(\text{Interdiction})$. \textbf{Violation flagrante du PNC.} Un même acte administratif ne peut être à la fois signé et inexistant « en même temps et sous le même rapport » (rentrée 2025, même ministère).

\vspace{0.3em}
\noindent\textbf{Violation du Principe du Tiers Exclu (PTE) / Faux Dilemme :}
« \emph{Soit on interdit tout, soit on autorise tout. Il n'y a pas de milieu.} »
Exclusion de la nuance : \textbf{Réglementation encadrée} (ex. : téléphones autorisés en cours sous surveillance, interdits en examen, déposés à l'entrée). Le texte force un choix binaire artificiel pour manipuler l'émotion.

\vspace{0.3em}
\noindent\textbf{Violation du Principe de Raison Suffisante (PRS) :}
\begin{itemize}
    \item « \emph{Une source bien introduite... nous confirme} » : Aucune source nominative, aucun document cité, aucune preuve accessible. Affirmation gratuite.
    \item « \emph{Le gouvernement nous ment encore une fois} » : Généralisation à partir d'un cas non établi (circulaire).
\end{itemize}

\vspace{0.5em}
\noindent\textbf{5. Verdict Logique (Modèle de rédaction)}

\vspace{0.3em}
\noindent\textbf{VERDICT DU TRIBUNAL DE LA LOGIQUE — AFFAIRE N° 2025/001}
\textit{Objet : Évaluation de la rigueur logique d'un corpus de discours publics (Campagne, Publicité, Commerce, Réseaux sociaux).}

\textbf{Qualification des infractions :}
Le dossier révèle une \textbf{pathologie logique systémique}.
\begin{enumerate}
    \item \textbf{Flou conceptuel (Pièces 1, 2) :} Les termes « gestion », « malédiction », « économie 40\% », « approuvé » sont utilisés en \emph{termes-valises} ou \emph{mots-passe-partout}, vidés de leur intension rigoureuse, pour déclencher l'affect plutôt qu'informer l'entendement.
    \item \textbf{Généralisations hâtives (Pièces 1, 2) :} Passage illégitime du singulier (un meeting pluvieux, un chauffeur Talla) à l'universel (tous les meetings, tous les flacons). Violation de la règle d'extension du jugement particulier à l'universel sans induction complète.
    \item \textbf{Validité formelle sans vérité matérielle (Pièce 3) :} Le vendeur construit des syllogismes \emph{formellement valides} (Barbara) mais reposant sur des \emph{majeures fausses} (Prix $\Rightarrow$ Solidité ; Solidité $\Rightarrow$ Garantie). C'est la \emph{sophistique marchande} : l'apparence de la rigueur (structure) masque l'absence de fondement (contenu).
    \item \textbf{Contradiction performative et Faux dilemme (Pièce 4) :} Le message WhatsApp s'autodétruit logiquement (PNC) tout en forçant un choix binaire exclu (PTE) pour court-circuiter l'esprit critique. L'absence de raison suffisante (source anonyme) achève de le disqualifier comme information.
\end{enumerate}

\textbf{Gravité pour le débat public :}
Ces infractions ne sont pas de simples « fautes de français ». Elles structurent la \emph{malinformation} : elles empêchent le citoyen de se forger un jugement éclairé (vote, achat, éducation). La confusion conceptuelle entretient la superstition politique ; le syllogisme fallacieux ruine le consommateur ; la contradiction virale paralyse l'action parentale.

\textbf{Recommandations :}
1) \textbf{Exiger la définition} : « Que voulez-vous dire exactement par \textit{malédiction} / \textit{économie 40\%} ? »
2) \textbf{Tester l'universel} : « Avez-vous vérifié \textit{tous} les meetings / \textit{tous} les utilisateurs ? »
3) \textbf{Vérifier les prémisses} : « Est-il vrai que \textit{tout} téléphone cher est solide ? »
4) \textbf{Traquer la contradiction} : « Les deux affirmations peuvent-elles être vraies ensemble ? »
5) \textbf{Exiger la source} : « Quel document officiel ? Quel nom, quelle date ? »

\textit{Fait à Douala-Bassa, le 15 mai 2025. Le Jury : [Noms des élèves].}
\end{exoresolu}

\courssub{Bilan de l'activité}
Cette activité a permis de mettre en évidence que la \cle{logique classique} n'est pas un formalisme abstrait déconnecté du réel, mais un \textbf{outil de survie intellectuelle} dans l'espace public camerounais contemporain.

\begin{aretenir}[Acquis majeurs]
\begin{itemize}
    \item \textbf{Concept :} L'imprécision définitionnelle (genre/différence spécifique) est la porte d'entrée de la manipulation (ex. « malédiction », « approuvé »).
    \item \textbf{Jugement :} Le glissement systématique I $\rightarrow$ A (particulier $\rightarrow$ universel) est la fallacie n°1 des discours politiques et publicitaires au Cameroun.
    \item \textbf{Raisonnement :} La \emph{validité formelle} (structure Barbara) ne garantit pas la \emph{vérité} (validité matérielle). Un syllogisme valide avec des prémisses fausses mène à n'importe quelle conclusion (ex. garantie du téléphone).
    \item \textbf{Principes :} La violation du PNC (Pièce 4) rend le discours \emph{inintelligible} (on ne sait plus quoi croire) ; la violation du PTE (faux dilemme) enferme dans l'alternative stérile ; l'absence de PRS (source anonyme) interdit toute vérification.
    \item \textbf{Compétence transférable :} La \textbf{Grille d'analyse logique} (5 critères : Clarté conceptuelle, Justesse des jugements, Validité du raisonnement, Respect des principes, Preuve/Source) est désormais un \emph{réflexe épistémique} pour tout élève face à un discours médiatique, un sujet de dissertation ou une situation professionnelle technique.
\end{itemize}
\end{aretenir}

\begin{attention}[Pièges évités grâce à l'activité]
\begin{itemize}
    \item Confondre \emph{validité} (liaison nécessaire prémisses-conclusion) et \emph{vérité} (conformité au réel). Le vendeur de Mokolo exploitait cette confusion.
    \item Croire qu'un diagramme de Venn « prouve » la vérité de la conclusion : il ne prouve que la \emph{forme}.
    \item Négliger la \emph{différence spécifique} dans la définition : définir « téléphone cher » par « qui coûte cher » (cercle vicieux) au lieu de « dont le prix dépasse la moyenne du marché pour des caractéristiques techniques données ».
    \item Accepter le \emph{faux dilemme} (PTE) comme une contrainte réelle : « Soit interdiction totale, soit anarchie » $\rightarrow$ Faux. La régulation est le tiers exclu réel.
\end{itemize}
\end{attention}

\begin{savaistu}[Logique et Citoyenneté numérique au Cameroun]
L'ANOR (Agence des Normes et de la Qualité) et l'ANTIC (Agence Nationale des Technologies de l'Information et de la Communication) luttent contre la désinformation. La \textbf{logique formelle} est la première ligne de défense : un citoyen qui repère une violation du PNC ou une généralisation hâtive (I $\rightarrow$ A) ne partage pas la \emph{fake news}. Au Probatoire, l'épreuve de dissertation (ex. : « La technique libère-t-elle l'homme ? ») exige une \emph{définition rigoureuse des concepts} (technique, liberté), des \emph{jugements nuancés} (éviter « La technique libère toujours » $\rightarrow$ A faux), un \emph{raisonnement valide} (pas de saut logique) et le \emph{respect des principes} (ne pas se contredire entre thèse et antithèse). Cette activité \textbf{est} l'entraînement direct à l'épreuve du Bac.
\end{savaistu}

\newpage

\coursec{Méthodes et exercices résolus}

\coursec{La logique classique}

\courssub{Définition et importance de la logique}

\begin{definition}[Logique]
La \cle{logique} (du grec \emph{logos} : raison, discours, science) est la science des lois de la pensée correcte. Elle étudie les opérations de l'esprit (concept, jugement, raisonnement) et les règles qui garantissent la validité des déductions, c'est-à-dire le passage nécessaire de prémisses vraies à une conclusion vraie.
\end{definition}

\begin{aretenir}[Objet et utilité de la logique]
\begin{itemize}
\item \textbf{Objet formel :} les lois de la pensée (identité, non-contradiction, tiers exclu) et les formes du raisonnement (syllogisme, induction).
\item \textbf{Objet matériel :} tout contenu de pensée soumis à ces lois.
\item \textbf{Importance :}
\begin{enumerate}
\item \emph{Critique :} distinguer le vrai du faux, le valide de l'invalide, démasquer les sophismes (publicité, rumeurs, débats politiques).
\item \emph{Constructive :} ordonner ses idées, définir rigoureusement, argumenter sans contradiction, rédiger une démonstration.
\item \emph{Universelle :} elle n'est pas réservée aux mathématiciens ; tout citoyen, commerçant, juriste, élève ou chef de famille a besoin de penser juste pour décider et agir.
\end{enumerate}
\end{itemize}
\end{aretenir}

\begin{savaistu}[Histoire brève]
Aristote (IV\textsuperscript{e} s. av. J.-C.) fonde la logique formelle (\emph{Organon}) : théorie du syllogisme, principes d'identité et de non-contradiction. Au XVII\textsuperscript{e} s., l'école de Port-Royal (\emph{Logique de Port-Royal}, 1662) renouvelle l'analyse du concept (compréhension/extension). Au XIX\textsuperscript{e} s., Boole et Frege créent la logique symbolique (algèbre de la logique, calcul des prédicats), base de l'informatique moderne. Au Cameroun, l'enseignement de la logique au lycée technique vise à former des techniciens capables de rigueur dans l'argumentation technique et la vie civique.
\end{savaistu}

\courssub{Le concept}

\begin{definition}[Concept]
Le \cle{concept} est une idée générale et abstraite qui représente une classe d'objets partageant des caractères communs. C'est l'opération par laquelle l'esprit saisit \emph{ce qu'est} une chose (son essence).
\end{definition}

\begin{propriete}[Compréhension et extension d'un concept]
\begin{itemize}
\item \textbf{Compréhension :} ensemble des caractères (propriétés) communs à tous les objets de la classe. C'est le \emph{contenu} du concept. Ex. : « triangle » = figure plane, trois côtés, trois angles.
\item \textbf{Extension :} ensemble des objets ou des espèces auxquels le concept s'applique. C'est la \emph{portée} du concept. Ex. : « triangle » = triangles rectangles, isocèles, équilatéraux, quelconques.
\item \textbf{Loi fondamentale (rapport inverse) :} Plus la compréhension est riche (nombreux caractères), plus l'extension est pauvre (peu d'objets), et inversement. Ex. : « être vivant » (compréhension pauvre, extension immense) $\rightarrow$ « mammifère » (compréhension plus riche, extension réduite) $\rightarrow$ « chien » (encore plus riche, extension encore plus réduite).
\end{itemize}
\end{propriete}

\begin{methode}[Analyser un concept (compréhension et extension)]
\begin{enumerate}
\item \textbf{Dégager la définition} : dire ce que le concept désigne (genre prochain + différence spécifique).
\item \textbf{Déterminer la compréhension} : lister les caractères essentiels communs à tous les objets.
\item \textbf{Déterminer l'extension} : énumérer les espèces ou individus concernés (dans le contexte précis).
\item \textbf{Appliquer la loi du rapport inverse} : vérifier qu'en ajoutant un caractère (restriction), l'extension diminue.
\item \textbf{Conclure} sur l'intérêt : éviter les équivoques dans un contrat, un règlement, un débat.
\end{enumerate}
\textbf{Réflexe Bac :} Une définition correcte n'est ni \emph{trop large} (inclut des étrangers : « l'homme est un animal ») ni \emph{trop étroite} (exclut des membres : « l'élève est un jeune en uniforme »).
\end{methode}

\begin{exoresolu}[Analyse du concept « élève »]
\textbf{Énoncé.} Un proviseur déclare : « Tous les élèves du lycée doivent porter l'uniforme. » Un élève d'une école privée voisine prétend être concerné. (1) Définissez le concept d'« élève ». (2) Donnez sa compréhension et son extension dans la phrase du proviseur. (3) L'élève de l'école voisine est-il concerné ? Justifiez.
\tcblower
\textbf{Solution détaillée.}
\begin{enumerate}
\item \textbf{Définition.} Un élève est une personne inscrite dans un établissement scolaire et qui y reçoit un enseignement structuré sous l'autorité d'un corps professoral.
\item \textbf{Compréhension et extension dans l'énoncé.}
\begin{itemize}
\item \emph{Compréhension} (caractères) : personne inscrite, recevant un enseignement, soumise au règlement intérieur de \emph{cet} établissement.
\item \emph{Extension} visée : l'ensemble des élèves \textbf{de ce lycée précis} (restriction par le complément du nom « du lycée »).
\end{itemize}
\item \textbf{Application.} L'élève de l'école voisine possède la compréhension \emph{générale} du concept « élève », mais il n'appartient pas à l'\emph{extension contextuelle} définie par le proviseur. Le caractère ajouté (« de ce lycée ») enrichit la compréhension et réduit l'extension (loi du rapport inverse).
\end{enumerate}
\textbf{Conclusion.} L'élève de l'école voisine n'est pas concerné. Cet exemple montre que préciser l'extension d'un concept (ici par un complément du nom) évite les conflits d'interprétation dans la vie administrative et sociale.
\end{exoresolu}

\courssub{Le jugement}

\begin{definition}[Jugement]
Le \cle{jugement} est l'opération de l'esprit qui affirme ou nie un attribut (prédicat) d'un sujet. Il unit ou sépare deux concepts par la copule (verbe « être » ou équivalent). Ex. : « L'or est un métal » (affirmation) ; « Le bois n'est pas un métal » (négation).
\end{definition}

\begin{propriete}[Structure et formes canoniques (A, E, I, O)]
Tout jugement se décompose en : \textbf{Sujet (S)} -- \textbf{Copule (est / n'est pas)} -- \textbf{Attribut (P)}.
On distingue deux dimensions :
\begin{itemize}
\item \textbf{Qualité :} \emph{Affirmatif} (S est P) ou \emph{Négatif} (S n'est pas P).
\item \textbf{Quantité :} \emph{Universel} (Tous les S / Aucun S) ou \emph{Particulier} (Certains S / Certains S ne sont pas P).
\end{itemize}
Croisant qualité et quantité, on obtient les \textbf{quatre formes logiques} :
\begin{center}
\begin{tabular}{|c|c|c|c|}\hline
 & \textbf{Universel} & \textbf{Particulier} \\ \hline
\textbf{Affirmatif} & \textbf{A} : Tout S est P & \textbf{I} : Quelque S est P \\ \hline
\textbf{Négatif} & \textbf{E} : Nul S n'est P & \textbf{O} : Quelque S n'est pas P \\ \hline
\end{tabular}
\end{center}
\textbf{Relations logiques (Carré d'opposition) :}
\begin{itemize}
\item \emph{Contradictoires} (ne peuvent être vrais ensemble ni faux ensemble) : A $\leftrightarrow$ O ; E $\leftrightarrow$ I.
\item \emph{Contraires} (ne peuvent être vrais ensemble, peuvent être faux ensemble) : A $\leftrightarrow$ E.
\item \emph{Subcontraires} (peuvent être vrais ensemble, ne peuvent être faux ensemble) : I $\leftrightarrow$ O.
\item \emph{Subalternation} : A $\rightarrow$ I (si A vrai, I vrai) ; E $\rightarrow$ O.
\end{itemize}
\end{propriete}

\begin{methode}[Identifier la forme d'un jugement (A, E, I, O)]
\begin{enumerate}
\item Repérer le \textbf{Sujet (S)}, la \textbf{Copule} (verbe « être » mis au présent) et l'\textbf{Attribut (P)}. Reformuler si nécessaire : « Paul mange » $\rightarrow$ « Paul est mangeur ».
\item Déterminer la \textbf{Qualité} : copule affirmative (est) ou négative (n'est pas).
\item Déterminer la \textbf{Quantité} : quantificateur universel (tous, aucun, nul) ou particulier (certains, quelques-uns, un).
\item Conclure par la \textbf{Forme} : A, E, I ou O.
\item Vérifier les \textbf{contradictions} : A contredit O ; E contredit I.
\end{enumerate}
\end{methode}

\begin{exoresolu}[Analyse de quatre jugements camerounais]
\textbf{Énoncé.} Analysez (Sujet, Copule, Attribut, Qualité, Quantité, Forme) : (a) « Tous les commerçants du marché de Mokolo paient leurs impôts. » (b) « Aucun élève sérieux ne triche. » (c) « Certains jeunes de Douala sont entrepreneurs. » (d) « Certains fonctionnaires ne sont pas ponctuels. » Indiquez les couples contradictoires.
\tcblower
\textbf{Solution détaillée.}
\begin{itemize}
\item \textbf{(a)} S = « commerçants du marché de Mokolo » ; Copule = « sont » (reformulation : « sont payeurs ») ; P = « payeurs de leurs impôts ». Qualité : affirmative. Quantité : universelle (« tous »). \textbf{Forme A.}
\item \textbf{(b)} S = « élèves sérieux » ; Copule = « ne sont pas » ; P = « tricheurs ». Qualité : négative. Quantité : universelle (« aucun »). \textbf{Forme E.}
\item \textbf{(c)} S = « jeunes de Douala » ; Copule = « sont » ; P = « entrepreneurs ». Qualité : affirmative. Quantité : particulière (« certains »). \textbf{Forme I.}
\item \textbf{(d)} S = « fonctionnaires » ; Copule = « ne sont pas » ; P = « ponctuels ». Qualité : négative. Quantité : particulière. \textbf{Forme O.}
\end{itemize}
\textbf{Contradictions.} A (a) contredit O (« Certains commerçants de Mokolo ne paient pas leurs impôts »). E (b) contredit I (« Certains élèves sérieux trichent »).
\textbf{Conclusion.} Ramener une phrase du langage courant à sa forme canonique A/E/I/O permet de voir immédiatement quelle proposition la réfute logiquement.
\end{exoresolu}

\courssub{Le raisonnement}

\begin{definition}[Raisonnement et syllogisme]
Le \cle{raisonnement} est l'opération qui tire une conclusion de prémisses. Le \cle{syllogisme} (raisonnement déductif type) comporte \textbf{deux prémisses} et \textbf{une conclusion}, et exactement \textbf{trois termes} :
\begin{itemize}
\item \textbf{Grand terme (P) :} attribut de la conclusion.
\item \textbf{Petit terme (S) :} sujet de la conclusion.
\item \textbf{Moyen terme (M) :} terme présent dans les deux prémisses, absent de la conclusion ; il assure la liaison entre S et P.
\end{itemize}
\end{definition}

\begin{propriete}[Règles de validité du syllogisme (catégorique)]
Un syllogisme est \textbf{valide} (forme correcte) s'il respecte \emph{toutes} les règles suivantes :
\begin{enumerate}
\item \textbf{Règle des trois termes :} Trois termes exactement, chacun pris dans le même sens (sinon : sophisme des quatre termes / équivocation).
\item \textbf{Règle du moyen terme :} Le moyen terme M doit être pris \emph{au moins une fois dans toute son extension} (c'est-à-dire comme sujet d'une universelle : « Tous M » ou « Aucun M »). Si M est attribut d'une affirmative, il n'est pas pris universellement.
\item \textbf{Règles de qualité :}
\begin{itemize}
\item De deux prémisses négatives, on ne peut rien conclure.
\item Si une prémisse est négative, la conclusion doit être négative.
\item De deux prémisses affirmatives, la conclusion est affirmative.
\end{itemize}
\item \textbf{Règles de quantité :}
\begin{itemize}
\item De deux prémisses particulières, on ne peut rien conclure.
\item Si une prémisse est particulière, la conclusion doit être particulière.
\item Aucun terme ne peut être pris dans la conclusion dans une extension plus grande que dans les prémisses (S et P ne peuvent être universels en conclusion s'ils sont particuliers dans les prémisses).
\end{itemize}
\end{enumerate}
\textbf{Distinction cruciale :} \emph{Validité} = forme correcte (la conclusion suit nécessairement des prémisses). \emph{Vérité} = conformité au réel. Un raisonnement valide peut avoir des prémisses fausses ; un raisonnement invalide peut avoir une conclusion vraie par hasard.
\end{propriete}

\begin{methode}[Vérifier la validité d'un syllogisme]
\begin{enumerate}
\item Identifier la \textbf{conclusion} (souvent après « donc », « par conséquent ») et les deux \textbf{prémisses}.
\item Repérer S, P, M. Vérifier qu'il n'y a que trois termes, univoques.
\item Vérifier que \textbf{M est universel au moins une fois} (sujet d'une A ou d'une E).
\item Vérifier les règles de \textbf{qualité} et de \textbf{quantité}.
\item Conclure : \textbf{Valide} ou \textbf{Invalide}. Préciser : « Le raisonnement est valide, mais les prémisses sont fausses » ou « Invalide (moyen terme non distribué), la conclusion ne s'impose pas ».
\end{enumerate}
\end{methode}

\begin{exoresolu}[Validité de deux syllogismes]
\textbf{Énoncé.} Examinez la validité :
(a) « Tous les mammifères sont des vertébrés ; or tous les chiens sont des mammifères ; donc tous les chiens sont des vertébrés. »
(b) « Tous les élèves de Première technique étudient la philosophie ; or Paul étudie la philosophie ; donc Paul est élève de Première technique. »
\tcblower
\textbf{Solution détaillée.}

\textbf{(a)} Conclusion : « Tous les chiens sont des vertébrés ». S = chiens ; P = vertébrés ; M = mammifères.
\begin{itemize}
\item 3 termes, univoques : OK.
\item M = « mammifères » est sujet de la 1\textsuperscript{re} prémisse (« \emph{Tous} les mammifères ») $\rightarrow$ M pris universellement : OK.
\item 2 prémisses A, conclusion A : règles qualité/quantité respectées.
\end{itemize}
\textbf{Valide} (Mode Barbara, 1\textsuperscript{re} figure). Prémisses vraies $\rightarrow$ conclusion vraie.

\textbf{(b)} Conclusion : « Paul est élève de Première technique ». S = Paul ; P = élève de 1\textsuperscript{re} tech ; M = étudier la philosophie.
\begin{itemize}
\item 3 termes, univoques : OK.
\item M = « étudier la philosophie » : dans P1 (« Tous les élèves... étudient »), M est attribut d'une affirmative $\rightarrow$ non universel. Dans P2 (« Paul étudie »), M est attribut d'une affirmative singulière $\rightarrow$ non universel. \textbf{M jamais universel.}
\end{itemize}
\textbf{Invalide} (Sophisme : moyen terme non distribué / affirmation du conséquent). Paul peut être en Terminale ou à l'université. La conclusion ne s'impose pas.
\end{exoresolu}

\courssub{Les principes logiques de la pensée}

\begin{definition}[Les trois principes premiers]
Ce sont les conditions \emph{a priori} de toute pensée cohérente. Ils s'appliquent \textbf{en même temps et sous le même rapport}.
\begin{enumerate}
\item \textbf{Principe d'identité :} « A est A ». Une chose est ce qu'elle est ; un terme doit garder le même sens tout au long du discours. Violation = \emph{équivocation} (changement de sens).
\item \textbf{Principe de non-contradiction :} « A n'est pas non-A » / « Il est impossible que A soit et ne soit pas en même temps et sous le même rapport ». Deux jugements contradictoires (A et O ; E et I) ne peuvent être vrais ensemble.
\item \textbf{Principe du tiers exclu :} « A est B ou A n'est pas B » / « Entre deux jugements contradictoires, il n'y a pas de troisième possibilité ; l'un est vrai, l'autre faux ». Attention : ne pas confondre \emph{contradictoires} (A/O, E/I) et \emph{contraires} (A/E : entre « tout » et « aucun », il y a « certains »).
\end{enumerate}
\end{definition}

\begin{methode}[Dépister les violations des principes dans un discours]
\begin{enumerate}
\item Citrer le passage exact qui pose problème.
\item Identifier le principe concerné : Identités (sens changeant), Non-contradiction (affirmation + négation simultanées), Tiers exclu (faux dilemme ou exclusion du milieu à tort).
\item Expliquer en une phrase \emph{pourquoi} le principe est violé (préciser « en même temps et sous le même rapport » pour la non-contradiction).
\item Proposer la correction logique.
\end{enumerate}
\end{methode}

\begin{exoresolu}[Dépistage des violations]
\textbf{Énoncé.} Indiquez si chaque énoncé viole un principe ; lequel ? Justifiez.
(a) « Ce téléphone est à la fois en vente et n'est pas en vente dans cette boutique. »
(b) « La démocratie, c'est le pouvoir du peuple ; or le peuple peut déléguer son pouvoir ; donc la démocratie est le pouvoir des délégués. »
(c) « De deux choses l'une : ou ce candidat est honnête, ou il ne l'est pas ; il faut choisir. »
\tcblower
\textbf{Solution détaillée.}

\textbf{(a)} Affirme « A » (en vente) et « non-A » (pas en vente) simultanément, même lieu, même moment. Viole le \textbf{principe de non-contradiction}. Nécessairement faux.

\textbf{(b)} Le terme « pouvoir » change de sens : P1 = « souveraineté » (pouvoir constituant) ; P2 = « mandat/délégation » (pouvoir constitué). Violation du \textbf{principe d'identité} (équivocation). Conclusion abusive.

\textbf{(c)} Pose l'alternative entre « honnête » (A) et « non honnête » (O). Ce sont des contradictoires. Application correcte du \textbf{principe du tiers exclu}. Aucune violation. (Note : « honnête » et « malhonnête » sont contraires, mais l'énoncé utilise bien la négation contradictoire « ne l'est pas »).
\end{exoresolu}

\courssub{Application : penser juste dans la vie quotidienne}

\begin{methode}[Traiter une situation concrète avec les outils logiques]
Face à un discours (marchandage, rumeur, débat, message WhatsApp, publicité) :
\begin{enumerate}
\item \textbf{Décrire la situation} et isoler le raisonnement implicite ou explicite.
\item \textbf{Analyser les concepts} : sont-ils définis ? Leur extension est-elle précise (pas de glissement) ?
\item \textbf{Formaliser les jugements} : les ramener aux formes A, E, I, O pour tester leur vérité.
\item \textbf{Reconstituer le syllogisme} : identifier prémisses, conclusion, termes S, P, M.
\item \textbf{Vérifier la validité} (règles du syllogisme) et les \textbf{principes} (identité, non-contradiction, tiers exclu).
\item \textbf{Produire un jugement argumenté} : le discours est-il logiquement recevable ? Si non, formuler la version correcte (probable, conditionnelle, ou rejet).
\end{enumerate}
\end{methode}

\begin{exoresolu}[La rumeur du marché de Bafoussam]
\textbf{Énoncé.} « Tous les produits importés sont chers ; or ce sac de riz est cher ; donc ce sac de riz est un produit importé. » Évaluez ce raisonnement.
\tcblower
\textbf{Solution détaillée.}
\begin{enumerate}
\item \textbf{Situation.} Un consommateur infère l'origine importée d'un produit à partir de son prix.
\item \textbf{Concepts.} « Produit importé » (compréhension : origine étrangère ; extension : riz, huile, médicaments importés...). « Produit cher » (compréhension : prix élevé ; extension : produits importés \emph{mais aussi} produits locaux rares, de qualité, ou taxés). Les extensions ne coïncident pas.
\item \textbf{Jugements.} P1 : « Tout produit importé est cher » (A). P2 : « Ce sac de riz est cher » (A singulière). C : « Ce sac de riz est importé » (A singulière).
\item \textbf{Validité.} M = « cher ». Dans P1 (A), « cher » est attribut affirmatif $\rightarrow$ non universel. Dans P2 (A), « cher » est attribut affirmatif $\rightarrow$ non universel. \textbf{M jamais pris universellement.} Raisonnement \textbf{invalide} (sophisme : affirmation du conséquent / moyen terme non distribué).
\item \textbf{Principes.} Pas de contradiction formelle, mais glissement sémantique : « cher » utilisé comme s'il était l'attribut exclusif des produits importés (violation de l'identité du concept « cher » dans ce contexte).
\item \textbf{Jugement argumenté.} La conclusion est \emph{possible} mais non \emph{nécessaire}. Le riz peut être local et cher (pénurie, qualité supérieure). Formulation correcte : « Ce sac est cher ; or \emph{beaucoup} de produits importés sont chers ; donc ce sac \emph{peut être} importé » (raisonnement inductif/analogique, conclusion probable).
\end{enumerate}
\textbf{Conclusion.} La logique protège le consommateur : le prix ne prouve pas l'origine. Seule la vérification de l'étiquette (traçabilité) tranche.
\end{exoresolu}

\begin{methode}[Rédiger une réponse de philosophie logique au Probatoire / Bac Tech]
\begin{enumerate}
\item \textbf{Annoncer le plan :} « Pour répondre, nous définirons d'abord les notions clés, puis nous analyserons le raisonnement, enfin nous conclurons. »
\item \textbf{Définir les notions} mobilisées (concept, jugement, syllogisme, validité, principes) \emph{avant} de les appliquer. Une définition précise vaut des points.
\item \textbf{Justifier chaque étape} en nommant la règle : « D'après la règle du moyen terme... », « En vertu du principe de non-contradiction... ».
\item \textbf{Illustrer par un exemple concret} camerounais (marché, école, administration, transport).
\item \textbf{Rédiger une conclusion en phrase complète} qui répond \emph{exactement} à la question, sans idée nouvelle.
\item \textbf{Relire} : vérifier l'univocité des termes (principe d'identité appliqué à sa propre copie).
\end{enumerate}
\end{methode}

\begin{exoresolu}[Rédaction d'une réponse argumentée : « La logique ne sert qu'aux mathématiciens »]
\textbf{Énoncé.} Rédigez une réponse justifiée (environ 15 lignes) en suivant la démarche complète.
\tcblower
\textbf{Solution détaillée (corrigé rédigé).}

La logique est la science des lois de la pensée correcte ; elle étudie les conditions de validité du concept, du jugement et du raisonnement. Affirmer qu'elle ne sert qu'aux mathématiciens, c'est en réduire abusivement l'extension en confondant l'usage technique d'un outil avec sa fonction universelle.

Or la logique est d'abord l'instrument de tout homme qui pense et qui parle. En vertu du \textbf{principe de non-contradiction}, nul ne peut soutenir sérieusement une thèse et son contraire : le commerçant qui fixe ses prix, le juriste qui plaide, l'élève qui compose, le chef de famille qui décide doivent éviter la contradiction pour être crédibles. De même, la maîtrise du \textbf{raisonnement} permet de distinguer une démonstration valide d'un sophisme : celui qui sait que le \textbf{moyen terme doit être pris universellement} ne se laissera pas tromper par l'argument publicitaire « Les grandes marques sont chères ; ce produit est cher ; donc c'est une grande marque » (moyen terme « cher » non distribué).

Un exemple camerounais l'illustre : lors des débats électoraux ou des assemblées générales de groupements d'initiative commune (GIC), le citoyen qui analyse les concepts employés (« développement », « transparence », « justice ») et qui vérifie la validité des raisonnements des candidats ou des dirigeants participe en connaissance de cause. La logique forme ainsi le \textbf{jugement critique}, condition de la vie démocratique et de la bonne gouvernance locale.

En conclusion, la logique n'est pas le privilège des mathématiciens : elle est l'instrument universel de la pensée juste, utile au marché comme au tribunal, à l'école comme à la famille. Dire qu'elle ne sert qu'aux mathématiciens, c'est méconnaître sa vocation anthropologique : \emph{organiser la raison pour agir avec justesse}.
\end{exoresolu}

\begin{attention}[Les 5 erreurs qui coûtent des points au Probatoire / Bac Tech]
\begin{enumerate}
\item \textbf{Confondre validité et vérité.} Un raisonnement peut être \emph{valide} avec des prémisses fausses, et \emph{invalide} avec une conclusion vraie par hasard. Le correcteur attend la distinction \emph{forme} (validité) / \emph{contenu} (vérité). N'écrivez jamais « le raisonnement est vrai » pour « valide ».
\item \textbf{Oublier de vérifier le moyen terme.} L'erreur classique : « Tous les A sont B ; C est B ; donc C est A ». Le moyen terme B, attribut d'une prémisse affirmative, n'est \emph{jamais} pris universellement. Toujours vérifier : M est-il sujet d'une universelle (A ou E) au moins une fois ?
\item \textbf{Définir un concept trop largement ou trop étroitement.} « L'homme est un animal » (trop large : inclut les bêtes) ou « L'élève est un jeune en uniforme » (trop étroit : exclut les élèves sans uniforme). Exigez la compréhension complète (genre + différence spécifique) et l'extension exacte.
\item \textbf{Appliquer la non-contradiction sans « en même temps et sous le même rapport ».} Dire « Ce magasin est ouvert et fermé » n'est contradictoire que si c'est au même moment et au même sens (ex. : porte ouverte / rideau baissé). Oublier cette précision invalide la justification.
\item \textbf{Conclure hors sujet ou par une opinion non argumentée.} La conclusion doit reprendre les termes exacts du sujet, en une phrase complète, justifiée par les étapes précédentes. Une conclusion « à l'emporte-pièce » annule le bénéfice de la démarche.
\end{enumerate}
\end{attention}

\newpage

\coursec{Exercices}

\courssub{Je vérifie mes connaissances}

\begin{exercice}[QCM : notions fondamentales]
Parmi les propositions suivantes, une seule est exacte par question. Recopiez le numéro de la bonne réponse.
\begin{enumerate}
    \item La logique formelle étudie principalement :
    \begin{enumerate}
        \item La vérité matérielle des prémisses.
        \item La structure et la validité du raisonnement.
        \item Le contenu psychologique de la pensée.
        \item L'origine historique des concepts.
    \end{enumerate}
    \item L'\cle{extension} d'un concept désigne :
    \begin{enumerate}
        \item L'ensemble des attributs essentiels qui le définissent.
        \item L'ensemble des objets auxquels le concept s'applique.
        \item La relation entre le sujet et l'attribut dans un jugement.
        \item La forme logique d'un syllogisme.
    \end{enumerate}
    \item Le jugement « Certains élèves de Première Technique ne sont pas boursiers » est de type :
    \begin{enumerate}
        \item A (universelle affirmative).
        \item E (universelle negative).
        \item I (particuliere affirmative).
        \item O (particuliere negative).
    \end{enumerate}
    \item Le principe selon lequel « une chose ne peut pas être et ne pas être en même temps et sous le même rapport » est :
    \begin{enumerate}
        \item Le principe d'identité.
        \item Le principe de non-contradiction.
        \item Le principe du tiers exclu.
        \item Le principe de raison suffisante.
    \end{enumerate}
\end{enumerate}
\end{exercice}

\begin{exercice}[Vrai ou Faux justifié]
Indiquez si chaque affirmation est vraie (V) ou fausse (F) et justifiez brièvement votre réponse en vous appuyant sur le cours.
\begin{enumerate}
    \item Un raisonnement valide sur le plan formel conduit nécessairement à une conclusion vraie dans la réalité.
    \item La compréhension d'un concept varie en sens inverse de son extension : plus le concept est étendu, plus il est pauvre en attributs.
    \item Tout syllogisme qui respecte la figure 1 (M-P, S-M, S-P) est automatiquement valide.
    \item Le principe du tiers exclu permet d'affirmer que toute proposition est soit vraie, soit fausse, sans troisième possibilité.
\end{enumerate}
\end{exercice}

\begin{exercice}[Définitions attendues]
Donnez une définition précise et philosophique des notions suivantes :
\begin{enumerate}
    \item Le \cle{concept} (en distinguant compréhension et extension).
    \item Le \cle{jugement} (en précisant sa structure sujet-copule-attribut).
    \item Le \cle{raisonnement} (en distinguant déduction et induction).
    \item Les trois \cle{principes logiques de la pensée} (identité, non-contradiction, tiers exclu).
\end{enumerate}
\end{exercice}

\begin{exercice}[Classification logique et carrés d'opposition]
\begin{enumerate}
    \item Classez chacune des propositions suivantes selon la quantité (universelle/particulière) et la qualité (affirmative/négative) en indiquant sa lettre canonique (A, E, I, O) :
    \begin{itemize}
        \item « Tous les ingénieurs ont suivi une formation scientifique. »
        \item « Aucune machine n'est un être vivant. »
        \item « Certains logiciels sont libres. »
        \item « Certains téléphones portables ne sont pas des smartphones. »
    \end{itemize}
    \item Pour la proposition A « Tous les élèves de 1re Tech étudient la philosophie », écrivez les trois propositions correspondantes E, I, O et précisez le type d'opposition (contraire, contradictoire, sous-alterne, sous-contrarie) qui lie A à chacune d'elles.
\end{enumerate}
\end{exercice}

\courssub{Je m'entraîne}

\begin{exercice}[Le carré d'opposition en situation]
On considère les quatre propositions suivantes :
\begin{itemize}
    \item P1 : « Tous les candidats au Probatoire sont inscrits dans un établissement. »
    \item P2 : « Aucun candidat au Probatoire n'est dispensé d'épreuves. »
    \item P3 : « Certains candidats au Probatoire sont des répétants. »
    \item P4 : « Certains candidats au Probatoire ne sont pas admis. »
\end{itemize}
\begin{enumerate}
    \item Identifiez le type (A, E, I, O) de chaque proposition.
    \item Construisez le carré d'opposition complet en reliant P1, P2, P3, P4 par les flèches d'opposition appropriées (contraire, contradictoire, sous-alterne, sous-contrarie).
    \item Si P1 est vraie, que peut-on déduire de la valeur de vérité de P2, P3, P4 ? Justifiez par les lois du carré.
    \item Si P3 est fausse, que devient la valeur de vérité des trois autres ?
\end{enumerate}
\end{exercice}

\begin{exercice}[Validité des syllogismes : règles et diagrammes de Venn]
Vérifiez la validité formelle des syllogismes suivants en appliquant les règles générales (terme moyen, distribution, prémisses négatives, etc.) et en traçant le diagramme de Venn correspondant.
\begin{enumerate}
    \item \textbf{Syllogisme 1 :}
    \begin{itemize}
        \item Tous les mammifères sont des vertébrés.
        \item Or, la baleine est un mammifère.
        \item Donc, la baleine est un vertébré.
    \end{itemize}
    \item \textbf{Syllogisme 2 :}
    \begin{itemize}
        \item Tous les fonctionnaires sont assidus.
        \item Or, Paul est assidu.
        \item Donc, Paul est fonctionnaire.
    \end{itemize}
    \item \textbf{Syllogisme 3 :}
    \begin{itemize}
        \item Aucun corrompu n'est honnête.
        \item Or, certains politiciens sont corrompus.
        \item Donc, certains politiciens ne sont pas honnêtes.
    \end{itemize}
\end{enumerate}
Pour chaque syllogisme : indiquez l'humeur et la figure, citez la/les règle(s) respectée(s) ou violée(s), et concluez sur la validité.
\end{exercice}

\begin{exercice}[Critique et correction de définitions]
Les définitions suivantes sont défectueuses. Pour chacune : identifiez le(s) vice(s) de définition (trop large, trop étroite, circulaire, obscure, négative là où l'affirmative est possible, etc.), puis proposez une définition correcte.
\begin{enumerate}
    \item « Un triangle est une figure qui a trois côtés. »
    \item « La liberté, c'est faire ce qu'on veut. »
    \item « Un menteur est celui qui dit des mensonges. »
    \item « L'école est un bâtiment où vont les enfants. »
\end{enumerate}
\end{exercice}

\begin{exercice}[Application des principes logiques à l'argumentation courante]
M. Kamga, délégué de classe, déclare lors du conseil de classe : « Si nous accordons la bourse d'excellence à Amadou, il faudra l'accorder à tout le monde, car Amadou n'est pas meilleur que les autres. Or, le budget ne permet pas de donner une bourse à tous. Donc, il ne faut pas la donner à Amadou. »
\begin{enumerate}
    \item Formalisez ce raisonnement (identifiez les prémisses et la conclusion, mettez en évidence la structure conditionnelle).
    \item Quel principe logique (identité, non-contradiction, tiers exclu) est implicitement mobilisé pour affirmer qu'on ne peut pas « donner et ne pas donner » la bourse simultanément ?
    \item Le raisonnement est-il valide ? Est-il solide (prémisses vraies) ? Justifiez.
    \item Reformulez l'argument en un syllogisme hypothético-disjonctif valide (modus ponens ou modus tollens).
\end{enumerate}
\end{exercice}

\courssub{Je me prépare au Bac}

\begin{exercice}[Sujet de dissertation type Baccalauréat Technique]
\textbf{Sujet :} « La logique suffit-elle à garantir la vérité d'un discours ? »
\begin{enumerate}
    \item Analysez le sujet : définissez les termes clés (\cle{logique}, \cle{vérité formelle}, \cle{vérité matérielle}, \cle{discours}), identifiez la tension (validité vs vérité) et formulez la problématique.
    \item Proposez un plan détaillé en trois parties (thèse, antithèse, synthèse) avec pour chaque partie : l'idée directrice, les arguments philosophiques (auteurs, concepts), et un exemple concret camerounais (scolaire, médiatique, social, technique).
    \item Rédigez l'introduction complète (accroche, définition des termes, problématique, annonce du plan).
    \item Traitez la première partie de la dissertation (thèse : la logique comme condition nécessaire de la vérité du discours) en environ 20 lignes, en intégrant une référence à Aristote ou aux principes logiques.
\end{enumerate}
\end{exercice}

\begin{exercice}[Analyse de document : débat public et sophismes]
On lit ci-dessous un extrait fictif d'un débat télévisé camerounais sur la réforme de l'enseignement technique :

\begin{quote}
\textbf{Animateur :} « Honorable Député, les syndicats d'enseignants dénoncent le manque d'équipements dans les lycées techniques. Que répondez-vous ? »

\textbf{Député X :} « Écoutez, mes chers compatriotes, ces syndicats sont dirigés par des gens qui n'ont jamais géré une entreprise. Ils ne connaissent rien à la gestion budgétaire. D'ailleurs, le Secrétaire Général du principal syndicat a été vu l'an dernier dans une manifestation de l'opposition. C'est donc une manœuvre politique pour déstabiliser le gouvernement. Si on écoute ces gens-là, on va augmenter les impôts de tout le monde pour acheter des machines qui rouilleront dans les coins. Est-ce que vous voulez payer plus d'impôts ? Moi, je dis non. Le gouvernement a déjà construit deux lycées techniques flambant neufs à Yaoundé et Douala, donc le problème est réglé. Ceux qui critiquent sont soit de mauvaise foi, soit ignorants. »
\end{quote}

\begin{enumerate}
    \item Relevez \textbf{deux sophismes} distincts commis par le Député X. Pour chacun : citez la phrase exacte, nommez le sophisme (ex: \emph{ad hominem}, \emph{épouvantail}, \emph{fausse alternative}, \emph{généralisation hâtive}, \emph{non sequitur}, \emph{appel à la peur}, etc.) et expliquez pourquoi c'est un sophisme.
    \item Identifiez le \textbf{principe logique de la pensée} (identité, non-contradiction, tiers exclu) qui est violé ou mis à mal par l'enchaînement argumentatif global du député. Justifiez.
    \item Rédigez un paragraphe argumenté (15-20 lignes) critiquant la démarche du Député X du point de vue de la rigueur logique, en distinguant la validité formelle de ses arguments et la vérité matérielle de ses affirmations.
\end{enumerate}
\end{exercice}

\begin{exercice}[Question de synthèse : logique et décision administrative]
Le conseil d'administration d'un Lycée Technique Industriel doit attribuer une unique bourse d'excellence technique (critères : moyenne $\ge$ 14/20, assiduité $\ge$ 95\%, projet technique innovant présenté). Trois candidats sont en lice :
\begin{itemize}
    \item \textbf{Candidat A :} Moyenne 15,5 ; Assiduité 98\% ; Prototype de sécheuse solaire pour produits agricoles (validé par un jury externe).
    \item \textbf{Candidat B :} Moyenne 14,0 ; Assiduité 96\% ; Aucun projet présenté (dossier incomplet).
    \item \textbf{Candidat C :} Moyenne 16,0 ; Assiduité 90\% ; Projet d'application mobile de gestion de stock (non finalisé).
\end{itemize}
Le règlement intérieur stipule : « La bourse est attribuée au candidat remplissant \textbf{tous} les critères. En cas d'égalité, le jury départage sur l'originalité du projet. »

\begin{enumerate}
    \item En utilisant le \textbf{principe d'identité} (A est A), montrez comment ce principe permet de définir sans ambiguïté les critères d'éligibilité et d'éviter l'arbitraire dans l'instruction des dossiers.
    \item En utilisant le \textbf{principe de non-contradiction}, expliquez pourquoi un candidat ne peut pas être simultanément « éligible » (remplit tous les critères) et « non éligible » (ne remplit pas tous les critères) au regard du même règlement, au même moment.
    \item En utilisant le \textbf{principe du tiers exclu}, montrez comment ce principe structure la décision finale : pour chaque candidat, la proposition « Ce candidat remplit tous les critères » est soit vraie, soit fausse, sans état intermédiaire. Appliquez-le aux trois candidats pour déterminer la liste des éligibles.
    \item Concluez : en quoi le respect rigoureux de ces trois principes garantit-il la \cle{cohérence}, la \cle{justice} et la \cle{transparence} de la décision administrative ? (Réponse synthétique en 10-15 lignes).
\end{enumerate}
\end{exercice}

\newpage

\coursec{Corrigés des exercices}

\begin{corrige}[Exercice 1 — QCM : notions fondamentales]
\begin{enumerate}
    \item \textbf{Réponse b.} La logique formelle étudie la \cle{structure} et la \cle{validité} du raisonnement, indépendamment du contenu (vérité matérielle) des prémisses. C'est précisément ce qui la distingue d'une enquête empirique : un raisonnement peut être valide même si ses prémisses sont fausses.
    \item \textbf{Réponse b.} L'\cle{extension} d'un concept est l'ensemble des objets (individus, classes) auxquels il s'applique. L'ensemble des attributs essentiels, c'est la \cle{compréhension} (réponse a, donc fausse).
    \item \textbf{Réponse d.} « Certains élèves... ne sont pas boursiers » : quantité \emph{particulière} (« certains ») et qualité \emph{négative} (« ne sont pas ») : c'est une proposition de type \textbf{O} (particulière négative).
    \item \textbf{Réponse b.} « Une chose ne peut pas être et ne pas être en même temps et sous le même rapport » est la formulation aristotélicienne du \cle{principe de non-contradiction}. Le tiers exclu (c) affirme qu'une proposition est soit vraie soit fausse ; l'identité (a) affirme que A est A.
\end{enumerate}
\textbf{Points de vigilance :} ne pas confondre extension et compréhension ; bien repérer la négation dans la question 3 ; distinguer non-contradiction (impossibilité d'être et de ne pas être) et tiers exclu (nécessité d'être vrai ou faux).
\end{corrige}

\begin{corrige}[Exercice 2 — Vrai ou Faux justifié]
\begin{enumerate}
    \item \textbf{Faux.} La validité formelle garantit seulement que \emph{si} les prémisses sont vraies, \emph{alors} la conclusion est vraie. Si une prémisse est fausse, un raisonnement parfaitement valide peut conclure faux. Exemple : « Tous les Camerounais sont des ministres ; or Paul est camerounais ; donc Paul est ministre » : raisonnement valide, conclusion fausse. La vérité matérielle exige des prémisses vraies.
    \item \textbf{Vrai.} C'est la loi de l'\cle{extension et de la compréhension} : plus un concept s'applique à d'objets (grande extension), moins il contient d'attributs (faible compréhension), et inversement. « Être vivant » a une extension immense mais une compréhension pauvre ; « élève de 1re technique du lycée de Douala » a une extension réduite mais une compréhension riche.
    \item \textbf{Faux.} La figure 1 (M-P, S-M, donc S-P) ne suffit pas : il faut aussi respecter les règles générales du syllogisme (distribution du moyen, pas de deux prémisses négatives, etc.). Contre-exemple : « Tous les chiens sont des animaux (M-P) ; tous les chats sont des animaux (S-M) ; donc tous les chats sont des chiens (S-P) » : figure 1, mais le moyen « animaux » n'est distribué dans aucune prémisse : raisonnement invalide.
    \item \textbf{Vrai.} Le \cle{principe du tiers exclu} affirme qu'entre une proposition et sa négation, il n'y a pas de troisième possibilité : toute proposition est vraie ou fausse. C'est le fondement du raisonnement par l'absurde et de la logique bivalente classique.
\end{enumerate}
\textbf{Points de vigilance :} pour la question 3, la justification par contre-exemple est exigée ; pour la question 1, bien distinguer \emph{validité} (forme) et \emph{solidité} (validité + prémisses vraies).
\end{corrige}

\begin{corrige}[Exercice 3 — Définitions attendues]
\begin{enumerate}
    \item \textbf{Le concept} est une représentation mentale générale et abstraite qui saisit les caractères essentiels communs à une classe d'objets. Il possède deux dimensions : la \cle{compréhension}, ensemble des attributs essentiels qui le définissent (ex. pour « triangle » : figure plane, fermée, trois côtés) ; et l'\cle{extension}, ensemble des objets auxquels il s'applique (tous les triangles particuliers). Ces deux dimensions varient en sens inverse.
    \item \textbf{Le jugement} est une opération de l'esprit qui énonce un rapport d'accord ou de désaccord entre deux concepts, et qui peut être vrai ou faux. Sa structure est : \cle{sujet} (ce dont on parle) — \cle{copule} (« est » / « n'est pas », qui affirme ou nie) — \cle{attribut} (ce qu'on en dit). Exemple : « La baleine \emph{est} un mammifère ». On le classe selon la quantité (universel/particulier) et la qualité (affirmatif/négatif) : types A, E, I, O.
    \item \textbf{Le raisonnement} est une opération de l'esprit qui tire d'un ou plusieurs jugements (les prémisses) un jugement nouveau (la conclusion). La \cle{déduction} va du général au particulier : si les prémisses sont vraies et la forme valide, la conclusion est nécessairement vraie (ex. le syllogisme). L'\cle{induction} va du particulier au général : elle conclut avec une probabilité seulement, jamais une nécessité absolue (ex. « tous les élèves observés travaillent, donc tous les élèves travaillent »).
    \item \textbf{Les trois principes logiques de la pensée} (formulés par Aristote) :
    \begin{itemize}
        \item \cle{Principe d'identité} : A est A ; toute chose est ce qu'elle est ; un terme doit garder le même sens dans tout un raisonnement.
        \item \cle{Principe de non-contradiction} : une chose ne peut pas être et ne pas être en même temps et sous le même rapport ; une proposition et sa négation ne peuvent être vraies ensemble.
        \item \cle{Principe du tiers exclu} : entre une proposition et sa négation, il n'y a pas de troisième possibilité ; toute proposition est soit vraie, soit fausse.
    \end{itemize}
\end{enumerate}
\textbf{Points de vigilance :} une définition philosophique doit être précise, non circulaire, et comporter les distinctions demandées ; citer Aristote pour les principes est un plus valorisé.
\end{corrige}

\begin{corrige}[Exercice 4 — Classification logique et carrés d'opposition]
\begin{enumerate}
    \item Classification :
    \begin{itemize}
        \item « Tous les ingénieurs ont suivi une formation scientifique » : universelle affirmative $\rightarrow$ \textbf{A}.
        \item « Aucune machine n'est un être vivant » : universelle négative $\rightarrow$ \textbf{E}.
        \item « Certains logiciels sont libres » : particulière affirmative $\rightarrow$ \textbf{I}.
        \item « Certains téléphones portables ne sont pas des smartphones » : particulière négative $\rightarrow$ \textbf{O}.
    \end{itemize}
    \item À partir de A : « Tous les élèves de 1re technique étudient la philosophie » :
    \begin{itemize}
        \item \textbf{E} : « Aucun élève de 1re technique n'étudie la philosophie » — relation de \textbf{contrariété} avec A (les deux ne peuvent être vraies ensemble, mais peuvent être fausses ensemble).
        \item \textbf{I} : « Certains élèves de 1re technique étudient la philosophie » — relation de \textbf{sous-alternation} (si A est vraie, I est vraie ; l'inverse ne tient pas).
        \item \textbf{O} : « Certains élèves de 1re technique n'étudient pas la philosophie » — relation de \textbf{contradiction} avec A (l'une est vraie si et seulement si l'autre est fausse).
    \end{itemize}
    Rappel : la \textbf{sous-contrariété} lie I et O entre elles (elles ne peuvent être fausses ensemble, mais peuvent être vraies ensemble).
\end{enumerate}
\textbf{Points de vigilance :} bien mémoriser les quatre relations du carré : A-E contraires, A-O et E-I contradictoires, A-I et E-O sous-alternes, I-O sous-contraires. Ne pas confondre contraire (incompatibles mais pas exhaustives) et contradictoire (incompatibles \emph{et} exhaustives).
\end{corrige}

\begin{corrige}[Exercice 5 — Le carré d'opposition en situation]
\begin{enumerate}
    \item Types : P1 = \textbf{A} (universelle affirmative) ; P2 = \textbf{E} (universelle négative) ; P3 = \textbf{I} (particulière affirmative) ; P4 = \textbf{O} (particulière négative).
    \item Carré d'opposition (P1 et P2 en haut, P3 et P4 en bas) :
    \begin{itemize}
        \item P1 (A) $\longleftrightarrow$ P2 (E) : \textbf{contraires}.
        \item P1 (A) $\longleftrightarrow$ P4 (O) : \textbf{contradictoires}.
        \item P2 (E) $\longleftrightarrow$ P3 (I) : \textbf{contradictoires}.
        \item P1 (A) $\longrightarrow$ P3 (I) : \textbf{sous-alternation}.
        \item P2 (E) $\longrightarrow$ P4 (O) : \textbf{sous-alternation}.
        \item P3 (I) $\longleftrightarrow$ P4 (O) : \textbf{sous-contraires}.
    \end{itemize}
    \item Si P1 (A) est vraie :
    \begin{itemize}
        \item P2 (E), contraire de P1, est \textbf{fausse} (deux contraires ne peuvent être vraies ensemble).
        \item P4 (O), contradictoire de P1, est \textbf{fausse}.
        \item P3 (I), sous-alterne de P1, est \textbf{vraie} (la vérité descend de l'universelle à la particulière).
    \end{itemize}
    \item Si P3 (I) est fausse :
    \begin{itemize}
        \item P2 (E), contradictoire de P3, est \textbf{vraie}.
        \item P1 (A) : si A était vraie, sa sous-alterne I serait vraie ; or I est fausse, donc P1 est \textbf{fausse} (la fausseté remonte de la particulière à l'universelle).
        \item P4 (O), sous-contraire de P3, est \textbf{vraie} (deux sous-contraires ne peuvent être fausses ensemble). On vérifie : P4 est aussi sous-alterne de P2 qui est vraie, donc vraie. Cohérent.
    \end{itemize}
\end{enumerate}
\textbf{Points de vigilance :} la vérité « descend » (A vraie $\Rightarrow$ I vraie), la fausseté « monte » (I fausse $\Rightarrow$ A fausse) ; attention à ne jamais déduire la vérité de l'universelle à partir de la particulière.
\end{corrige}

\begin{corrige}[Exercice 6 — Validité des syllogismes : règles et diagrammes de Venn]
\textbf{Syllogisme 1 :} Tous les mammifères (M) sont des vertébrés (P) ; la baleine (S) est un mammifère (M) ; donc la baleine (S) est un vertébré (P).
\begin{itemize}
    \item Mode : \textbf{AAA} (Barbara) ; \textbf{figure 1} (M-P, S-M, donc S-P).
    \item Règles : trois termes seulement ; le moyen « mammifères » est distribué dans la majeure (universelle) ; aucune prémisse négative ; le terme S est distribué en conclusion comme en prémisse.
    \item \textbf{Valide.} Diagramme de Venn : le cercle « mammifères » est entièrement inclus dans « vertébrés » ; le point « baleine » placé dans « mammifères » se trouve nécessairement dans « vertébrés ».
\end{itemize}

\textbf{Syllogisme 2 :} Tous les fonctionnaires (P) sont assidus (M) ; Paul (S) est assidu (M) ; donc Paul est fonctionnaire (P).
\begin{itemize}
    \item Structure : P-M, S-M, donc S-P : \textbf{figure 2}, mode AAA.
    \item Règle violée : le \cle{terme moyen} « assidu » n'est \textbf{distribué dans aucune prémisse} (attribut de deux propositions affirmatives, donc non distribué). C'est le sophisme du \emph{moyen non distribué}.
    \item \textbf{Invalide.} Venn : le cercle « fonctionnaires » est inclus dans « assidus », mais le point « Paul » peut être placé dans « assidus » \emph{en dehors} de « fonctionnaires » : la conclusion ne s'impose pas. Paul peut être un commerçant assidu.
\end{itemize}

\textbf{Syllogisme 3 :} Aucun corrompu (M) n'est honnête (P) ; certains politiciens (S) sont corrompus (M) ; donc certains politiciens (S) ne sont pas honnêtes (P).
\begin{itemize}
    \item Mode : \textbf{EIO} (Ferio) ; \textbf{figure 1} (M-P, S-M, donc S-P).
    \item Règles : une seule prémisse négative, conclusion négative (conforme) ; le moyen « corrompu » est distribué dans la majeure (E distribue les deux termes) ; P, distribué en conclusion (O), l'est aussi dans la majeure E.
    \item \textbf{Valide.} Venn : les cercles « corrompus » et « honnêtes » sont disjoints ; la zone commune à « politiciens » et « corrompus » est non vide et se trouve hors de « honnêtes » : la conclusion O est forcée.
\end{itemize}
\textbf{Points de vigilance :} toujours identifier le moyen (absent de la conclusion) ; vérifier sa distribution ; rappeler que d'une prémisse particulière on ne peut tirer une conclusion universelle, et que de deux prémisses négatives on ne tire rien.
\end{corrige}

\begin{corrige}[Exercice 7 — Critique et correction de définitions]
\begin{enumerate}
    \item « Un triangle est une figure qui a trois côtés » : \textbf{trop large} (un angle à trois côtés non fermé, ou une figure à trois côtés courbes, y satisferait). Définition correcte : « Un triangle est une figure géométrique plane, fermée, formée de trois segments de droite (trois côtés) et de trois angles. »
    \item « La liberté, c'est faire ce qu'on veut » : \textbf{trop large et fausse} (elle inclut le caprice et la licence, et rend impossible toute vie en société : si chacun fait ce qu'il veut, la liberté d'autrui est détruite). Définition correcte (sens philosophique) : « La liberté est le pouvoir d'agir selon sa volonté éclairée par la raison, dans le respect de la loi et de la liberté d'autrui » (cf. Rousseau : « obéir à la loi qu'on s'est prescrite »).
    \item « Un menteur est celui qui dit des mensonges » : \textbf{circulaire} (on définit le menteur par « mensonge », terme de la même famille : le défini revient dans la définition). Définition correcte : « Un menteur est une personne qui énonce sciemment une proposition qu'elle sait fausse dans l'intention de tromper autrui. »
    \item « L'école est un bâtiment où vont les enfants » : \textbf{trop large et trop étroite à la fois} : trop large (un cinéma est un bâtiment où vont les enfants) et trop étroite (l'école n'est pas d'abord un bâtiment : on peut faire l'école sous un hangar ; et elle accueille aussi des adolescents et des adultes). Définition correcte : « L'école est une institution sociale chargée de transmettre aux jeunes générations des savoirs, des savoir-faire et des valeurs par un enseignement organisé. »
\end{enumerate}
\textbf{Points de vigilance :} une bonne définition doit être \emph{convertible} (ni trop large ni trop étroite), \emph{non circulaire}, \emph{claire} (pas obscure) et \emph{affirmative} quand c'est possible. Nommer précisément le vice est exigé.
\end{corrige}

\begin{corrige}[Exercice 8 — Application des principes logiques à l'argumentation courante]
\begin{enumerate}
    \item Formalisation. Soit $p$ : « on accorde la bourse à Amadou » ; $q$ : « on accorde la bourse à tout le monde » ; $r$ : « le budget permet de donner une bourse à tous ».
    \begin{itemize}
        \item Prémisse 1 : $p \Rightarrow q$ (donner à Amadou oblige à donner à tous, par équité).
        \item Prémisse 2 : $\neg r$ (le budget ne permet pas de donner à tous), donc on ne peut pas donner à tous : $\neg q$.
        \item Conclusion : donc $\neg p$ (il ne faut pas donner à Amadou).
    \end{itemize}
    Structure : $[(p \Rightarrow q) \land \neg q] \Rightarrow \neg p$.
    \item Le \cle{principe de non-contradiction} est mobilisé : on ne peut pas à la fois « donner la bourse à tous » et « ne pas la donner à tous » (donner et ne pas donner en même temps et sous le même rapport). C'est ce qui permet de passer de « le budget ne le permet pas » à « il est impossible de donner à tous ».
    \item Le raisonnement est \textbf{formellement valide} : c'est un \emph{modus tollens} ($p \Rightarrow q$ ; or $\neg q$ ; donc $\neg p$). Mais sa \textbf{solidité} est discutable : la prémisse 1 est contestable, car l'équité n'exige pas de traiter de la même façon des situations différentes ; si Amadou remplit des critères d'excellence que les autres ne remplissent pas, donner à Amadou n'oblige pas à donner à tous. Le principe d'égalité porte sur les cas \emph{identiques}, non sur tous les cas. Le raisonnement est donc valide mais non nécessairement solide.
    \item Reformulation en syllogisme hypothétique (modus tollens) :
    \begin{itemize}
        \item Majeure : Si l'on accorde la bourse à Amadou sans critère distinctif, alors l'équité oblige à l'accorder à tous.
        \item Mineure : Or il est impossible de l'accorder à tous (budget insuffisant).
        \item Conclusion : Donc il ne faut pas l'accorder à Amadou sans critère distinctif.
    \end{itemize}
\end{enumerate}
\textbf{Points de vigilance :} distinguer explicitement validité (forme) et solidité (vérité des prémisses) ; nommer le \emph{modus tollens} ; montrer que la critique porte sur la prémisse, non sur la forme.
\end{corrige}

\begin{corrige}[Exercice 9 — Sujet de dissertation type Baccalauréat Technique]
\begin{enumerate}
    \item \textbf{Analyse du sujet.} \cle{Logique} : science des lois formelles de la pensée, qui garantit la validité (cohérence) du raisonnement. \cle{Vérité formelle} : accord du discours avec les règles de la logique. \cle{Vérité matérielle} : accord du discours avec la réalité. \cle{Discours} : ensemble organisé d'énoncés visant à affirmer ou à persuader. Tension : un discours peut être parfaitement cohérent (valide) et pourtant faux dans les faits ; inversement, peut-on atteindre la vérité sans logique ? Problématique : la logique est-elle une condition suffisante de la vérité, ou seulement une condition nécessaire que doit compléter l'expérience ?
    \item \textbf{Plan détaillé.}
    \begin{itemize}
        \item \textbf{I. Thèse : la logique est la condition première de tout discours vrai.} Sans les principes d'identité, de non-contradiction et du tiers exclu (Aristote), aucune pensée cohérente n'est possible ; un discours contradictoire ne peut être vrai. Exemple : un élève de 1re technique qui affirme « ce circuit est fermé et ouvert en même temps » ne dit rien de pensable ; le jury du Probatoire exige la cohérence de la copie.
        \item \textbf{II. Antithèse : la logique ne suffit pas, car elle ne garantit que la validité, non la vérité.} Un syllogisme valide à prémisses fausses conclut faux (« Tous les élèves de ce lycée sont boursiers ; or Amadou est élève ; donc il est boursier »). La vérité matérielle exige l'expérience et la vérification des faits. Exemple : les rumeurs sur les réseaux sociaux au Cameroun enchaînent parfois des raisonnements cohérents sur des prémisses inventées.
        \item \textbf{III. Synthèse : la logique est une condition nécessaire mais non suffisante.} Le discours vrai exige à la fois la rigueur formelle (logique) et la vérité des prémisses (expérience, enquête, honnêteté intellectuelle). Exemple : l'ingénieur camerounais qui diagnostique une panne combine déduction rigoureuse et observation concrète de la machine.
    \end{itemize}
    \item \textbf{Introduction rédigée.} Tout débat, du conseil de classe à l'Assemblée nationale, prétend dire le vrai. Mais qu'est-ce qui garantit la vérité d'un discours ? La \emph{logique}, science des lois formelles de la pensée, assure la cohérence et la validité de nos raisonnements ; la \emph{vérité}, elle, peut s'entendre de deux façons : vérité formelle (accord du discours avec ses propres règles) et vérité matérielle (accord du discours avec la réalité). Dès lors, un problème se pose : la logique suffit-elle à garantir la vérité d'un discours ? Autrement dit, un raisonnement correctement construit est-il nécessairement un discours vrai ? Nous montrerons d'abord que la logique est la condition première de toute pensée vraie ; nous verrons ensuite qu'elle ne suffit pas, car la validité n'est pas la vérité ; nous conclurons enfin qu'elle est une condition nécessaire mais non suffisante, que doit compléter l'expérience.
    \item \textbf{Première partie (thèse).} La logique est d'abord ce sans quoi aucun discours ne peut prétendre à la vérité. Aristote, fondateur de la logique, établit que toute pensée repose sur trois principes : le principe d'identité (une chose est ce qu'elle est), le principe de non-contradiction (une chose ne peut être et ne pas être en même temps et sous le même rapport) et le principe du tiers exclu (une proposition est vraie ou fausse, sans troisième possibilité). Ces principes ne sont pas des conventions arbitraires : ils sont les conditions mêmes de l'intelligibilité. Celui qui les viole ne dit plus rien : affirmer qu'un moteur « fonctionne et ne fonctionne pas » dans le même sens et au même moment, c'est détruire toute possibilité de diagnostic, donc toute vérité technique. De même, le syllogisme aristotélicien montre comment une conclusion vraie peut être nécessairement tirée de prémisses vraies : « Tous les métaux conduisent l'électricité ; or le cuivre est un métal ; donc le cuivre conduit l'électricité ». Ici, la logique n'ajoute rien au contenu, mais elle \emph{transmet} la vérité des prémisses à la conclusion : elle est le canal de la vérité. Sans elle, même des prémisses vraies resteraient stériles, incapables de produire un savoir ordonné. C'est pourquoi toute science, toute technique et toute justice supposent la rigueur logique : un jugement contradictoire est injuste, un calcul incohérent est dangereux, un raisonnement invalide est nul. La logique apparaît ainsi comme la condition nécessaire de tout discours vrai.
\end{enumerate}
\textbf{Barème indicatif (sur 20) :} analyse et problématique 4 pts ; plan cohérent et progressif 4 pts ; introduction rédigée 4 pts ; développement de la thèse (arguments, référence à Aristote, exemple) 6 pts ; correction de la langue 2 pts. \textbf{Points de vigilance :} la problématique doit être une question ; chaque partie doit contenir une idée, un argument, une référence et un exemple ; éviter la paraphrase du sujet.
\end{corrige}

\begin{corrige}[Exercice 10 — Analyse de document : débat public et sophismes]
\begin{enumerate}
    \item Deux sophismes (parmi d'autres possibles) :
    \begin{itemize}
        \item \textbf{Sophisme \emph{ad hominem} :} « le Secrétaire Général du principal syndicat a été vu l'an dernier dans une manifestation de l'opposition. C'est donc une manœuvre politique ». Au lieu de répondre à l'argument (le manque d'équipements), le député attaque la \emph{personne} et les intentions supposées des syndicalistes. Or la valeur d'une critique ne dépend pas de l'identité de celui qui la formule : même un opposant peut dire vrai.
        \item \textbf{Sophisme de l'épouvantail (ou de l'homme de paille) :} « Si on écoute ces gens-là, on va augmenter les impôts de tout le monde pour acheter des machines qui rouilleront dans les coins. » Le député déforme la revendication (demander des équipements) en une caricature absurde (gaspillage fiscal) plus facile à réfuter. On peut aussi y voir un \textbf{appel à la peur} (« Est-ce que vous voulez payer plus d'impôts ? »).
        \item Autres réponses acceptables : \textbf{généralisation hâtive} (« deux lycées flambant neufs à Yaoundé et Douala, donc le problème est réglé » : deux cas ne suffisent pas à conclure pour tout le pays) ; \textbf{fausse alternative} (« Ceux qui critiquent sont soit de mauvaise foi, soit ignorants » : il existe une troisième possibilité, la critique fondée et sincère).
    \end{itemize}
    \item C'est le \cle{principe du tiers exclu} qui est mis à mal par la fausse alternative finale : « soit de mauvaise foi, soit ignorants » prétend épuiser les possibilités alors qu'une troisième existe (être de bonne foi et informé). Plus largement, le principe de \cle{non-contradiction} est fragilisé : le député laisse entendre que le problème est à la fois réel (il faut construire des lycées, ce qu'il vante) et réglé (il n'y a plus de problème), sans trancher. Enfin, le \cle{principe d'identité} est bafoué par le glissement de sens : « régler le problème de l'équipement » devient « construire deux bâtiments », alors que ce n'est pas la même chose (un bâtiment n'est pas un équipement pédagogique).
    \item \textbf{Paragraphe critique.} La démonstration du Député X illustre la différence capitale entre validité formelle et vérité matérielle. Sur le plan formel, aucun de ses enchaînements n'est valide : de « le syndicaliste a manifesté avec l'opposition », on ne peut déduire « sa critique est fausse » (non sequitur) ; de « deux lycées ont été construits », on ne peut déduire « tous les lycées techniques sont équipés » (induction abusive). Sur le plan matériel, ses prémisses elles-mêmes sont douteuses : l'affirmation selon laquelle les machines « rouilleront » est une pure conjecture, et l'existence de deux lycées neufs ne dit rien de leur équipement réel. Le député substitue ainsi la \emph{persuasion} à la \emph{démonstration} : il joue sur les émotions (peur des impôts), attaque les personnes et déforme la thèse adverse au lieu d'examiner les faits. Or, comme l'enseigne la logique classique, un discours vrai exige à la fois des prémisses vraies et une forme valide : ici, les deux font défaut. La rigueur logique aurait exigé de répondre à la question posée — les lycées techniques sont-ils suffisamment équipés ? — par des données vérifiables et un raisonnement correct. Ce discours est donc un cas typique de rhétorique sophistique : il peut convaincre un public inattentif, mais il ne prouve rien.
\end{enumerate}
\textbf{Barème indicatif :} deux sophismes correctement nommés, cités et expliqués : 6 pts ; principe logique identifié et justifié : 4 pts ; paragraphe argumenté distinguant validité et vérité : 8 pts ; langue : 2 pts. \textbf{Points de vigilance :} toujours citer la phrase exacte ; nommer le sophisme en latin ou en français ; ne pas confondre « argument invalide » et « affirmation fausse ».
\end{corrige}

\begin{corrige}[Exercice 11 — Question de synthèse : logique et décision administrative]
\begin{enumerate}
    \item \textbf{Principe d'identité (A est A).} Ce principe impose que chaque critère garde un sens fixe et déterminé tout au long de la procédure : « moyenne $\ge$ 14/20 » signifie exactement cela, et non « environ 14 » ou « 14 si le candidat est sympathique » ; « assiduité $\ge$ 95\% » ne peut pas devenir « 90\% » en cours de route ; « projet innovant présenté » veut dire un projet effectivement déposé et finalisé. Grâce à l'identité, les critères sont des concepts à extension claire : on peut déterminer sans ambiguïté quels dossiers y appartiennent. Toute modification interprétative en faveur d'un candidat violerait le principe et introduirait l'arbitraire.
    \item \textbf{Principe de non-contradiction.} Un même candidat, au même moment et au regard du même règlement, ne peut pas être à la fois « éligible » (il remplit tous les critères) et « non éligible » (il ne les remplit pas tous). Affirmer les deux serait une contradiction logique qui détruirait la décision : la bourse ne peut être à la fois attribuée et refusée au même candidat. Ce principe oblige donc le jury à trancher chaque dossier de manière cohérente et interdit toute décision « à double détente » (par exemple déclarer C éligible pour son projet tout en le recalant pour son assiduité).
    \item \textbf{Principe du tiers exclu.} Pour chaque candidat, la proposition « ce candidat remplit tous les critères » est soit vraie, soit fausse : il n'existe pas de troisième état (« à moitié éligible »). Application :
    \begin{itemize}
        \item \textbf{Candidat A :} moyenne 15,5 $\ge$ 14 (vrai) ; assiduité 98\% $\ge$ 95\% (vrai) ; projet présenté et validé (vrai). La proposition est \textbf{vraie} : A est \textbf{éligible}.
        \item \textbf{Candidat B :} moyenne 14,0 $\ge$ 14 (vrai) ; assiduité 96\% (vrai) ; mais aucun projet présenté (faux). La conjonction est \textbf{fausse} : B est \textbf{non éligible}.
        \item \textbf{Candidat C :} moyenne 16,0 (vrai) ; projet présenté mais non finalisé (discutable, admettons présenté) ; mais assiduité 90\% $<$ 95\% (faux). La conjonction est \textbf{fausse} : C est \textbf{non éligible}.
    \end{itemize}
    Une conjonction n'est vraie que si \emph{tous} ses termes sont vrais : il suffit d'un critère non rempli pour que la proposition globale soit fausse. Seul \textbf{A} est éligible ; la bourse lui est donc attribuée, sans qu'il soit besoin du critère de départage (pas d'égalité).
    \item \textbf{Conclusion.} Le respect des trois principes logiques fait de la décision administrative un acte rationnel et non un acte de pouvoir. Le principe d'identité garantit la \cle{transparence} : les critères, définis une fois pour toutes, sont publics et stables, chacun peut vérifier leur application. Le principe de non-contradiction garantit la \cle{cohérence} : le jury ne peut pas tenir deux discours incompatibles sur le même dossier, ce qui interdit le favoritisme dissimulé. Le principe du tiers exclu garantit la \cle{justice} : chaque candidat reçoit une réponse claire, éligible ou non, fondée sur les mêmes règles pour tous ; l'égalité de traitement est ainsi assurée. La logique apparaît donc comme l'instrument de la bonne administration : en soumettant la décision à des règles formelles, elle protège les citoyens — ici les élèves — contre l'arbitraire et rend la décision explicable, contestable et légitime.
\end{enumerate}
\textbf{Barème indicatif :} application correcte de chaque principe (3 $\times$ 3 pts = 9 pts) ; détermination exacte des éligibles avec justification logique (conjonction) : 4 pts ; conclusion synthétique reliant les trois principes aux trois valeurs : 5 pts ; langue : 2 pts. \textbf{Points de vigilance :} préciser « en même temps et sous le même rapport » pour la non-contradiction ; rappeler qu'une conjonction est fausse dès qu'un seul critère manque ; ne pas attribuer la bourse à C malgré sa meilleure moyenne : le règlement prime.
\end{corrige}

\newpage

\coursec{Fiche bilan}

\begin{popfigure}
\begin{tikzpicture}[mindmap, grow cyclic, every node/.style={concept, text=white, font=\small\bfseries},
root concept/.append style={concept color=popdark, font=\normalsize\bfseries},
level 1/.append style={level distance=3.4cm, sibling angle=60},
level 2/.append style={level distance=2.2cm, sibling angle=45, font=\scriptsize}]
\node[root concept]{La logique classique}
  child[concept color=popblue]{ node{Définition et importance}
    child{ node{Science du raisonnement correct} }
    child{ node{Outil de pensée claire} }
  }
  child[concept color=poporange]{ node{Le concept}
    child{ node{Compréhension et extension} }
    child{ node{Définition : genre prochain + différence spécifique} }
  }
  child[concept color=popgreen]{ node{Le jugement}
    child{ node{Affirme ou nie} }
    child{ node{Vrai ou faux} }
  }
  child[concept color=poppurple]{ node{Le raisonnement}
    child{ node{Déductif : syllogisme} }
    child{ node{Inductif} }
  }
  child[concept color=poppink]{ node{Principes logiques}
    child{ node{Identité} }
    child{ node{Non-contradiction} }
    child{ node{Tiers exclu} }
  }
  child[concept color=popgold]{ node{Application}
    child{ node{Penser juste au quotidien} }
    child{ node{Démasquer les sophismes} }
  };
\end{tikzpicture}
\legende{Carte mentale de la leçon : la logique classique, science du raisonnement correct, repose sur le concept, le jugement, le raisonnement et les principes fondamentaux de la pensée.}
\end{popfigure}

\begin{bilan}
\textbf{Définitions clés à connaître}
\begin{itemize}
  \item \cle{Logique} : science qui étudie les conditions de validité du raisonnement ; elle enseigne l'art de penser correctement (du grec \emph{logos} : parole, raison).
  \item \cle{Concept} : représentation générale et abstraite d'un objet de pensée (ex. : « homme », « justice »). Il possède une \cle{compréhension} (ensemble des caractères) et une \cle{extension} (ensemble des êtres visés).
  \item \cle{Définition} (d'un concept) : elle se fait par le \cle{genre prochain} et la \cle{différence spécifique} (ex. : « l'homme est un animal raisonnable »).
  \item \cle{Jugement} : opération de l'esprit qui affirme ou nie un rapport entre deux concepts (ex. : « l'homme est mortel »). Il est vrai ou faux.
  \item \cle{Raisonnement} : opération qui tire d'un ou plusieurs jugements (les prémisses) un jugement nouveau (la conclusion).
\end{itemize}

\textbf{Les trois principes logiques de la pensée}
\begin{center}
\begin{tabularx}{\linewidth}{|l|X|X|}
\hline
\rowcolor{popblueL} \textbf{Principe} & \textbf{Formule} & \textbf{Signification} \\
\hline
Identité & $A$ est $A$ & Un concept garde le même sens dans tout le raisonnement. \\
\hline
Non-contradiction & $A$ n'est pas non-$A$ & On ne peut affirmer et nier la même chose en même temps et sous le même rapport. \\
\hline
Tiers exclu & $A$ ou non-$A$ & Entre une affirmation et sa négation, il n'y a pas de troisième possibilité. \\
\hline
\end{tabularx}
\end{center}

\textbf{Les formes de raisonnement}
\begin{itemize}
  \item \cle{Raisonnement déductif} : du général au particulier ; la conclusion est nécessaire si les prémisses sont vraies. Forme type : le \cle{syllogisme} (majeure, mineure, conclusion).
  \item \cle{Raisonnement inductif} : du particulier au général ; la conclusion est probable, non certaine.
\end{itemize}

\textbf{Méthode express : analyser un syllogisme}
\begin{enumerate}
  \item Repérer la conclusion et les deux prémisses (majeure et mineure).
  \item Vérifier que les termes gardent le même sens (principe d'identité).
  \item Vérifier que la conclusion découle nécessairement des prémisses.
  \item Conclure : raisonnement \emph{valide} (forme correcte) ou \emph{sophisme} (raisonnement faux).
\end{enumerate}
\end{bilan}

\begin{aretenir}[Checklist avant le Bac]
\begin{itemize}
  \item[$\square$] Je sais définir la logique et expliquer son importance pour la pensée et la vie pratique.
  \item[$\square$] Je distingue clairement concept, jugement et raisonnement, avec un exemple pour chacun.
  \item[$\square$] Je connais la différence entre compréhension et extension d'un concept.
  \item[$\square$] Je sais construire une définition par le genre prochain et la différence spécifique.
  \item[$\square$] Je sais qu'un jugement affirme ou nie, et qu'il est vrai ou faux.
  \item[$\square$] Je distingue raisonnement déductif (nécessaire) et raisonnement inductif (probable).
  \item[$\square$] Je sais identifier les trois parties d'un syllogisme : majeure, mineure, conclusion.
  \item[$\square$] Je connais par cœur les trois principes logiques : identité, non-contradiction, tiers exclu.
  \item[$\square$] Je sais illustrer chaque principe par un exemple concret de la vie quotidienne camerounaise.
  \item[$\square$] Je sais repérer un sophisme et expliquer pourquoi le raisonnement est invalide.
  \item[$\square$] Je sais rédiger une introduction de dissertation qui problématise à partir d'une situation concrète.
  \item[$\square$] Je justifie toujours mes réponses par des définitions précises et des arguments ordonnés.
\end{itemize}
\end{aretenir}

\findecours

\end{document}
